% Options for packages loaded elsewhere
\PassOptionsToPackage{unicode}{hyperref}
\PassOptionsToPackage{hyphens}{url}
\PassOptionsToPackage{dvipsnames,svgnames,x11names}{xcolor}
%
\documentclass[
  10pt,
]{article}
\usepackage{amsmath,amssymb}
\usepackage{iftex}
\ifPDFTeX
  \usepackage[T1]{fontenc}
  \usepackage[utf8]{inputenc}
  \usepackage{textcomp} % provide euro and other symbols
\else % if luatex or xetex
  \usepackage{unicode-math} % this also loads fontspec
  \defaultfontfeatures{Scale=MatchLowercase}
  \defaultfontfeatures[\rmfamily]{Ligatures=TeX,Scale=1}
\fi
\usepackage{lmodern}
\ifPDFTeX\else
  % xetex/luatex font selection
  \setmainfont[]{DejaVu Serif}
  \setsansfont[]{DejaVu Sans}
  \setmonofont[]{DejaVu Sans Mono}
  \setmathfont[]{Latin Modern Math}
\fi
% Use upquote if available, for straight quotes in verbatim environments
\IfFileExists{upquote.sty}{\usepackage{upquote}}{}
\IfFileExists{microtype.sty}{% use microtype if available
  \usepackage[]{microtype}
  \UseMicrotypeSet[protrusion]{basicmath} % disable protrusion for tt fonts
}{}
\makeatletter
\@ifundefined{KOMAClassName}{% if non-KOMA class
  \IfFileExists{parskip.sty}{%
    \usepackage{parskip}
  }{% else
    \setlength{\parindent}{0pt}
    \setlength{\parskip}{6pt plus 2pt minus 1pt}}
}{% if KOMA class
  \KOMAoptions{parskip=half}}
\makeatother
\usepackage{xcolor}
\usepackage[margin=2.0cm]{geometry}
\usepackage{longtable,booktabs,array}
\usepackage{calc} % for calculating minipage widths
% Correct order of tables after \paragraph or \subparagraph
\usepackage{etoolbox}
\makeatletter
\patchcmd\longtable{\par}{\if@noskipsec\mbox{}\fi\par}{}{}
\makeatother
% Allow footnotes in longtable head/foot
\IfFileExists{footnotehyper.sty}{\usepackage{footnotehyper}}{\usepackage{footnote}}
\makesavenoteenv{longtable}
\setlength{\emergencystretch}{3em} % prevent overfull lines
\providecommand{\tightlist}{%
  \setlength{\itemsep}{0pt}\setlength{\parskip}{0pt}}
\setcounter{secnumdepth}{-\maxdimen} % remove section numbering
% doc_header.tex -- LaTeX preamble additions used by build_docs.sh for every
% Markdown -> LaTeX -> PDF conversion (README and provenance documents of Revision_SuperLoop).
\usepackage{etoolbox}
% longtable may not be loaded by pandoc when a document has no table
\usepackage{longtable}
% bold vectors: \bm -> \symbf under unicode-math (XeLaTeX), bm package otherwise
\makeatletter
\@ifpackageloaded{unicode-math}{\providecommand{\bm}[1]{\symbf{#1}}}{\usepackage{bm}}
\makeatother
% smaller type in long tables so that wide audit tables fit the text width;
% inside a table a line break is allowed after an underscore of plain text
% (pandoc writes it as \_), so that long identifiers wrap inside their column
% rows are set a little apart, so that symbols with indices do not touch
\AtBeginEnvironment{longtable}{\footnotesize\renewcommand{\arraystretch}{1.45}%
  \let\sptunderscore\_\def\_{\sptunderscore\allowbreak}}
\setlength{\LTcapwidth}{\textwidth}
% allow line breaks after slashes and hyphens inside table cells
\usepackage{xurl}
% code blocks: break over-long lines instead of running into the margin.
% Plain ``` blocks are typeset by pandoc as {verbatim}, highlighted ones
% (```bash) as {Highlighting}; both are re-declared as breakable fancyvrb
% environments.  The continuation line is marked with a small hooked arrow.
\usepackage{fvextra}
\RecustomVerbatimEnvironment{verbatim}{Verbatim}{breaklines,breakanywhere,fontsize=\small}
\makeatletter
\@ifundefined{Highlighting}{}{%
  \RecustomVerbatimEnvironment{Highlighting}{Verbatim}{breaklines,breakanywhere,fontsize=\small,commandchars=\\\{\}}}
\makeatother
% inline code (pandoc writes it as \texttt{...}): never hyphenate -- an inserted
% hyphen could be read as part of a file name -- but allow a line break after
% an underscore, slash, comma, colon or equals sign, so that long paths and
% column lists wrap inside the text width and inside table columns.
% Inside a formula (e.g. the name of a notebook symbol in an equation) the same
% command must first leave math mode, which \text of amsmath does.
\ExplSyntaxOn
\tl_new:N \l__spt_code_tl
\cs_new_protected:Npn \__spt_texttt:n #1
  {
    \tl_set:Nn \l__spt_code_tl {#1}
    \regex_replace_all:nnN { ([\/\,\:\=]) } { \1 \c{allowbreak} } \l__spt_code_tl
    \regex_replace_all:nnN { \c{_} } { \c{_} \c{allowbreak} } \l__spt_code_tl
    \group_begin: \ttfamily \hyphenchar \font = -1 \scan_stop: \l__spt_code_tl \group_end:
  }
\RenewDocumentCommand \texttt { m }
  {
    \mode_if_math:TF { \text { \__spt_texttt:n {#1} } } { \__spt_texttt:n {#1} }
  }
\ExplSyntaxOff
% A figure is written in the Markdown source as an image paragraph followed by
% a caption paragraph ("**Figure S1.** ..."), so that the caption is also
% visible on GitHub.  An image that starts a paragraph is centered.
\usepackage{graphicx}% pandoc loads it only for documents that contain an image
\NewCommandCopy\sptincludegraphics\includegraphics
\RenewDocumentCommand\includegraphics{O{} m}{%
  \ifvmode\begin{center}\sptincludegraphics[#1]{#2}\end{center}%
  \else\sptincludegraphics[#1]{#2}\fi}
% keep figures where they are referenced
\usepackage{float}
\floatplacement{figure}{H}
% running header with the document title
\usepackage{fancyhdr}
\pagestyle{fancy}
\fancyhf{}
\fancyhead[L]{\small PROJECT spacetimes, Revision\_SuperLoop}
\fancyhead[R]{\small\thepage}
\renewcommand{\headrulewidth}{0.4pt}
% 16x16 matrices: tighter columns so that a full matrix fits the text width
\setlength{\arraycolsep}{3.2pt}
% The documents of this project contain long tagged symbols, file paths and
% identifiers that cannot be hyphenated.  A paragraph is therefore allowed to
% have looser lines before TeX gives up and lets a line stick out of the margin.
\AtBeginDocument{\tolerance=2000 \emergencystretch=5em \hbadness=2000}
\ifLuaTeX
  \usepackage{selnolig}  % disable illegal ligatures
\fi
\IfFileExists{bookmark.sty}{\usepackage{bookmark}}{\usepackage{hyperref}}
\IfFileExists{xurl.sty}{\usepackage{xurl}}{} % add URL line breaks if available
\urlstyle{same}
\hypersetup{
  pdftitle={The calculations of Revision\_SuperLoop: formulas, code, commands, and the output of every check},
  pdfauthor={Prepared for Patrick Lee Nash, Ph.D.},
  colorlinks=true,
  linkcolor={blue},
  filecolor={Maroon},
  citecolor={blue},
  urlcolor={blue},
  pdfcreator={LaTeX via pandoc}}

\title{The calculations of Revision\_SuperLoop: formulas, code,
commands, and the output of every check}
\usepackage{etoolbox}
\makeatletter
\providecommand{\subtitle}[1]{% add subtitle to \maketitle
  \apptocmd{\@title}{\par {\large #1 \par}}{}{}
}
\makeatother
\subtitle{Provenance document \texttt{P2} of PROJECT spacetimes,
Revision\_SuperLoop}
\author{Prepared for Patrick Lee Nash, Ph.D.}
\date{November 16, 2026}

\begin{document}
\maketitle

\hypertarget{what-this-document-is}{%
\section{1. What this document is}\label{what-this-document-is}}

This document states, for every computation of the folder
\texttt{Revision\_SuperLoop}, the formulas that the programs evaluate,
how they evaluate them, the command that runs them, and the complete
output of every check, as the programs printed it in the last run.
Section 2 says how to reproduce everything; Section 3 describes the
programs; Sections 4 to 8 go through the computations (the algebra, the
classification, the manifolds, the superfluid, renormalizability);
Section 9 describes the program of the examples of the explanation of
the sentence on the superfluid (folder \texttt{explanations}) and gives
its output; Section 10 does the same for the program of the document on
the \(3+1\) manifold and the renormalizability of the Yukawa form
(folder \texttt{explanations}); Section 11 reproduces the checks of the
documents; Section 12 lists the data files and the figures; Section 13
gives the checksums; Section 14 is the complete table of the symbols,
which defines every symbol used here.

\hypertarget{how-to-reproduce-everything}{%
\section{2. How to reproduce
everything}\label{how-to-reproduce-everything}}

From the folder \texttt{Revision\_SuperLoop} of the repository:

\begin{verbatim}
cargo run --release --manifest-path code/superloop-rs/Cargo.toml -- all   # data files and checks of the program
python3 code/scripts/malcev_groebner.py                                  # Groebner bases, checks G-1 .. G-12
python3 code/scripts/make_figures.py                                     # the seven figures
(cd explanations && cargo run --release --manifest-path code/explain-rs/Cargo.toml)   # checks X-1 .. X-10
python3 explanations/code/make_explain_figures.py                        # the three figures of the explanation
(cd explanations && cargo run --release --manifest-path code/r13-rs/Cargo.toml)      # checks Z-1 .. Z-8
python3 explanations/code/make_r13_figures.py                            # the four figures of the 3+1 document
bash code/scripts/check_all.sh                                           # everything, with one verdict
\end{verbatim}

\texttt{check\_all.sh} runs, in this order: the program (all five
commands); the Gröbner script; the figures; the program of the examples
of the explanation and its figures; the program of the document on the
\(3+1\) manifold and its figures; twice in a row the registry of the
symbols, the fragments of the documents, the checks of the symbols and
of self-containedness, the assembly of the Markdown documents, the
manuscript (.tex, .pdf, .md) and the provenance documents (.tex, .pdf)
(the second pass makes every document contain the reports of the first);
then it verifies that every report passes, that no placeholder is left
unresolved, that every report printed in a document is the current one,
that no formula ends directly before a digit, that every document exists
as .md, .tex and .pdf, and that no document names a software tool that
does not belong in a manuscript, and prints one verdict, ``ALL PASS'' or
the list of failures. Requirements: a Rust toolchain for the edition
2021; Python 3 with SymPy, NumPy and Matplotlib; TeX Live with REVTeX
4.2, XeLaTeX and latexmk, the DejaVu fonts and the Latin Modern Math
font; pandoc. The program refuses to run unless it is started in the
folder \texttt{Revision\_SuperLoop} (it writes its data files into
\texttt{data/} of that folder).

\hypertarget{the-programs}{%
\section{3. The programs}\label{the-programs}}

\texttt{code/superloop-rs} (9 source files, 3305 lines) is a Rust
program with the commands \texttt{algebra}, \texttt{malcev},
\texttt{manifolds}, \texttt{superfluid}, \texttt{renorm} and
\texttt{all}. Its modules:

\begin{longtable}[]{@{}
  >{\raggedright\arraybackslash}p{(\columnwidth - 2\tabcolsep) * \real{0.5000}}
  >{\raggedright\arraybackslash}p{(\columnwidth - 2\tabcolsep) * \real{0.5000}}@{}}
\toprule\noalign{}
\begin{minipage}[b]{\linewidth}\raggedright
module
\end{minipage} & \begin{minipage}[b]{\linewidth}\raggedright
content
\end{minipage} \\
\midrule\noalign{}
\endhead
\bottomrule\noalign{}
\endlastfoot
\texttt{main.rs} & the commands; the guard on the working folder \\
\texttt{checks.rs} & a check is a numbered statement, PASS or FAIL; each
command writes its report
\texttt{data/sl\_\textless{}command\textgreater{}\_checks.txt}, whose
second line reads ``N of M checks passed''; the program exits with
status 1 if a check fails \\
\texttt{bioct.rs} & quaternions with complex coefficients; the
Cayley-Dickson product (with the sign of the octonions or of the split
octonions); the conjugations, the norm, the commutator, the associator,
the Jacobiator and the linearized Malcev expression; left and right
multiplications as real \(16\times16\) integer matrices; the
coefficients are exact Gaussian integers (128-bit integer real and
imaginary parts; an overflow stops the program) or double-precision
complex numbers; a deterministic pseudo-random generator (SplitMix64) \\
\texttt{exact.rs} & exact linear algebra on integer matrices: products
and commutators; spans and ranks over the rationals by fraction-free
Gaussian elimination (each row divided by the greatest common divisor of
its entries); Lie closures (commutators are added until the exact span
stops growing); the signature of a symmetric integer matrix by symmetric
elimination in exact rational arithmetic; ranks modulo the prime
\(2^{61}-1\) for the large systems \\
\texttt{algebra.rs} & command \texttt{algebra}: checks A, L, Q and C \\
\texttt{malcev.rs} & command \texttt{malcev}: the covariant families,
Jacobi and Malcev identities on all basis triples and 4-tuples; checks
M \\
\texttt{manifolds.rs} & command \texttt{manifolds}: checks N \\
\texttt{superfluid.rs} & command \texttt{superfluid}: the Dirac
matrices, the pair channels, the Bogoliubov-de Gennes matrices, the
grand potential, the gap equation; checks S \\
\texttt{renorm.rs} & command \texttt{renorm}: power counting, the
contact coupling, the bubble, the invariants; checks R \\
\end{longtable}

A rank computed modulo a prime can only be smaller than the rank over
the rationals, because vectors that are independent modulo the prime are
independent over the rationals; so a rank found equal to its largest
possible value is exact, and the dimension of a solution space computed
from such a rank is an upper bound. Every such bound in this folder is
matched by the same number of explicitly verified solutions. The library
crate \texttt{code/spacetimes-rs} of the repository supplies complex
numbers, the diagonalization of complex Hermitian matrices,
Gauss-Legendre nodes and weights, Brent's root finder and the writing of
data files.

\texttt{code/scripts/malcev\_groebner.py} (486 lines) computes the
Gröbner bases with the computer-algebra library SymPy (Buchberger's
algorithm over the rationals, Section 5). The figures are drawn by
\texttt{code/scripts/make\_figures.py} from the data files.

\textbf{How to read the reports.} Each check prints one line: its label,
the statement, and, after two dashes, the numbers that decide it. The
reports write formulas in plain text, with the names listed in the last
column of the table of symbols (Section 12). The following conventions
complete that column.

\begin{itemize}
\tightlist
\item
  \textbf{Elements of the algebras}: \texttt{x}, \texttt{y}, \texttt{z},
  \texttt{w} are the bi-octonions \(\mathbf x\), \(\mathbf y\),
  \(\mathbf z\), \(\mathbf w\); \texttt{aH}, \texttt{bS}, \texttt{cH},
  \texttt{dS}, \texttt{bS\textquotesingle{}} are the bi-quaternions
  \(\mathbf a\), \(\mathbf b\), \(\mathbf c\), \(\mathbf d\),
  \(\mathbf b'\) of the form \(\mathbf a+\mathbf b\,\mathbf e_4\) (the
  letter H recalls that the element lies in
  \(\mathbb C\otimes\mathbb H\), the letter S that it multiplies
  \(\mathbf e_4\) to give a spinor charge); \texttt{uL},
  \texttt{uL\textquotesingle{}}, \texttt{uL1}, \texttt{uL2} are
  \(\mathbf u\), \(\mathbf u'\), \(\mathbf u_1\), \(\mathbf u_2\)
  (Lorentz part); \texttt{tI}, \texttt{tI\textquotesingle{}} are
  \(\mathbf t\), \(\mathbf t'\) (isospin); \texttt{gL}, \texttt{hL} are
  \(\mathbf g\), \(\mathbf h\) (elements of \(SL(2,\mathbb C)\));
  \texttt{vI} is \(\mathbf v\) (a unit quaternion, an isospin rotation);
  \texttt{rP}, \texttt{rP\textquotesingle{}} are \(\mathbf r\),
  \(\mathbf r'\) (translations). \texttt{e0}, \ldots, \texttt{e7} are
  the units \(\mathbf e_0,\dots,\mathbf e_7\) (and \texttt{1} in
  \texttt{gL\ =\ 1\ +\ i\ e1\ +\ e2} is \(\mathbf e_0\)); \texttt{ie1},
  \ldots, \texttt{ie7} are \(i\mathbf e_1,\dots,i\mathbf e_7\);
  \texttt{e\_p}, \texttt{e\_j}, \texttt{e\_iota} are \(\mathbf e_p\),
  \(\mathbf e_j\), \(\mathbf e_\iota\).
\item
  \textbf{Operations}: a final \texttt{bar} (\texttt{xbar},
  \texttt{bSbar}, \texttt{gLbar}) is the conjugation
  \(\bar{\mathbf x}\); \texttt{\^{}*} the complex conjugation;
  \texttt{\^{}dag} the Hermitian conjugation; \texttt{\^{}T} the
  transpose; \texttt{\^{}-1} the inverse; \texttt{(x)} the tensor
  product \(\otimes\); \texttt{o} the composition \(\circ\);
  \texttt{{[}x,y{]}} the commutator and \texttt{{[}x,y,z{]}} the
  associator; \texttt{Jac} the Jacobiator; \texttt{L\_x}, \texttt{R\_x}
  the left and right multiplications \(L^{\rm LL}_{\mathbf x}\),
  \(R^{\rm LL}_{\mathbf x}\), and \texttt{L\_{[}e1,e4{]}} the left
  multiplication by the commutator \([\mathbf e_1,\mathbf e_4]\);
  \texttt{cross(a,\ b)} is the cross product of two complex
  three-component vectors; \texttt{\textbar{}...\textbar{}} an absolute
  value or the length of a vector; the dot in \texttt{kE.(kE+pE)} the
  Euclidean scalar product; \texttt{x\ -\textgreater{}\ y} a map;
  \texttt{x(y)} a product or a function value, as the context of the
  formula shows.
\item
  \textbf{Coefficients of the families}: \texttt{c\_tw}, \texttt{c\_sl},
  \texttt{c\_PP}, \texttt{c\_Ps}, \texttt{c\_sP}, \texttt{c\_TT},
  \texttt{c\_Ts}, \texttt{c\_sT} are \(c^{\rm tw}_{\rm LL}\), \ldots,
  \(c^{\rm sT}_{\rm LL}\); a complex one is split as
  \texttt{ctw\_re\ +\ i\ ctw\_im} into its real and imaginary parts (and
  so for the others); \texttt{csP} and \texttt{cTT} are real. In checks
  M-3, M-4, M-7 and M-9 a failing triple or 4-tuple is given by the
  positions of its elements in the basis of the family: positions 0, 1,
  2 are \(\mathbf e_1\), \(\mathbf e_2\), \(\mathbf e_3\) and 3, 4, 5
  are \(i\mathbf e_1\), \(i\mathbf e_2\), \(i\mathbf e_3\) (the Lorentz
  part); in \(\mathcal F^{\rm B}_{\rm LL}\) positions 6 to 9 are the
  translations \(\mathbf e_0\), \(i\mathbf e_1\), \(i\mathbf e_2\),
  \(i\mathbf e_3\), and in \(\mathcal F^{\rm C}_{\rm LL}\) positions 6
  to 8 are the isospin elements \(\mathbf e_1\), \(\mathbf e_2\),
  \(\mathbf e_3\); the last eight positions are the spinor charges
  \(\mathbf e_4,\dots,\mathbf e_7\),
  \(i\mathbf e_4,\dots,i\mathbf e_7\).
\item
  \textbf{The superfluid}: \texttt{dvec} is \(\vec d_{\rm LL}\) and
  \texttt{d\_iota} its components; \texttt{eps2} is
  \(\mathsf E_{\rm LL}\), \texttt{eps3} the Levi-Civita symbol
  \(\epsilon^{\rm LL}_{jj'j''}\); \texttt{E\^{}(+-)} stands for
  \(E^+_{\rm LL}\) and \(E^-_{\rm LL}\), and \texttt{E\^{}(+-)(Delta)}
  for these energies with the gap \texttt{Delta} in place of
  \(\Delta_{\rm LL}\); \texttt{Delta\_(+-)} are the gaps
  \(\Delta^\pm_{\rm LL}\); \texttt{Etilde\^{}(+-)} are
  \(\tilde E^\pm_{\rm LL}\); \texttt{eps} is \(\varepsilon_{\rm LL}\),
  \texttt{q} the momentum \(\vec q_{\rm LL}\) or its length;
  \texttt{Q\_LL(Kch)} is the electric charge of the two spinor charges
  that the pair operator with the matrix \texttt{Kch} annihilates;
  \texttt{0+} and \texttt{0-} are spin zero with positive and negative
  parity; \texttt{spin(spin+1)} in check Q-3 is the spin quantum number
  times itself plus one, the value 0, 2 or \(\tfrac34\) of the Casimir
  for the spins 0, 1 and \(\tfrac12\).
\item
  \textbf{Abbreviations}: \texttt{BdG} is Bogoliubov-de Gennes,
  \texttt{BCS} Bardeen-Cooper-Schrieffer; \texttt{M\_2(C)} the complex
  \(2\times2\) matrices; \texttt{R\^{}8}, \texttt{C\^{}3},
  \texttt{C\^{}8} the real or complex vector spaces of that dimension;
  \texttt{Im\ H}, \texttt{Im\ O} and \texttt{C\ (x)\ H} are
  \({\rm Im}\,\mathbb H\), \({\rm Im}\,\mathbb O\) and
  \(\mathbb C\otimes\mathbb H\); the groups and algebras
  (\texttt{SL(2,C)}, \texttt{so(3,1)}, \texttt{g2(C)},
  \texttt{SU(2)\_v}, \texttt{U(1)\_F\_LL}, \texttt{Z2}, \ldots) are
  written as in the text; \texttt{4x4} means four by four; ``exact'' in
  a check means integer or rational arithmetic, ``double precision''
  floating-point arithmetic with 53-bit mantissas.
\item
  \textbf{Comments of the source code} also call the quaternion units
  \(\mathbf e_1\), \(\mathbf e_2\), \(\mathbf e_3\) \texttt{i\_q},
  \texttt{j\_q}, \texttt{k\_q}, and the unit \(\mathbf e_4\) of the
  Cayley-Dickson construction \texttt{l}; these names do not occur in
  the reports, the data files or the figures.
\end{itemize}

\hypertarget{command-algebra-the-bi-octonions-the-super-loop-and-its-operators}{%
\section{\texorpdfstring{4. Command \texttt{algebra}: the bi-octonions,
the super-loop and its
operators}{4. Command algebra: the bi-octonions, the super-loop and its operators}}\label{command-algebra-the-bi-octonions-the-super-loop-and-its-operators}}

\textbf{Quaternions and bi-quaternions.} The units \(\mathbf e_0=1\),
\(\mathbf e_1\), \(\mathbf e_2\), \(\mathbf e_3\) obey
\(\mathbf e_j^2=-\mathbf e_0\) and
\(\mathbf e_1\mathbf e_2=-\mathbf e_2\mathbf e_1=\mathbf e_3\) with its
cyclic permutations; the bi-quaternions \(\mathbf a\), \(\mathbf b\),
\(\mathbf c\), \(\mathbf d\) are their complex combinations. The map
\(M_{\rm LL}(\mathbf e_0)=\sigma^{\rm LL}_0\),
\(M_{\rm LL}(\mathbf e_j)=-i\sigma^{\rm LL}_j\) (complex-linear) is an
isomorphism onto the complex \(2\times2\) matrices, with
\(\det M_{\rm LL}(\mathbf b)=N_{\rm LL}(\mathbf b)\).

\textbf{The bi-octonions} are the pairs
\(\mathbf x=\mathbf a+\mathbf b\,\mathbf e_4\) with the Cayley-Dickson
product
\[(\mathbf a+\mathbf b\,\mathbf e_4)(\mathbf c+\mathbf d\,\mathbf e_4)=(\mathbf a\mathbf c-\bar{\mathbf d}\mathbf b)+(\mathbf d\mathbf a+\mathbf b\bar{\mathbf c})\,\mathbf e_4 ,\]
units \(\mathbf e_5=\mathbf e_1\mathbf e_4\),
\(\mathbf e_6=\mathbf e_2\mathbf e_4\),
\(\mathbf e_7=\mathbf e_3\mathbf e_4\), and the special cases
\(\mathbf a(\mathbf b\,\mathbf e_4)=(\mathbf b\mathbf a)\,\mathbf e_4\),
\((\mathbf b\,\mathbf e_4)\mathbf c=(\mathbf b\bar{\mathbf c})\,\mathbf e_4\),
\((\mathbf b\,\mathbf e_4)(\mathbf d\,\mathbf e_4)=-\bar{\mathbf d}\mathbf b\).
The conjugation \(\bar{\mathbf x}\) changes the signs of the
coefficients of \(\mathbf e_1,\dots,\mathbf e_7\); \(\mathbf x^*\)
conjugates every complex coefficient;
\(\mathbf x^\dagger=(\bar{\mathbf x})^*\). With the coefficients
\(\xi^{\rm LL}_p(\mathbf x)\), the norm is
\(N_{\rm LL}(\mathbf x)=\sum_p\xi^{\rm LL}_p(\mathbf x)^2\), so that
\(\mathbf x\bar{\mathbf x}=N_{\rm LL}(\mathbf x)\,\mathbf e_0\), and its
polarization is
\(N_{\rm LL}(\mathbf x,\mathbf y)=\sum_p\xi^{\rm LL}_p(\mathbf x)\xi^{\rm LL}_p(\mathbf y)\).
The program builds the integer table of the 64 products of the units
(data file \texttt{sl\_multiplication\_table.csv}) and computes
everything else from it.

\textbf{Where associativity enters (checks A).} The associator is
\([\mathbf x,\mathbf y,\mathbf z]=(\mathbf x\mathbf y)\mathbf z-\mathbf x(\mathbf y\mathbf z)\)
and the Jacobiator of the commutator
\([\mathbf x,\mathbf y]=\mathbf x\mathbf y-\mathbf y\mathbf x\) is
\(\mathrm{Jac}_{\rm LL}(\mathbf x,\mathbf y,\mathbf z)=[[\mathbf x,\mathbf y],\mathbf z]+[[\mathbf y,\mathbf z],\mathbf x]+[[\mathbf z,\mathbf x],\mathbf y]\).
Expanding its twelve products gives, for any product,
\[\begin{aligned}\mathrm{Jac}_{\rm LL}(\mathbf x,\mathbf y,\mathbf z)={}&[\mathbf x,\mathbf y,\mathbf z]+[\mathbf y,\mathbf z,\mathbf x]+[\mathbf z,\mathbf x,\mathbf y]\\&-[\mathbf y,\mathbf x,\mathbf z]-[\mathbf x,\mathbf z,\mathbf y]-[\mathbf z,\mathbf y,\mathbf x].\end{aligned}\]
For operators (A-6, on 50 random triples of \(4\times4\) integer
matrices with entries from \(-5\) to \(5\)) every associator vanishes,
so the Jacobi identity holds. For the bi-octonions the associator is
alternating (A-1, all 512 ordered triples of units), so the six terms
are equal and
\(\mathrm{Jac}_{\rm LL}(\mathbf x,\mathbf y,\mathbf z)=6[\mathbf x,\mathbf y,\mathbf z]\)
(A-4); exactly 168 ordered triples of units have a non-zero associator
(A-2; data file \texttt{sl\_associators.csv}). The norm is
multiplicative and the three Moufang identities
\(\mathbf z(\mathbf x(\mathbf z\mathbf y))=((\mathbf z\mathbf x)\mathbf z)\mathbf y\),
\(((\mathbf y\mathbf z)\mathbf x)\mathbf z=\mathbf y(\mathbf z(\mathbf x\mathbf z))\)
and
\((\mathbf z\mathbf x)(\mathbf y\mathbf z)=(\mathbf z(\mathbf x\mathbf y))\mathbf z\)
hold (A-3, exact, on 300 random triples with Gaussian-integer
coefficients whose real and imaginary parts run from \(-3\) to \(3\)).
So the unit bi-octonions
\(\mathbb L_{\rm LL}=\{\mathbf w\in\mathbb B_{\rm LL}:N_{\rm LL}(\mathbf w)=1\}\)
form a Moufang loop, with the identity \(\mathbf e_0\) and the inverse
\(\bar{\mathbf w}\). The commutator algebra of the imaginary
bi-octonions, \(\mathfrak m_{\rm LL}\), obeys the linearized Malcev
identity
\[\begin{aligned}&\mathrm{Jac}_{\rm LL}(\mathbf x,\mathbf y,[\mathbf w,\mathbf z])+\mathrm{Jac}_{\rm LL}(\mathbf w,\mathbf y,[\mathbf x,\mathbf z])\\&\quad=[\mathrm{Jac}_{\rm LL}(\mathbf x,\mathbf y,\mathbf z),\mathbf w]+[\mathrm{Jac}_{\rm LL}(\mathbf w,\mathbf y,\mathbf z),\mathbf x]\end{aligned}\]
on all \(7^4=2401\) 4-tuples of imaginary units (A-5); both sides are
complex-multilinear, so this proves it. \textbf{Theorem R} (A-8): a
linear map \(\varrho_{\rm LL}\) of \(\mathfrak m_{\rm LL}\) into
operators with
\(\varrho_{\rm LL}([\mathbf x,\mathbf y])=\varrho_{\rm LL}(\mathbf x)\varrho_{\rm LL}(\mathbf y)-\varrho_{\rm LL}(\mathbf y)\varrho_{\rm LL}(\mathbf x)\)
maps every Jacobiator to the Jacobiator of operators, which is zero; the
program computes the Jacobiators of all triples of the real basis
\(\mathbf e_p\), \(i\mathbf e_p\) (\(p=1,\dots,7\)) and finds that they
span a space of real dimension 14, all of \(\mathfrak m_{\rm LL}\);
hence \(\varrho_{\rm LL}=0\). Check A-7 shows that the left
multiplications \(L^{\rm LL}_{\mathbf x}\) do not represent
\(\mathfrak m_{\rm LL}\):
\([L^{\rm LL}_{\mathbf e_1},L^{\rm LL}_{\mathbf e_4}]\neq L^{\rm LL}_{[\mathbf e_1,\mathbf e_4]}\).

\textbf{The Lorentz part and the spinor charges (checks L).} With
\(J^j_{\rm LL}=\tfrac12\mathbf e_j\) and
\(K^j_{\rm LL}=\tfrac12 i\,\mathbf e_j\),
\[[J^j_{\rm LL},J^{j'}_{\rm LL}]=\textstyle\sum_{j''}\epsilon^{\rm LL}_{jj'j''}J^{j''}_{\rm LL},\quad [J^j_{\rm LL},K^{j'}_{\rm LL}]=\textstyle\sum_{j''}\epsilon^{\rm LL}_{jj'j''}K^{j''}_{\rm LL},\quad [K^j_{\rm LL},K^{j'}_{\rm LL}]=-\textstyle\sum_{j''}\epsilon^{\rm LL}_{jj'j''}J^{j''}_{\rm LL}\]
(L-1). The Killing form
\(\mathrm{Kil}_{\rm LL}(\mathbf x,\mathbf y)=\mathrm{tr}(\mathrm{ad}\,\mathbf x\,\mathrm{ad}\,\mathbf y)\)
of \(\mathfrak l_{\rm LL}\), with \(\mathrm{ad}\,\mathbf x\) the matrix
of \(\mathbf y\mapsto[\mathbf x,\mathbf y]\) in the basis
\(\mathbf e_1\), \(\mathbf e_2\), \(\mathbf e_3\), \(i\mathbf e_1\),
\(i\mathbf e_2\), \(i\mathbf e_3\), has the signature (3, 3), with
\(\mathrm{Kil}_{\rm LL}(\mathbf e_1,\mathbf e_1)=-16\) and
\(\mathrm{Kil}_{\rm LL}(i\mathbf e_1,i\mathbf e_1)=16\) (L-2). A Lorentz
element \(\mathbf u\) acts on a spinor charge by
\[[\mathbf u,\mathbf b\,\mathbf e_4]=2(\mathbf b\mathbf u)\,\mathbf e_4=-2\,(W_{\rm LL}(\mathbf u)\mathbf b)\,\mathbf e_4,\qquad W_{\rm LL}(\mathbf u)\mathbf b=-\mathbf b\mathbf u ;\]
\(W_{\rm LL}\) is a representation,
\([W_{\rm LL}(\mathbf u_1),W_{\rm LL}(\mathbf u_2)]=W_{\rm LL}([\mathbf u_1,\mathbf u_2])\),
but the inner action \(\rho_{\rm LL}=-2W_{\rm LL}\) obeys
\([\rho_{\rm LL}(\mathbf u_1),\rho_{\rm LL}(\mathbf u_2)]=-2\rho_{\rm LL}([\mathbf u_1,\mathbf u_2])\)
(L-3). Two spinor charges have the bracket
\([\mathbf b\,\mathbf e_4,\mathbf b'\mathbf e_4]=\bar{\mathbf b}\mathbf b'-\bar{\mathbf b}'\mathbf b\),
an element of \(\mathfrak l_{\rm LL}\); these brackets span
\(\mathfrak l_{\rm LL}\) (L-4, rank 6). If \(\mathbf u\) acts on the
spinor charges by \(\Theta_{\rm LL}W_{\rm LL}(\mathbf u)\), then
\([\mathbf u_2,\mathbf b\,\mathbf e_4]=\Theta_{\rm LL}(W_{\rm LL}(\mathbf u_2)\mathbf b)\,\mathbf e_4\)
and
\[\mathrm{Jac}_{\rm LL}(\mathbf u_1,\mathbf u_2,\mathbf b\,\mathbf e_4)=(\Theta_{\rm LL}-\Theta_{\rm LL}^2)\,(W_{\rm LL}([\mathbf u_1,\mathbf u_2])\mathbf b)\,\mathbf e_4 ,\]
because the Jacobiator is
\([[\mathbf u_1,\mathbf u_2],\mathbf b\,\mathbf e_4]-[\mathbf u_1,[\mathbf u_2,\mathbf b\,\mathbf e_4]]+[\mathbf u_2,[\mathbf u_1,\mathbf b\,\mathbf e_4]]\);
for the super-loop, \(\Theta_{\rm LL}=-2\), it is
\(-6(W_{\rm LL}([\mathbf u_1,\mathbf u_2])\mathbf b)\,\mathbf e_4\)
(L-5, exact on all 288 basis triples), and
\(\mathrm{Jac}_{\rm LL}(\mathbf e_4,\mathbf e_5,\mathbf e_6)=-12\,\mathbf e_7\).
With the Lie action \(\Theta_{\rm LL}=1\) and the same bracket of spinor
charges the Jacobi identity fails on triples of spinor charges (L-6;
Theorem L).

\textbf{Derivations and automorphisms.} The Lorentz and isospin
derivations
\[D_{\rm LL}(\mathbf u)(\mathbf a+\mathbf b\,\mathbf e_4)=[\mathbf u,\mathbf a]+(-\mathbf b\mathbf u)\,\mathbf e_4,\qquad D^{\rm v}_{\rm LL}(\mathbf t)(\mathbf a+\mathbf b\,\mathbf e_4)=(\mathbf t\mathbf b)\,\mathbf e_4\]
are derivations,
\(\mathcal D_{\rm LL}(\mathbf x\mathbf y)=(\mathcal D_{\rm LL}\mathbf x)\mathbf y+\mathbf x(\mathcal D_{\rm LL}\mathbf y)\)
(L-7, on 100 random pairs), they commute, and the inner action factors
as \([\mathbf u,\cdot\,]=D_{\rm LL}(\mathbf u)\circ\Theta_{\rm LL}\)
with the twist charge \(\Theta_{\rm LL}\) equal to 1 on
\(\mathfrak l_{\rm LL}\) and \(-2\) on \(\mathfrak s_{\rm LL}\). The
automorphisms
\(\varphi^{\rm LL}_{\mathbf g,\mathbf v}(\mathbf a+\mathbf b\,\mathbf e_4)=\mathbf g\mathbf a\bar{\mathbf g}+(\mathbf v\mathbf b\bar{\mathbf g})\,\mathbf e_4\)
are checked for \(\mathbf g=\mathbf e_0+i\mathbf e_1+\mathbf e_2\) and
\(\mathbf v=\mathbf e_3\) (L-9). \textbf{All derivations} (L-8): a real
\(16\times16\) matrix \(\mathcal D_{\rm LL}\) is a derivation if the
condition above holds for all 256 pairs of real basis elements
\(\mathbf e_p\), \(i\mathbf e_p\); these are 4096 linear equations for
its 256 entries, of rank 228 (modulo the prime), so there are at most 28
independent derivations; the inner derivations
\([L^{\rm LL}_{\mathbf x},L^{\rm LL}_{\mathbf y}]+[L^{\rm LL}_{\mathbf x},R^{\rm LL}_{\mathbf y}]+[R^{\rm LL}_{\mathbf x},R^{\rm LL}_{\mathbf y}]\),
each checked to be a derivation on random pairs, span 28 dimensions: the
derivations form \(\mathfrak g_2(\mathbb C)\), of real dimension 28.

\textbf{Operators and eigenvalues (checks Q).} The isospin generators
\(T^\iota_{\rm LL}=i\,D^{\rm v}_{\rm LL}(\mathbf e_\iota/2)\) obey
\([T^1_{\rm LL},T^2_{\rm LL}]=iT^3_{\rm LL}\) and its cyclic
permutations; the electric charge is \(Q_{\rm LL}=T^3_{\rm LL}\), with
the eigenvalues 0 (four times) on \(\mathbb C\otimes\mathbb H\) and
\(\pm\tfrac12\) (twice each) on the spinor charges (Q-1); the isospin
Casimir \((T_{\rm LL})^2=\sum_\iota(T^\iota_{\rm LL})^2\) is 0 and
\(\tfrac34\) (Q-2). The spin Casimir
\(\mathrm{Cas}_{\rm LL}=-\sum_jD_{\rm LL}(J^j_{\rm LL})^2\) is 0 on
\(\mathbf e_0\), 2 on the Lorentz part and \(\tfrac34\) on the spinor
charges (Q-3, exact: \(4\,\mathrm{Cas}_{\rm LL}\) is an integer matrix).
Chirality: \(D_{\rm LL}(K^j_{\rm LL})=i\,D_{\rm LL}(J^j_{\rm LL})\) on
every bi-octonion (Q-4). Spinor number: \(iD_{\rm LL}(J^3_{\rm LL})\)
has the eigenvalues \(-1,0,0,1\) on \(\mathbb C\otimes\mathbb H\) and
\(\pm\tfrac12\) on the spinor charges, so
\(\exp(2\pi D_{\rm LL}(J^3_{\rm LL}))\) is \(+1\) and \(-1\) (Q-5).
Conjugation by an element of \(SL(2,\mathbb C)\) inside the loop:
\(\mathbf g(\mathbf b\,\mathbf e_4)\bar{\mathbf g}=(\mathbf b\,\mathbf g^2)\,\mathbf e_4\),
with both orders of the products (Q-6, on 50 random \(\mathbf b\)). The
eigenvalues of Q-1, Q-2 and Q-5 are computed in double precision (data
file \texttt{sl\_quantum\_numbers.csv}).

\textbf{The inversion element (checks C).}
\(\mathbf e_4^2=-\mathbf e_0\), \(\mathbf e_4^{-1}=-\mathbf e_4\) and
\(\mathbf e_4\,\mathbf a\,\mathbf e_4^{-1}=\bar{\mathbf a}\) (C-1, 100
random \(\mathbf a\)): on \(SL(2,\mathbb C)\) the inversion, on the
Hermitian bi-quaternions
\(x^0_{\rm LL}\mathbf e_0+i\sum_jx^j_{\rm LL}\mathbf e_j\) the parity.
The union of \(SL(2,\mathbb C)\) and \(SL(2,\mathbb C)\,\mathbf e_4\) is
closed:
\((\mathbf g\,\mathbf e_4)(\mathbf h\,\mathbf e_4)=-\bar{\mathbf h}\mathbf g\),
\(\mathbf g(\mathbf h\,\mathbf e_4)=(\mathbf h\mathbf g)\,\mathbf e_4\),
\((\mathbf g\,\mathbf e_4)\mathbf h=(\mathbf g\bar{\mathbf h})\,\mathbf e_4\)
(C-2). The inversion is not an automorphism of \(SL(2,\mathbb C)\):
\((\mathbf e_1\mathbf e_2)^{-1}=-\mathbf e_3\), but
\(\mathbf e_1^{-1}\mathbf e_2^{-1}=+\mathbf e_3\) (C-3); and
\(\mathbf x\mapsto\mathbf e_4\mathbf x\mathbf e_4^{-1}\) is not an
automorphism of \(\mathbb B_{\rm LL}\) (C-4).

\begin{verbatim}
Revision_SuperLoop, command `algebra`: bi-octonions, the Lorentz super-loop, its Malcev algebra, its operators and eigenvalues
27 of 27 checks passed

[PASS] A-1  the associator [x,y,z] = (xy)z - x(yz) of the bi-octonions B_LL is alternating (it changes sign under every exchange of two of its arguments) on all 512 ordered triples of the units e0..e7: B_LL is an ALTERNATIVE algebra
[PASS] A-2  B_LL is NOT associative: exactly 168 of the 512 ordered triples of units have a non-zero associator (the triples of distinct imaginary units outside the 7 quaternionic triples); example [e1,e2,e4] = (e1 e2) e4 - e1 (e2 e4) = 2 e7  -- non-zero associators: 168; coefficients of [e1,e2,e4] on e0..e7: [0, 0, 0, 0, 0, 0, 0, 2]
[PASS] A-3  on 300 random bi-octonions x, y, z with Gaussian-integer coefficients the complex norm is multiplicative, N_LL(xy) = N_LL(x) N_LL(y), and the three Moufang identities z(x(zy)) = ((zx)z)y, ((yz)x)z = y(z(xz)), (zx)(yz) = (z(xy))z hold exactly: the unit bi-octonions form a Moufang LOOP, the super-loop L_LL
[PASS] A-4  for the commutator [x,y] = xy - yx of B_LL the Jacobiator Jac(x,y,z) = [[x,y],z] + [[y,z],x] + [[z,x],y] equals six times the associator on all 512 triples of units, Jac(x,y,z) = 6 [x,y,z]; it is non-zero on 168 triples: the Jacobi identity fails exactly where associativity fails  -- non-zero Jacobiators: 168
[PASS] A-5  the imaginary bi-octonions m_LL = C (x) Im O with the commutator satisfy the linearized Malcev identity Jac(x,y,[w,z]) + Jac(w,y,[x,z]) = [Jac(x,y,z),w] + [Jac(w,y,z),x] on all 2401 4-tuples of imaginary units: m_LL is a Malcev algebra that is not a Lie algebra
[PASS] A-6  for 50 random triples of 4x4 integer matrices (operators, whose product is associative) the Jacobiator of the commutator vanishes exactly; this is the step that every no-go proof takes when it treats the symmetry generators as operators
[PASS] A-7  left multiplication x -> L_x is not a representation of the Malcev algebra m_LL: [L_e1, L_e4] - L_[e1,e4] is a non-zero matrix
[PASS] A-8  Theorem R: the Jacobiators of the real basis elements e_p, i e_p (p = 1..7) of m_LL span all of m_LL (real rank 14, exact); because operators obey the Jacobi identity (A-6), every linear map from m_LL to operators on a vector space that turns commutators into commutators maps every Jacobiator to 0 and therefore vanishes: m_LL has no non-zero realization by operators  -- rank of the span of the Jacobiators: 14
[PASS] L-1  the Lorentz part l_LL = C (x) Im H (real basis J_LL^j = e_j/2, K_LL^j = i e_j/2, j = 1,2,3) is a Lie subalgebra of m_LL isomorphic to so(3,1): [J_LL^j, J_LL^j'] = sum_j'' eps3_(j j' j'') J_LL^j'', [J_LL^j, K_LL^j'] = sum_j'' eps3_(j j' j'') K_LL^j'', [K_LL^j, K_LL^j'] = -sum_j'' eps3_(j j' j'') J_LL^j''; the unit bi-quaternions SL(2,C) = Spin(3,1)^+ form a subgroup of the super-loop L_LL
[PASS] L-2  the Killing form Killing(x,y) = tr(ad x ad y) of l_LL (6x6, exact) is non-degenerate with signature (3 positive, 3 negative), negative on the rotations and positive on the boosts, as for so(3,1)  -- signature (3, 3, 0), Killing(e1,e1) = -16, Killing(i e1,i e1) = 16
[PASS] L-3  a Lorentz element uL of l_LL acts on a spinor charge bS e4 of s_LL = (C (x) H) e4 by [uL, bS e4] = 2 (bS uL) e4 = -2 (W_LL(uL) bS) e4, where W_LL(uL) bS = -bS uL is a Lie representation (the Weyl spinor); the inner action rho_LL(uL) = -2 W_LL(uL) obeys [rho_LL(uL1), rho_LL(uL2)] = -2 rho_LL([uL1,uL2]): it is NOT a representation (the twist)
[PASS] L-4  the bracket of two spinor charges lies in the Lorentz part, [bS e4, bS' e4] = bSbar bS' - bS'bar bS in l_LL, and these brackets span all of l_LL (rank 6): spinor charges close into Lorentz generators  -- rank 6
[PASS] L-5  Jacobiators by sector: Jac(uL1, uL2, bS e4) = -6 (W_LL([uL1,uL2]) bS) e4 = (Theta_LL - Theta_LL^2) (W_LL([uL1,uL2]) bS) e4 with the twist charge Theta_LL = -2 on s_LL (exact on all 288 triples of basis elements); Jac(s_LL, s_LL, s_LL) (e.g. Jac(e4,e5,e6) = -12 e7) and Jac(l_LL, s_LL, s_LL) (e.g. Jac(e1,e4,e6)) are non-zero as well
[PASS] L-6  with the Lie action (twist charge 1 instead of -2) and the same non-zero bracket [s_LL, s_LL] in l_LL, the Jacobi identity fails on triples of spinor charges: no Lie algebra extends so(3,1) by bosonic spinor charges that close into so(3,1) (Theorem L; the classification G-1, G-5, G-10, M-3 shows that every Lie algebra of the covariant families has no Lorentz part in [s_LL, s_LL])
[PASS] L-7  the Lorentz action by automorphisms, D_LL(uL)(aH + bS e4) = [uL, aH] + (-bS uL) e4, and the isospin action Dv_LL(tI)(aH + bS e4) = (tI bS) e4 (tI in Im H) are derivations of B_LL (100 random pairs), they commute, the isospin derivations vanish on C (x) H, and the inner action factors as [uL, .] = D_LL(uL) o Theta_LL with the twist charge Theta_LL = 1 on l_LL and -2 on s_LL
[PASS] L-8  the real derivations of B_LL form a Lie algebra of real dimension 28 (the complex Lie algebra g2(C)): the linear system has rank 228 for 256 unknowns (at most 28 solutions), and the inner derivations [L_x,L_y] + [L_x,R_y] + [R_x,R_y] are verified derivations spanning 28 dimensions  -- rank of the system 228, explicit derivations span 28
[PASS] L-9  varphi_(gL,vI)(aH + bS e4) = gL aH gLbar + (vI bS gLbar) e4 is an automorphism of B_LL for the unit bi-quaternions gL = 1 + i e1 + e2 (a Lorentz transformation) and vI = e3 (an isospin rotation), on 100 random pairs: (SL(2,C) x SL(2,C))/Z2 = SO(4,C) inside G2(C)
[PASS] Q-1  the electric charge Q_LL = T_LL^3 = i Dv_LL(e3/2) (the third isospin generator, an automorphism of B_LL) is a Hermitian operator on B_LL with the eigenvalues 0 (four times) on C (x) H (the Lorentz part and the scalar) and +1/2, -1/2 (twice each) on the spinor charges (double precision)  -- on C (x) H: +0.000000 +0.000000 +0.000000 +0.000000; on s_LL: -0.500000 -0.500000 +0.500000 +0.500000
[PASS] Q-2  the isospin Casimir T_LL^2 = sum_iota (T_LL^iota)^2 = -sum_iota Dv_LL(e_iota/2)^2 equals 0 on C (x) H and 3/4 = (1/2)(1/2 + 1) on the spinor charges: the spinor charges form an isospin doublet (double precision)  -- on s_LL: +0.750000 +0.750000 +0.750000 +0.750000
[PASS] Q-3  the Lorentz Casimir Cas_LL = -sum_j D_LL(J_LL^j)^2 (exact; 4 Cas_LL is an integer matrix) has the value spin(spin+1) = 0 on the scalar e0, 2 (spin 1) on the Lorentz part C (x) Im H and 3/4 (spin 1/2) on the spinor charges
[PASS] Q-4  chirality: on every bi-octonion the boost generators act as i times the rotation generators, D_LL(K_LL^j) = i D_LL(J_LL^j); with respect to the complex structure of B_LL every sector has one chirality, type (spin, 0): l_LL = (1,0), s_LL = (1/2,0); the complex-conjugate super-loop has the other chirality (0, spin)
[PASS] Q-5  spinor number: the rotation generator i D_LL(J_LL^3) has the integer eigenvalues {-1, 0, 0, 1} on C (x) H and the half-odd eigenvalues {-1/2, -1/2, 1/2, 1/2} on the spinor charges, so the rotation by 2 pi, exp(2 pi D_LL(J_LL^3)), is +1 on the Lorentz part and -1 on the spinor charges (double precision)  -- on C (x) H: -1.000000 +0.000000 +0.000000 +1.000000; on s_LL: -0.500000 -0.500000 +0.500000 +0.500000
[PASS] Q-6  inside the loop, conjugation by an element gL of SL(2,C) acts on a spinor charge as gL (bS e4) gLbar = (bS gL^2) e4 (both bracketings agree), whereas the Lorentz automorphism acts as bS -> bS gL^-1: the loop's inner action turns a spinor twice as fast and backwards (twist charge -2)
[PASS] C-1  the unit e4 of the super-loop (e4^2 = -1, e4^-1 = -e4) acts on bi-quaternions by e4 aH e4^-1 = aHbar (100 random aH): on SL(2,C) it is the inversion gL -> gL^-1, and on Minkowski space R^(1,3)_LL (Hermitian bi-quaternions) it is the parity (x_LL^0, x_LL^1, x_LL^2, x_LL^3) -> (x_LL^0, -x_LL^1, -x_LL^2, -x_LL^3)
[PASS] C-2  the union of SL(2,C) and SL(2,C) e4 is closed under the product: (gL e4)(hL e4) = -hLbar gL, gL (hL e4) = (hL gL) e4, (gL e4) hL = (gL hLbar) e4; it is a Moufang loop (a Chein-type double of the Lorentz group) inside L_LL, and it is not associative because SL(2,C) is not commutative
[PASS] C-3  inversion is not an automorphism of SL(2,C): for gL = e1, hL = e2, (gL hL)^-1 = -e3 but gL^-1 hL^-1 = +e3; hence no operator on any vector space can implement gL -> gL^-1 by conjugation in a faithful representation (conjugation is multiplicative): the extension by e4 exists only as a loop
[PASS] C-4  the inner map x -> e4 x e4^-1 is not an automorphism of B_LL (it fails on random pairs): inner maps of the loop are not symmetries of its multiplication, another place where associativity is missing
\end{verbatim}

\hypertarget{the-classification-of-the-covariant-brackets}{%
\section{5. The classification of the covariant
brackets}\label{the-classification-of-the-covariant-brackets}}

\textbf{The families.} The bracket of \(\mathfrak l_{\rm LL}\) with
itself is the Lorentz bracket, and
\([\mathbf u,\mathbf r]=\mathbf u\mathbf r+\mathbf r\mathbf u^\dagger\)
on the translations \(\mathbf r=\mathbf r^\dagger\) (Hermitian
bi-quaternions, the sector \(\mathfrak p_{\rm LL}\)). The other brackets
are \[\begin{aligned}
[\mathbf u,\mathbf b\,\mathbf e_4]&=c^{\rm tw}_{\rm LL}(W_{\rm LL}(\mathbf u)\mathbf b)\,\mathbf e_4, &
[\mathbf r,\mathbf r']&=c^{\rm PP}_{\rm LL}(\mathbf r\bar{\mathbf r}'-\mathbf r'\bar{\mathbf r}),\\
[\mathbf r,\mathbf b\,\mathbf e_4]&=c^{\rm Ps}_{\rm LL}(\mathbf b^*\bar{\mathbf r})\,\mathbf e_4, &
[\mathbf t,\mathbf t']&=c^{\rm TT}_{\rm LL}(\mathbf t\mathbf t'-\mathbf t'\mathbf t),\\
[\mathbf t,\mathbf b\,\mathbf e_4]&=c^{\rm Ts}_{\rm LL}(\mathbf t\mathbf b)\,\mathbf e_4, &
[\mathbf b\,\mathbf e_4,\mathbf b'\mathbf e_4]&=c^{\rm sl}_{\rm LL}(\bar{\mathbf b}\mathbf b'-\bar{\mathbf b}'\mathbf b)+c^{\rm sP}_{\rm LL}\,i(\bar{\mathbf b}\mathbf b'^*-\bar{\mathbf b}'\mathbf b^*)\\
&&&\quad+\mathrm{Re}\big[c^{\rm sT}_{\rm LL}(\mathbf b\bar{\mathbf b}'-\mathbf b'\bar{\mathbf b})\big],
\end{aligned}\] a term being present only if its target sector is:
\(\mathcal F^{\rm A}_{\rm LL}\) has the sectors
\(\mathfrak l_{\rm LL}\oplus\mathfrak s_{\rm LL}\) (real dimension 14),
\(\mathcal F^{\rm B}_{\rm LL}\) adds \(\mathfrak p_{\rm LL}\) (18),
\(\mathcal F^{\rm C}_{\rm LL}\) adds the isospin generators
\(\mathbf t\in{\rm Im}\,\mathbb H\), \(\mathfrak{su}(2)\) (17). The real
part in the last term is taken coefficient by coefficient;
\(c^{\rm sP}_{\rm LL}\) and \(c^{\rm TT}_{\rm LL}\) are real, the others
complex. The bi-octonions are \(\mathcal F^{\rm A}_{\rm LL}\) with
\(c^{\rm tw}_{\rm LL}=-2\), \(c^{\rm sl}_{\rm LL}=1\) (M-1).

\textbf{Why these are the most general covariant brackets.} After
complexification the Lorentz algebra is two commuting copies of
\(\mathfrak{sl}(2)\), and a complex representation is labelled by two
spins. The spinor charges are \((\tfrac12,0)\) times an isodoublet, the
Lorentz part \((1,0)\), the translations \((\tfrac12,\tfrac12)\), the
isospin generators \((0,0)\) times an isotriplet, together with the
complex conjugates of these representations. A covariant real bilinear
bracket is a sum of covariant complex-bilinear maps between these
representations and their conjugates, and the number of such maps is the
number of times the target occurs in the tensor product of the sources:
\((\tfrac12,0)\) occurs once in \((1,0)\otimes(\tfrac12,0)\) and never
in \((0,1)\otimes(\tfrac12,0)\) {[}\(c^{\rm tw}_{\rm LL}\){]}; \((1,0)\)
times the isosinglet occurs once in the antisymmetric square of
\((\tfrac12,0)\) times the isodoublet and never in
\((\tfrac12,0)\otimes(0,\tfrac12)\) {[}\(c^{\rm sl}_{\rm LL}\){]}; the
antisymmetric square of \((\tfrac12,\tfrac12)\) is \((1,0)\oplus(0,1)\)
{[}\(c^{\rm PP}_{\rm LL}\){]}; \((\tfrac12,0)\) occurs once in
\((\tfrac12,\tfrac12)\otimes(0,\tfrac12)\) and never in
\((\tfrac12,\tfrac12)\otimes(\tfrac12,0)\), so the bracket of a
translation with a spinor charge is antilinear in the spinor
{[}\(c^{\rm Ps}_{\rm LL}\){]}; \((\tfrac12,\tfrac12)\) times the
isosinglet occurs once in \((\tfrac12,0)\otimes(0,\tfrac12)\) times
isodoublet \(\otimes\) isodoublet, and one real combination is
antisymmetric {[}\(c^{\rm sP}_{\rm LL}\){]}; the isotriplet occurs once
in the antisymmetric square of the spinor charges
{[}\(c^{\rm sT}_{\rm LL}\){]}; the isodoublet occurs once in isotriplet
\(\otimes\) isodoublet {[}\(c^{\rm Ts}_{\rm LL}\){]}. All other brackets
between different sectors vanish by covariance.

\textbf{The Gröbner-basis method} (\texttt{malcev\_groebner.py}). The
Jacobi identity and the linearized Malcev identity are multilinear, so
they hold for all elements if they hold for basis elements. Evaluated on
elements with integer coordinates, each component of an identity is a
polynomial in the real and imaginary parts of the coefficients (the
variables \texttt{ctw\_re}, \texttt{ctw\_im}, \texttt{csl\_re}, \ldots,
degree 2 for Jacobi, 3 for Malcev). The script evaluates the identities
on random integer tuples (each element with four non-zero coordinates
from \(-2\), \(-1\), \(1\), \(2\) at random positions) and inserts the
resulting polynomials into an exact rational linear span, until twelve
consecutive tuples add nothing new (at most 400 tuples). It then
computes the reduced Gröbner basis of the ideal generated by the span,
in the graded reverse lexicographic order (grevlex),
\(\mathcal I^{\rm Jac}_{\rm LL}\) for Jacobi and
\(\mathcal I^{\rm Mal}_{\rm LL}\) for Malcev. Every polynomial found is
a true consequence of the identity, so every polynomial of the ideal
vanishes on every solution: the statements G-1 to G-12 are necessary
conditions, proved by ideal membership (reduction to zero by the basis).
To restrict to the solutions on which a polynomial does not vanish (for
example \(c^{\rm sl}_{\rm LL}\neq0\)), the ideal is enlarged by that
polynomial times the auxiliary variable \(w^{\rm sat}_{\rm LL}\) minus 1
(Rabinowitsch's trick). The command \texttt{malcev} of the Rust program
then verifies every solution found here on all basis triples and
4-tuples in exact arithmetic (\(14^4=38416\) 4-tuples for
\(\mathcal F^{\rm A}_{\rm LL}\), \(18^4=104976\) for
\(\mathcal F^{\rm B}_{\rm LL}\), \(17^4=83521\) for
\(\mathcal F^{\rm C}_{\rm LL}\)), so the solutions are proved to satisfy
the identities, and it exhibits explicit failing triples or 4-tuples for
the excluded members (data file \texttt{sl\_malcev\_families.csv}).

\textbf{Theorem M.} (a) In \(\mathcal F^{\rm A}_{\rm LL}\) the Jacobi
identity forces \(c^{\rm sl}_{\rm LL}=0\) and
\(c^{\rm tw}_{\rm LL}\in\{0,1\}\). (b) The Malcev identity forces
\(c^{\rm tw}_{\rm LL}(c^{\rm tw}_{\rm LL}-1)(c^{\rm tw}_{\rm LL}+2)=0\);
if \(c^{\rm sl}_{\rm LL}\neq0\) then \(c^{\rm tw}_{\rm LL}=-2\), and the
rescaling \(\mathbf b\mapsto\lambda^{\rm sc}_{\rm LL}\mathbf b\), which
maps \(c^{\rm sl}_{\rm LL}\) to
\((\lambda^{\rm sc}_{\rm LL})^2c^{\rm sl}_{\rm LL}\), makes
\(c^{\rm sl}_{\rm LL}=1\): the algebra is \(\mathfrak m_{\rm LL}\). (c)
In \(\mathcal F^{\rm B}_{\rm LL}\) even the Malcev identity forces
\(c^{\rm sl}_{\rm LL}=0\); the only non-Lie member has
\(c^{\rm tw}_{\rm LL}=-2\) and all other coefficients zero; the Lie
members are the Poincaré algebra with commuting spinor charges, the
Poincaré algebra with
\([\mathfrak s_{\rm LL},\mathfrak s_{\rm LL}]\subset\mathfrak p_{\rm LL}\),
and, for \(c^{\rm PP}_{\rm LL}=\lvert c^{\rm Ps}_{\rm LL}\rvert^2>0\),
the de Sitter algebra \(\mathfrak{so}(4,1)\) on its real spinor
representation. (d) In \(\mathcal F^{\rm C}_{\rm LL}\) the Jacobi
identity forces \(c^{\rm sl}_{\rm LL}=0\), and if
\(c^{\rm sl}_{\rm LL}\neq0\) the Malcev identity forces
\(c^{\rm Ts}_{\rm LL}=c^{\rm sT}_{\rm LL}=0\) and
\(c^{\rm tw}_{\rm LL}=-2\).

\begin{longtable}[]{@{}
  >{\raggedright\arraybackslash}p{(\columnwidth - 8\tabcolsep) * \real{0.2000}}
  >{\raggedright\arraybackslash}p{(\columnwidth - 8\tabcolsep) * \real{0.2000}}
  >{\raggedright\arraybackslash}p{(\columnwidth - 8\tabcolsep) * \real{0.2000}}
  >{\raggedright\arraybackslash}p{(\columnwidth - 8\tabcolsep) * \real{0.2000}}
  >{\raggedright\arraybackslash}p{(\columnwidth - 8\tabcolsep) * \real{0.2000}}@{}}
\toprule\noalign{}
\begin{minipage}[b]{\linewidth}\raggedright
family
\end{minipage} & \begin{minipage}[b]{\linewidth}\raggedright
identity
\end{minipage} & \begin{minipage}[b]{\linewidth}\raggedright
independent constraint polynomials
\end{minipage} & \begin{minipage}[b]{\linewidth}\raggedright
random tuples used
\end{minipage} & \begin{minipage}[b]{\linewidth}\raggedright
elements of the reduced Groebner basis
\end{minipage} \\
\midrule\noalign{}
\endhead
\bottomrule\noalign{}
\endlastfoot
\(\mathcal F^{\rm A}_{\rm LL}\) & Jacobi & 6 & 13 & 5 \\
\(\mathcal F^{\rm A}_{\rm LL}\) & Malcev & 8 & 13 & 9 \\
\(\mathcal F^{\rm B}_{\rm LL}\) & Jacobi & 19 & 15 & 19 \\
\(\mathcal F^{\rm B}_{\rm LL}\) & Malcev & 85 & 20 & 20 \\
\(\mathcal F^{\rm C}_{\rm LL}\) & Jacobi & 12 & 15 & 25 \\
\(\mathcal F^{\rm C}_{\rm LL}\) & Malcev & 43 & 16 & 113 \\
\end{longtable}

The output of the script:

\begin{verbatim}
Revision_SuperLoop, script `malcev_groebner.py`: Lie and Malcev classification of the covariant families by Groebner bases
12 of 12 checks passed

[PASS] G-1  F_A (l_LL + s_LL), Jacobi identity: csl_re and csl_im lie in the ideal, so every LIE algebra of the family has [s_LL, s_LL] = 0 (no Lorentz part): bosonic spinor charges cannot close into so(3,1) (Theorem L)
[PASS] G-2  F_A, Jacobi identity: Re and Im of c_tw (c_tw - 1) with c_tw = ctw_re + i ctw_im, i.e. ctw_re^2 - ctw_re - ctw_im^2 and ctw_im (2 ctw_re - 1), lie in the ideal; their only real common zeros are c_tw = 0 and c_tw = 1: the Lie action of so(3,1) on the spinor charges is 0 or the Weyl representation
[PASS] G-3  F_A, Malcev identity: Re and Im of c_tw (c_tw - 1)(c_tw + 2) with c_tw = ctw_re + i ctw_im lie in the ideal; this complex cubic has only the real roots 0, 1, -2: the twist is 0, 1 (Lie) or -2 (non-Lie), and ctw_im = 0
[PASS] G-4  F_A, Malcev identity: on every real solution with [s_LL, s_LL] != 0 (csl_re^2 + csl_im^2 != 0) the twist is c_tw = -2 (ctw_re + 2 and ctw_im lie in the saturated ideal): the bi-octonion algebra m_LL is the only solution with c_sl != 0 (up to the rescaling bS -> lambda_sc bS by a complex number lambda_sc != 0, which maps c_sl to lambda_sc^2 c_sl)
[PASS] G-5  F_B (Poincare or (anti-)de Sitter + spinor charges), Jacobi identity: csl_re, csl_im, cPP_im lie in the ideal; ctw_im = 0 and ctw_re (ctw_re - 1) = 0 on real solutions (ctw_im (4 ctw_im^2 + 1) lies in the ideal)
[PASS] G-6  F_B, Jacobi identity: cPP_re ctw_re - (cPs_re^2 + cPs_im^2), csP (ctw_re - 1), csP cPP_re, csP cPs_re, csP cPs_im lie in the ideal: the Lie solutions are (i) Poincare with abelian spinor charges, (ii) Poincare with [s_LL, s_LL] in p_LL (csP != 0, a bosonic analogue of supersymmetry), (iii) curvature cPP_re = |c_Ps|^2 >= 0 with [p_LL, s_LL] != 0 and [s_LL, s_LL] = 0
[PASS] G-7  F_B, Malcev identity: csl_re and csl_im lie in the ideal: in the presence of translations (flat or curved) even a MALCEV algebra cannot close spinor charges into the Lorentz algebra
[PASS] G-8  F_B, Malcev identity: Re and Im of c_tw (c_tw - 1)(c_tw + 2) lie in the ideal (c_tw = 0, 1 or -2 on real points), and cPs_re (ctw_re - 1), cPs_im (ctw_re - 1), csP (ctw_re - 1), cPP_re ctw_re - |c_Ps|^2 lie in it: for the twist c_tw = -2 one gets c_Ps = 0, csP = 0, cPP_re = 0, c_sl = 0
[PASS] G-9  F_B, Malcev identity with the twist c_tw = -2: the ideal contains ctw_im, cPP_re, cPP_im, cPs_re, cPs_im, csl_re, csl_im, csP: the only non-Lie Malcev extension of the Poincare algebra by spinor charges in this family is the ABELIAN twisted module (spinor charges commute with each other and with the translations)
[PASS] G-10  F_C (l_LL + su(2) isospin + spinor charges), Jacobi identity: csl_re^2 + csl_im^2 lies in the ideal, so c_sl = 0 on every real solution: every Lie algebra of the family has no Lorentz part in [s_LL, s_LL], even when isospin is present
[PASS] G-11  F_C, Malcev identity: csl_re |c_Ts|^2, csl_im |c_Ts|^2, csl_re |c_sT|^2, csl_im |c_sT|^2 lie in the ideal: whenever the spinor charges close into the Lorentz algebra (c_sl != 0), isospin acts trivially on them (c_Ts = 0) and [s_LL, s_LL] has no isospin part (c_sT = 0) -- isospin can join the super-loop only as an outer automorphism
[PASS] G-12  F_C, Malcev identity: on real solutions with c_sl != 0, G-11 gives |c_Ts|^2 = 0 and |c_sT|^2 = 0, hence c_Ts = c_sT = 0; adding cTs_re, cTs_im, csT_re, csT_im to the ideal and saturating by csl_re^2 + csl_im^2 != 0, the ideal contains ctw_re + 2 and ctw_im: the twist is -2 -- the spinor charges with the Lorentz bracket form the bi-octonion algebra m_LL, and su(2) isospin decouples (it can act on m_LL only by automorphisms)
\end{verbatim}

The reduced Gröbner bases of all six ideals (data file
\texttt{sl\_groebner\_bases.txt}):

\begin{verbatim}
Revision_SuperLoop, script code/scripts/malcev_groebner.py: reduced Groebner bases of the constraint ideals
(sympy 1.14.0); variables and brackets are defined in the docstring of the script

== family F_A, jacobi identity: 6 independent constraint polynomials from 13 random tuples; reduced Groebner basis (grevlex, variables ctw_re, ctw_im, csl_re, csl_im):
   ctw_im*(4*ctw_im**2 + 1)
   -ctw_im**2 + ctw_re**2 - ctw_re
   ctw_im*(2*ctw_re - 1)
   csl_re
   csl_im
== family F_A, malcev identity: 8 independent constraint polynomials from 13 random tuples; reduced Groebner basis (grevlex, variables ctw_re, ctw_im, csl_re, csl_im):
   ctw_im*(12*ctw_im**4 + 35*ctw_im**2 - 10*ctw_re + 13)
   ctw_im*(12*ctw_im**2*ctw_re + 4*ctw_im**2 + 7*ctw_re - 1)
   csl_im*(2*ctw_im**2*ctw_re - ctw_im**2 - 3*ctw_re - 6)
   csl_im*ctw_im*(2*ctw_im**2 + 5*ctw_re + 7)
   -3*ctw_im**2*ctw_re - ctw_im**2 + ctw_re**3 + ctw_re**2 - 2*ctw_re
   -ctw_im*(ctw_im**2 - 3*ctw_re**2 - 2*ctw_re + 2)
   csl_im*(ctw_im**2 + ctw_re**2 + 4*ctw_re + 4)
   -csl_im*ctw_im + csl_re*ctw_re + 2*csl_re
   csl_im*ctw_re + 2*csl_im + csl_re*ctw_im
== family F_B, jacobi identity: 19 independent constraint polynomials from 15 random tuples; reduced Groebner basis (grevlex, variables ctw_re, ctw_im, cPP_re, cPP_im, cPs_re, cPs_im, csl_re, csl_im, csP):
   ctw_im*(4*ctw_im**2 + 1)
   -cPs_re*(cPP_re - cPs_im**2 - cPs_re**2)
   -cPs_im*(cPP_re - cPs_im**2 - cPs_re**2)
   -ctw_im**2 + ctw_re**2 - ctw_re
   ctw_im*(2*ctw_re - 1)
   cPP_re*ctw_re - cPs_im**2 - cPs_re**2
   cPP_re*ctw_im
   cPs_re*(ctw_re - 1)
   cPs_re*ctw_im
   cPs_im*(ctw_re - 1)
   cPs_im*ctw_im
   csP*(ctw_re - 1)
   csP*ctw_im
   cPP_re*csP
   cPs_re*csP
   cPs_im*csP
   cPP_im
   csl_re
   csl_im
== family F_B, malcev identity: 85 independent constraint polynomials from 20 random tuples; reduced Groebner basis (grevlex, variables ctw_re, ctw_im, cPP_re, cPP_im, cPs_re, cPs_im, csl_re, csl_im, csP):
   ctw_im*(12*ctw_im**4 + 35*ctw_im**2 - 10*ctw_re + 13)
   ctw_im*(12*ctw_im**2*ctw_re + 4*ctw_im**2 + 7*ctw_re - 1)
   -3*ctw_im**2*ctw_re - ctw_im**2 + ctw_re**3 + ctw_re**2 - 2*ctw_re
   -ctw_im*(ctw_im**2 - 3*ctw_re**2 - 2*ctw_re + 2)
   -cPs_re*(cPP_re - cPs_im**2 - cPs_re**2)
   -cPs_im*(cPP_re - cPs_im**2 - cPs_re**2)
   cPP_re*ctw_re - cPs_im**2 - cPs_re**2
   cPP_re*ctw_im
   cPs_re*(ctw_re - 1)
   cPs_re*ctw_im
   cPs_im*(ctw_re - 1)
   cPs_im*ctw_im
   csP*(ctw_re - 1)
   csP*ctw_im
   cPP_re*csP
   cPs_re*csP
   cPs_im*csP
   cPP_im
   csl_re
   csl_im
== family F_C, jacobi identity: 12 independent constraint polynomials from 15 random tuples; reduced Groebner basis (grevlex, variables cTT, ctw_re, ctw_im, cTs_re, cTs_im, csl_re, csl_im, csT_re, csT_im):
   ctw_im*(4*ctw_im**2 + 1)
   cTs_im*(cTs_im**2 + cTs_re**2)
   csl_im*(2*ctw_im**2 - ctw_re + 1)
   -csl_im*(cTs_im*ctw_im - cTs_re*ctw_re + cTs_re)
   csl_im*(cTs_im*ctw_re - cTs_im + cTs_re*ctw_im)
   csl_im*(cTs_im**2 + cTs_re**2)
   -ctw_im**2 + ctw_re**2 - ctw_re
   ctw_im*(2*ctw_re - 1)
   cTT*cTs_re + cTs_im**2 - cTs_re**2
   cTs_im*(cTT - 2*cTs_re)
   cTT*csl_re + 2*cTs_im*csl_im
   -csl_im*ctw_im + csl_re*ctw_re - csl_re
   csl_im*ctw_re - csl_im + csl_re*ctw_im
   cTs_im*csl_im + cTs_re*csl_re
   cTs_im*csl_re - cTs_re*csl_im
   csl_im**2 + csl_re**2
   csl_im*(cTT - 2*cTs_re)
   cTT*csT_re - 2*csl_re
   cTs_re*csT_re - csl_re
   cTs_im*csT_re - csl_im
   csT_im*csl_im + csT_re*csl_re
   -csT_im*csl_re + csT_re*csl_im
   cTT*csT_im - 2*csl_im
   cTs_re*csT_im - csl_im
   cTs_im*csT_im + csl_re
== family F_C, malcev identity: 43 independent constraint polynomials from 16 random tuples; reduced Groebner basis (grevlex, variables cTT, ctw_re, ctw_im, cTs_re, cTs_im, csl_re, csl_im, csT_re, csT_im):
   ctw_im*(12*ctw_im**4 + 35*ctw_im**2 - 10*ctw_re + 13)
   cTs_im*(6*cTT*ctw_im**2*ctw_re + 3*cTT*ctw_im**2 + 3*cTT*ctw_re - 4*cTs_im*ctw_im**3 - 10*cTs_im*ctw_im*ctw_re + 4*cTs_im*ctw_im - 12*cTs_re*ctw_im**2 - 6*cTs_re*ctw_re)
   cTs_im*(cTT*ctw_im + 6*cTs_im*ctw_re)*(2*ctw_im**2 - ctw_re + 1)
   cTs_im*(2*ctw_im**2 - 9*ctw_re)*(2*ctw_im**2 - ctw_re + 1)
   cTs_im*(cTs_im**2 + cTs_re**2)*(2*ctw_im**2 + ctw_re)
   -cTs_im*(cTs_im**2 + cTs_re**2)*(cTs_im*ctw_im - 3*cTs_re*ctw_re)
   cTs_im*(cTs_im**2 + cTs_re**2)*(cTs_im*ctw_re + 3*cTs_re*ctw_im)
   cTs_im*(cTs_im**2 + cTs_re**2)*(cTs_im**2 + 9*cTs_re**2)
   csl_im*(2*cTT*ctw_im**2*ctw_re - cTT*ctw_im**2 - 3*cTT*ctw_re - 6*cTT - 6*cTs_im*ctw_im*ctw_re + 6*cTs_im*ctw_im + 3*cTs_re*ctw_re + 15*cTs_re)
   csl_im*(2*cTT*ctw_im**3 + 5*cTT*ctw_im*ctw_re + 7*cTT*ctw_im - 18*cTs_im*ctw_im**2 - 9*cTs_im*ctw_re + 9*cTs_im - 36*cTs_re*ctw_im)
   csl_im*(2*ctw_im**2 - ctw_re + 1)*(2*ctw_im**2 + 3*ctw_re + 6)
   csl_im*(4*cTs_im*ctw_im**2*ctw_re + 4*cTs_im*ctw_im**2 + 3*cTs_im*ctw_re - 3*cTs_im + 6*cTs_re*ctw_im)
   csl_im*(4*cTs_im*ctw_im**3 - 8*cTs_im*ctw_im*ctw_re + 14*cTs_im*ctw_im - 9*cTs_re*ctw_re + 9*cTs_re)
   -csl_im*(6*cTs_im*csT_im*ctw_im*ctw_re - 9*cTs_im*csT_im*ctw_im - 15*cTs_re*csT_im - 4*csl_im*ctw_im**2*ctw_re + 2*csl_im*ctw_im**2 + 5*csl_im*ctw_re + 10*csl_im - csl_re*ctw_im)
   -csl_im*(18*cTs_im*csT_im*ctw_im**2 - 27*cTs_im*csT_im*ctw_re - 9*cTs_im*csT_im - 4*csl_im*ctw_im**3 - 10*csl_im*ctw_im*ctw_re - 2*csl_im*ctw_im - 12*csl_re*ctw_re - 24*csl_re)
   6*cTT*csT_im*ctw_im**2*ctw_re + 3*cTT*csT_im*ctw_im**2 + 3*cTT*csT_im*ctw_re - 4*csl_im*ctw_im**2*ctw_re - csl_im*ctw_im**2 - 3*csl_im*ctw_re**2 - 5*csl_im*ctw_re + 2*csl_im + csl_re*ctw_im
   ctw_im*(6*cTT*csT_im*ctw_im**2 - 3*cTT*csT_im*ctw_re + 3*cTT*csT_im - 4*csl_im*ctw_im**2 - 4*csl_im*ctw_re - 2*csl_im + 6*csl_re*ctw_im)
   csT_im*(2*ctw_im**2 - 9*ctw_re)*(2*ctw_im**2 - ctw_re + 1)
   18*cTs_im*csT_im*ctw_im**2*ctw_re + 9*cTs_im*csT_im*ctw_im**2 + 9*cTs_im*csT_im*ctw_re + 2*csl_im*ctw_im**3 + 8*csl_im*ctw_im*ctw_re + 10*csl_im*ctw_im + 3*csl_re*ctw_im**2 + 3*csl_re*ctw_re + 6*csl_re
   4*cTs_im*csT_im*ctw_im**3 - 2*cTs_im*csT_im*ctw_im*ctw_re + 2*cTs_im*csT_im*ctw_im - 4*csl_im*ctw_im**2*ctw_re + csl_im*ctw_im**2 - csl_im*ctw_re**2 - csl_im*ctw_re + 2*csl_im - 3*csl_re*ctw_im
   csl_im*(4*csT_im*ctw_im**2*ctw_re + 4*csT_im*ctw_im**2 + 3*csT_im*ctw_re - 3*csT_im + 6*csT_re*ctw_im)
   csl_im*(4*csT_im*ctw_im**3 - 8*csT_im*ctw_im*ctw_re + 14*csT_im*ctw_im - 9*csT_re*ctw_re + 9*csT_re)
   csl_im*(4*cTs_im**2*csT_im*ctw_re + 4*cTs_im**2*csT_im + 2*cTs_im*csl_im*ctw_im + cTs_im*csl_re + 2*cTs_re*csl_im*ctw_re + 5*cTs_re*csl_im)
   csl_im*(4*cTs_im**2*csT_im*ctw_im + 4*cTs_im*cTs_re*csT_im - 2*cTs_im*csl_im*ctw_re - 3*cTs_im*csl_im + 2*cTs_re*csl_im*ctw_im - cTs_re*csl_re)
   csl_im*(4*cTs_im*csT_im**2*ctw_re + 4*cTs_im*csT_im**2 + 2*csT_im*csl_im*ctw_im + csT_im*csl_re + 2*csT_re*csl_im*ctw_re + 5*csT_re*csl_im)
   csl_im*(4*cTs_im*csT_im**2*ctw_im + 4*cTs_re*csT_im**2 - 2*csT_im*csl_im*ctw_re - 3*csT_im*csl_im + 2*csT_re*csl_im*ctw_im - csT_re*csl_re)
   ctw_im*(12*ctw_im**2*ctw_re + 4*ctw_im**2 + 7*ctw_re - 1)
   12*cTs_im*ctw_im**2*ctw_re - cTs_im*ctw_im**2 - 5*cTs_im*ctw_re**2 + 5*cTs_im*ctw_re + 4*cTs_re*ctw_im**3 + cTs_re*ctw_im
   cTs_im*(cTT*ctw_im**2 + cTT*ctw_re**2 - 4*cTs_im*ctw_im*ctw_re + 2*cTs_im*ctw_im - 4*cTs_re*ctw_im**2 - 2*cTs_re*ctw_re)
   cTs_im*(cTT + 2*cTs_re)*(cTs_im**2 + cTs_re**2)
   12*csl_im*ctw_im**2*ctw_re + 5*csl_im*ctw_im**2 + csl_im*ctw_re**2 + csl_im*ctw_re - 2*csl_im + 4*csl_re*ctw_im**3 + 9*csl_re*ctw_im
   csl_im*(cTT*ctw_im**2 + cTT*ctw_re**2 + 4*cTT*ctw_re + 4*cTT - 6*cTs_im*ctw_im - 6*cTs_re*ctw_re - 12*cTs_re)
   csl_im*(cTs_im*ctw_im*ctw_re - 2*cTs_im*ctw_im + cTs_re*ctw_im**2 + cTs_re*ctw_re - cTs_re)
   cTs_im*csl_im*(cTT - 2*cTs_re)*(ctw_re + 1)
   csl_im*(cTs_im*ctw_im**2 + cTs_im*ctw_re**2 + cTs_im*ctw_re - 2*cTs_im + 3*cTs_re*ctw_im)
   csl_im*(cTT*cTs_im*ctw_im + cTT*cTs_re + 2*cTs_im**2 - 2*cTs_im*cTs_re*ctw_im)
   csl_im*(3*cTs_im*csT_im*ctw_re + 3*cTs_im*csT_im + csl_im*ctw_im*ctw_re + csl_re*ctw_im**2 + 2*csl_re*ctw_re + 4*csl_re)
   -csl_im*(6*cTs_im*csT_im*ctw_im + 6*cTs_re*csT_im - csl_im*ctw_im**2 - csl_im*ctw_re**2 - 3*csl_im*ctw_re - 2*csl_im + csl_re*ctw_im)
   12*csT_im*ctw_im**2*ctw_re - csT_im*ctw_im**2 - 5*csT_im*ctw_re**2 + 5*csT_im*ctw_re + 4*csT_re*ctw_im**3 + csT_re*ctw_im
   csl_im*(csT_im*ctw_im*ctw_re - 2*csT_im*ctw_im + csT_re*ctw_im**2 + csT_re*ctw_re - csT_re)
   3*cTT*csT_im*ctw_im**2 + 3*cTT*csT_im*ctw_re**2 - 2*csl_im*ctw_im**2 - 2*csl_im*ctw_re**2 - 4*csl_im*ctw_re + 4*csl_re*ctw_im
   3*cTs_im*csT_im*ctw_im**2 + 3*cTs_im*csT_im*ctw_re**2 + 2*csl_im*ctw_im*ctw_re + 3*csl_im*ctw_im + 2*csl_re*ctw_im**2 + csl_re*ctw_re + 2*csl_re
   -2*cTs_im**2*csT_im*csl_im + 2*cTs_im*cTs_re*csT_im*csl_re - 2*cTs_im*csl_im*csl_re - cTs_re*csl_im**2 + cTs_re*csl_re**2
   2*cTs_im**2*csT_im*csl_re + 2*cTs_im*cTs_re*csT_im*csl_im - cTs_im*csl_im**2 + cTs_im*csl_re**2 + 2*cTs_re*csl_im*csl_re
   csT_im*(cTT - 2*cTs_re)*(csl_im**2 + csl_re**2)
   csl_im*(3*cTT*csT_im*ctw_re + 3*cTT*csT_im - 6*cTs_im*csT_im*ctw_im - 6*cTs_re*csT_im - 2*csl_im*ctw_re - 4*csl_im - 2*csl_re*ctw_im)
   csl_im*(csT_im*ctw_im**2 + csT_im*ctw_re**2 + csT_im*ctw_re - 2*csT_im + 3*csT_re*ctw_im)
   3*cTT*csT_im*csl_im*ctw_im - 3*cTT*csT_im*csl_re + 6*cTs_im*csT_im*csl_im*ctw_re + 6*cTs_re*csT_im*csl_re - 2*csl_im**2*ctw_im + 2*csl_im*csl_re*ctw_re + 4*csl_im*csl_re
   4*cTT*csT_im**2*ctw_re + 4*cTs_im*csT_im**2*ctw_im - 6*csT_im*csl_im*ctw_re - 7*csT_im*csl_im + 6*csT_re*csl_im*ctw_im - 5*csT_re*csl_re
   4*cTT*csT_im**2*ctw_im - 4*cTs_im*csT_im**2*ctw_re - 6*csT_im*csl_im*ctw_im + 7*csT_im*csl_re - 6*csT_re*csl_im*ctw_re - 5*csT_re*csl_im
   2*cTT*csT_im**2*csl_re + 2*cTs_im*csT_im**2*csl_im - csT_im*csl_im*csl_re + csT_re*csl_im**2 + 2*csT_re*csl_re**2
   -2*cTs_im*csT_im**2*csl_im + 2*cTs_re*csT_im**2*csl_re - 2*csT_im*csl_im*csl_re - csT_re*csl_im**2 + csT_re*csl_re**2
   2*cTs_im*csT_im**2*csl_re + 2*cTs_re*csT_im**2*csl_im - csT_im*csl_im**2 + csT_im*csl_re**2 + 2*csT_re*csl_im*csl_re
   csl_im*(2*cTT*csT_im**2 + 2*cTs_re*csT_im**2 - 3*csT_im*csl_im + 3*csT_re*csl_re)
   -3*ctw_im**2*ctw_re - ctw_im**2 + ctw_re**3 + ctw_re**2 - 2*ctw_re
   -ctw_im*(ctw_im**2 - 3*ctw_re**2 - 2*ctw_re + 2)
   2*cTT**2*cTs_re + cTT*cTs_im**2 - cTT*cTs_re**2 + 3*cTs_im**2*cTs_re - cTs_re**3
   -cTT*cTs_im*ctw_im + cTT*cTs_re*ctw_re + cTs_im**2*ctw_re + 2*cTs_im*cTs_re*ctw_im - cTs_re**2*ctw_re
   -2*cTs_im*ctw_im*ctw_re + cTs_im*ctw_im - cTs_re*ctw_im**2 + cTs_re*ctw_re**2 - cTs_re*ctw_re
   cTT*cTs_im*ctw_re + cTT*cTs_re*ctw_im + cTs_im**2*ctw_im - 2*cTs_im*cTs_re*ctw_re - cTs_re**2*ctw_im
   -cTs_im*ctw_im**2 + cTs_im*ctw_re**2 - cTs_im*ctw_re + 2*cTs_re*ctw_im*ctw_re - cTs_re*ctw_im
   cTs_im*(2*cTT**2 - 2*cTT*cTs_re + cTs_im**2 - 3*cTs_re**2)
   -cTT*csl_im*ctw_im + cTT*csl_re*ctw_re + 2*cTT*csl_re + 3*cTs_im*csl_im - 3*cTs_re*csl_re
   -2*csl_im*ctw_im*ctw_re - csl_im*ctw_im - csl_re*ctw_im**2 + csl_re*ctw_re**2 + csl_re*ctw_re - 2*csl_re
   cTT*csl_im*ctw_re + 2*cTT*csl_im + cTT*csl_re*ctw_im - 3*cTs_im*csl_re - 3*cTs_re*csl_im
   -csl_im*ctw_im**2 + csl_im*ctw_re**2 + csl_im*ctw_re - 2*csl_im + 2*csl_re*ctw_im*ctw_re + csl_re*ctw_im
   -cTT*cTs_im*csl_im + cTT*cTs_re*csl_re + 2*cTs_im**2*csl_re + 2*cTs_im*cTs_re*csl_im
   3*cTs_im*csl_im - 2*cTs_re*csl_im*ctw_im + 2*cTs_re*csl_re*ctw_re + cTs_re*csl_re
   -3*cTs_im*csl_re + 2*cTs_re*csl_im*ctw_re + cTs_re*csl_im + 2*cTs_re*csl_re*ctw_im
   csl_re*(cTs_im**2 + cTs_re**2)
   cTT*cTs_im*csl_re + cTT*cTs_re*csl_im + 2*cTs_im**2*csl_im - 2*cTs_im*cTs_re*csl_re
   -2*cTs_im*csl_im*ctw_im + 2*cTs_im*csl_re*ctw_re + cTs_im*csl_re - 3*cTs_re*csl_im
   2*cTs_im*csl_im*ctw_re + cTs_im*csl_im + 2*cTs_im*csl_re*ctw_im + 3*cTs_re*csl_re
   3*cTs_im*csT_im*csl_re + 3*cTs_re*csT_im*csl_im - csl_im**2*ctw_re - 2*csl_im**2 - 2*csl_im*csl_re*ctw_im + csl_re**2*ctw_re + 2*csl_re**2
   3*cTs_im*csT_im*csl_im - 3*cTs_re*csT_im*csl_re - csl_im**2*ctw_im + 2*csl_im*csl_re*ctw_re + 4*csl_im*csl_re + csl_re**2*ctw_im
   csl_im*(cTs_im**2 + cTs_re**2)
   2*cTT**2*csT_re + 2*cTs_im*cTs_re*csT_im + 3*cTs_im*csl_im - 3*cTs_re*csl_re
   -3*cTT*csT_im*ctw_im + 3*cTT*csT_re*ctw_re + 2*csl_im*ctw_im - 2*csl_re*ctw_re - 4*csl_re
   -2*csT_im*ctw_im*ctw_re + csT_im*ctw_im - csT_re*ctw_im**2 + csT_re*ctw_re**2 - csT_re*ctw_re
   3*cTT*csT_im*ctw_re + 3*cTT*csT_re*ctw_im - 2*csl_im*ctw_re - 4*csl_im - 2*csl_re*ctw_im
   -csT_im*ctw_im**2 + csT_im*ctw_re**2 - csT_im*ctw_re + 2*csT_re*ctw_im*ctw_re - csT_re*ctw_im
   2*cTT*cTs_re*csT_re - 2*cTs_im*cTs_re*csT_im + 3*cTs_im*csl_im - 3*cTs_re*csl_re
   ctw_re*(cTs_im*csT_im + cTs_re*csT_re)
   ctw_im*(cTs_im*csT_im + cTs_re*csT_re)
   cTs_re*(cTs_im*csT_im + cTs_re*csT_re)
   2*cTT*cTs_im*csT_re - 2*cTs_im**2*csT_im - 3*cTs_im*csl_re - 3*cTs_re*csl_im
   -3*cTs_im*csT_im*ctw_im + 3*cTs_im*csT_re*ctw_re - csl_im*ctw_re - 2*csl_im - csl_re*ctw_im
   3*cTs_im*csT_im*ctw_re + 3*cTs_im*csT_re*ctw_im - csl_im*ctw_im + csl_re*ctw_re + 2*csl_re
   cTs_im*(cTs_im*csT_im + cTs_re*csT_re)
   cTs_im*(cTs_im*csT_re - cTs_re*csT_im)
   -cTT*csT_im*csl_im + cTT*csT_re*csl_re + 2*cTs_im*csT_im*csl_re + 2*cTs_re*csT_im*csl_im
   3*csT_im*csl_im - 2*csT_re*csl_im*ctw_im + 2*csT_re*csl_re*ctw_re + csT_re*csl_re
   -3*csT_im*csl_re + 2*csT_re*csl_im*ctw_re + csT_re*csl_im + 2*csT_re*csl_re*ctw_im
   csl_re*(cTs_im*csT_im + cTs_re*csT_re)
   csl_re*(cTs_im*csT_re - cTs_re*csT_im)
   cTT*csT_im*csl_re + cTT*csT_re*csl_im + 2*cTs_im*csT_im*csl_im - 2*cTs_re*csT_im*csl_re
   csl_im*(cTs_im*csT_im + cTs_re*csT_re)
   csl_im*(cTs_im*csT_re - cTs_re*csT_im)
   cTT*(csT_im**2 + csT_re**2)
   cTs_re*(csT_im**2 + csT_re**2)
   cTs_im*(csT_im**2 + csT_re**2)
   csl_re*(csT_im**2 + csT_re**2)
   csl_im*(csT_im**2 + csT_re**2)
   2*cTT**2*csT_im + 2*cTs_im**2*csT_im - 3*cTs_im*csl_re - 3*cTs_re*csl_im
   2*cTT*cTs_re*csT_im - 2*cTs_im**2*csT_im - 3*cTs_im*csl_re - 3*cTs_re*csl_im
   -3*cTs_im*csT_im*ctw_im + 3*cTs_re*csT_im*ctw_re - csl_im*ctw_re - 2*csl_im - csl_re*ctw_im
   3*cTs_im*csT_im*ctw_re + 3*cTs_re*csT_im*ctw_im - csl_im*ctw_im + csl_re*ctw_re + 2*csl_re
   csT_im*(cTs_im**2 + cTs_re**2)
   2*cTT*cTs_im*csT_im + 2*cTs_im*cTs_re*csT_im - 3*cTs_im*csl_im + 3*cTs_re*csl_re
   -2*csT_im*csl_im*ctw_im + 2*csT_im*csl_re*ctw_re + csT_im*csl_re - 3*csT_re*csl_im
   2*csT_im*csl_im*ctw_re + csT_im*csl_im + 2*csT_im*csl_re*ctw_im + 3*csT_re*csl_re
   2*cTT*csT_im*csT_re - cTs_im*csT_im**2 + cTs_re*csT_im*csT_re - 3*csT_im*csl_re - 3*csT_re*csl_im
   2*cTT*csT_im**2 + cTs_im*csT_im*csT_re + cTs_re*csT_im**2 - 3*csT_im*csl_im + 3*csT_re*csl_re
\end{verbatim}

The exhaustive verification by the command \texttt{malcev}:

\begin{verbatim}
Revision_SuperLoop, command `malcev`: exhaustive exact verification of the Lie and Malcev classification of Lorentz-covariant extensions by spinor charges
9 of 9 checks passed

[PASS] M-1  F_A with the twist c_tw = -2 and c_sl = 1 is exactly the commutator algebra m_LL of the imaginary bi-octonions (all 196 basis brackets agree); it satisfies the Malcev identity on all 38416 basis 4-tuples and violates the Jacobi identity
[PASS] M-2  F_A with c_sl = 0: the twists c_tw = 0, 1, -2 give Malcev algebras (exhaustive); c_tw = 0 and c_tw = 1 give Lie algebras, c_tw = -2 does not (Carlsson's non-Lie Malcev module of sl(2))
[PASS] M-3  excluded members of F_A fail exhaustively: (c_tw, c_sl) = (1, 1) [Lie action and spinor charges closing into Lorentz], (0, 1), (2, 0), (-1, 0), (3, 0) violate both identities (in agreement with G-1, G-3, G-4)  -- first failing Malcev 4-tuple (positions in the basis) (0, 6, 6, 8)
[PASS] M-4  F_B with the twist c_tw = -2 and all other coefficients 0 (the abelian twisted module of the Poincare algebra) satisfies the Malcev identity on all 104976 basis 4-tuples and violates the Jacobi identity: internal spinor charges, commuting with the translations p_LL, with the non-Lie Lorentz action  -- first failing Jacobi triple (positions in the basis) (0, 1, 10)
[PASS] M-5  the Lie solutions of F_B found by G-6 are Lie algebras (exhaustive): Poincare with abelian spinor charges (c_tw = 1), Poincare with [s_LL, s_LL] in p_LL (c_tw = 1, c_sP = 1: bosonic spinor charges whose brackets are translations), and curved space-time with c_PP = |c_Ps|^2 (c_tw = 1, c_PP = 1, c_Ps = 1 and c_PP = 2, c_Ps = 1 + i): spinor charges that do not commute with the translations
[PASS] M-6  for c_PP > 0 the space-time part l_LL + p_LL (with [p_LL, p_LL] = c_PP x Lorentz generator) has a Killing form of signature (4 positive, 6 negative): it is so(4,1), the de Sitter algebra; the spinor charges with [p_LL, s_LL] != 0 form its real 8-dimensional spinor representation  -- signature (4, 6, 0)
[PASS] M-7  excluded members of F_B fail exhaustively: the bi-octonion bracket together with translations (c_tw = -2, c_sl = 1), the Lie action with c_sl = 1 (with or without c_sP = 1), and the twist together with c_sP = 1 or with the curvature c_PP = 1: with translations present, spinor charges never close into the Lorentz algebra (G-7, G-9)  -- first failing Malcev 4-tuple (positions in the basis) (0, 6, 10, 14)
[PASS] M-8  F_C: the bi-octonion algebra m_LL with a decoupled isospin algebra su(2) (c_tw = -2, c_sl = 1, c_TT = 1, c_Ts = 0, c_sT = 0) is Malcev and not Lie; isospin acting on the spinor charges with the Lie Lorentz action and [s_LL, s_LL] = 0 (c_tw = 1, c_Ts = 1, c_TT = 1) is a Lie algebra
[PASS] M-9  isospin cannot act on the spinor charges of m_LL inside a Malcev algebra: adding c_Ts = 1 (isospin acts on s_LL), or c_sT = 1 ([s_LL, s_LL] gets an isospin part), or both, to m_LL + su(2) violates the Malcev identity (exhaustive; G-11): electric charge enters the super-loop as an outer operator (an automorphism), not as an element  -- first failing Malcev 4-tuple (positions in the basis) (0, 1, 6, 9)
\end{verbatim}

\hypertarget{command-manifolds}{%
\section{\texorpdfstring{6. Command
\texttt{manifolds}}{6. Command manifolds}}\label{command-manifolds}}

\begin{itemize}
\tightlist
\item
  \textbf{\(\mathbb R^{1,7}_{\rm LL}\) and \(\mathbb R^{1,3}_{\rm LL}\)}
  (N-1): the Hermitian bi-octonions
  \(\mathbf x=x^0_{\rm LL}\mathbf e_0+i\sum_{p=1}^7x^p_{\rm LL}\mathbf e_p\)
  with real coordinates have the real norm
  \(N_{\rm LL}(\mathbf x)=(x^0_{\rm LL})^2-\sum_{p=1}^7(x^p_{\rm LL})^2\),
  signature (1, 7); the Hermitian bi-quaternions (\(p\le3\)) have the
  signature (1, 3), and \(SL(2,\mathbb C)\) acts on them by
  \(\mathbf x\mapsto\mathbf g\mathbf x\mathbf g^\dagger\). The
  signatures are computed exactly from the Gram matrix of the norm.
\item
  \textbf{The conjugations} (N-4, N-5):
  \(\mathbf g=\exp(\theta_{\rm LL}\mathbf e_p)\) has
  \(\mathbf g^\dagger=\exp(-\theta_{\rm LL}\mathbf e_p)\) and
  \(\mathbf g=\exp(\theta_{\rm LL}i\mathbf e_p)\) has
  \(\mathbf g^\dagger=\mathbf g\), so the generators of
  \(\mathbf x\mapsto\mathbf g\mathbf x\mathbf g^\dagger\) are
  \(\mathbf x\mapsto[\mathbf e_p,\mathbf x]\) and
  \(\mathbf x\mapsto i(\mathbf e_p\mathbf x+\mathbf x\mathbf e_p)\). The
  program writes them as real \(8\times8\) integer matrices on the
  coordinates \(x^0_{\rm LL},\dots,x^7_{\rm LL}\) and computes their Lie
  closure exactly: 28 for \(p=1,\dots,7\) (all of
  \(\mathfrak{so}(1,7)\)), 9 for \(p=1,2,3\) (on
  \(\mathbb R^{1,3}_{\rm LL}\) alone: 6).
\item
  \textbf{The split octonions} (N-2): the map
  \((\mathbf a,\mathbf b)\mapsto\mathbf a+(i\mathbf b)\,\mathbf e_4\)
  from pairs of real quaternions with the product
  \((\mathbf a,\mathbf b)(\mathbf c,\mathbf d)=(\mathbf a\mathbf c+\bar{\mathbf d}\mathbf b,\ \mathbf d\mathbf a+\mathbf b\bar{\mathbf c})\)
  is multiplicative, because
  \(\overline{i\mathbf d}\,(i\mathbf b)=-\bar{\mathbf d}\mathbf b\); on
  its image
  \(\mathbb O^{\rm s}_{\rm LL}=\mathbb H+(i\mathbb H)\,\mathbf e_4\) the
  norm is
  \(N_{\rm LL}(\mathbf a+(i\mathbf b)\mathbf e_4)=N_{\rm LL}(\mathbf a)-N_{\rm LL}(\mathbf b)\),
  of signature (4, 4).
\item
  \textbf{The space of spinor charges} (N-3): for \(\mathbf b\) with
  real part \(\mathbf b_1\) and imaginary part \(\mathbf b_2\) (real
  quaternions),
  \(N_{\rm LL}(\mathbf b\,\mathbf e_4)=N_{\rm LL}(\mathbf b_1)-N_{\rm LL}(\mathbf b_2)+2i\,N_{\rm LL}(\mathbf b_1,\mathbf b_2)\):
  \(\mathrm{Re}\,N_{\rm LL}\) and \(\mathrm{Im}\,N_{\rm LL}\) both have
  the signature (4, 4).
\item
  \textbf{The frame fields} (N-7):
  \(X^{\rm LL}_{\mathbf x}(\mathbf w)=\mathbf w\mathbf x\) is a linear
  vector field, and the Lie bracket of two of them is
  \([X^{\rm LL}_{\mathbf x},X^{\rm LL}_{\mathbf y}](\mathbf w)=(\mathbf w\mathbf x)\mathbf y-(\mathbf w\mathbf y)\mathbf x=\mathbf w[\mathbf x,\mathbf y]+2[\mathbf w,\mathbf x,\mathbf y]\),
  where the last step writes
  \((\mathbf w\mathbf x)\mathbf y=\mathbf w(\mathbf x\mathbf y)+[\mathbf w,\mathbf x,\mathbf y]\)
  and uses
  \([\mathbf w,\mathbf y,\mathbf x]=-[\mathbf w,\mathbf x,\mathbf y]\);
  tangency is \(N_{\rm LL}(\mathbf w,\mathbf w\mathbf x)=0\) for
  imaginary \(\mathbf x\). Both are checked on 100 random \(\mathbf w\),
  \(\mathbf x\), \(\mathbf y\) with Gaussian-integer coefficients.
\item
  \textbf{Closures} (N-6, N-8, N-9): the left multiplications by the
  imaginary bi-octonions generate \(\mathfrak{so}(8,\mathbb C)\) (real
  dimension 56), by the imaginary octonions \(\mathfrak{so}(8)\) (28),
  by the imaginary split units on \(\mathbb O^{\rm s}_{\rm LL}\)
  \(\mathfrak{so}(4,4)\) (28); the products
  \(L^{\rm LL}_{\mathbf e_p}L^{\rm LL}_{\mathbf e_{p'}}\) (\(p<p'\))
  span \(\mathfrak{spin}(7)\) (21), whose stabilizer of \(\mathbf e_0\)
  is \(\mathfrak g_2\) (14); \(\mathbf e_1\), \(i\mathbf e_4\),
  \(i\mathbf e_5\) have the brackets
  \([\mathbf e_1,i\mathbf e_4]=2i\mathbf e_5\),
  \([\mathbf e_1,i\mathbf e_5]=-2i\mathbf e_4\),
  \([i\mathbf e_4,i\mathbf e_5]=-2\mathbf e_1\), a Lie algebra whose
  Killing form has the signature (2, 1):
  \(\mathfrak{sl}(2,\mathbb R)=\mathfrak{so}(2,1)\).
\end{itemize}

\begin{verbatim}
Revision_SuperLoop, command `manifolds`: the manifolds that carry the Lorentz super-loop
9 of 9 checks passed

[PASS] N-1  the Hermitian bi-octonions x = x_LL^0 + i sum_p x_LL^p e_p (x^dag = x, real x_LL^0, x_LL^p, p = 1..7) form the 7+1 Minkowski space R^(1,7)_LL: the complex norm N_LL is real on them, with signature (1 positive, 7 negative); the Hermitian bi-quaternions (p = 1,2,3) form R^(1,3)_LL, signature (1,3), on which SL(2,C) acts by x -> gL x gL^dag  -- R^(1,7)_LL: (1, 7, 0), R^(1,3)_LL: (1, 3, 0)
[PASS] N-2  the split octonions sit inside B_LL as Os = H + (i H) e4: this real 8-dimensional subspace is closed under the product, the map from the split Cayley-Dickson algebra (pairs (aH, bS) of real quaternions with the sign +1) to B_LL, (aH, bS) -> aH + (i bS) e4, is an isomorphism onto Os, and N_LL restricted to Os has signature (4,4): the 4+4 manifold is the real form Os of the super-loop's algebra; its unit elements form the pseudo-sphere S^(3,4), a Moufang loop  -- signature (4, 4, 0)
[PASS] N-3  the space of spinor charges s_LL (real dimension 8) carries the quadratic forms Re N_LL and Im N_LL (invariant under the Lorentz and isospin automorphisms, because N_LL is multiplicative), both of signature (4,4): the base of the fibration of L_LL over s_LL (aH + bS e4 -> bS, fibre a copy of SL(2,C) wherever N_LL(aH) = 1 - N_LL(bS) is not zero) is a 4+4 space  -- Re N_LL (4, 4, 0), Im N_LL (4, 4, 0)
[PASS] N-4  on the 7+1 manifold the one-parameter conjugations x -> gL x gL^dag by the loop elements gL = exp(theta e_p) (rotations) and gL = exp(theta i e_p) (boosts), p = 1..7, have the generators x -> [e_p, x] and x -> i (e_p x + x e_p), which lie in so(1,7), and their Lie closure is all of so(1,7) (dimension 28, exact): the super-loop acts on R^(1,7)_LL through the Lorentz group SO(1,7)  -- closure dimension 28
[PASS] N-5  conjugation by the loop's own Lorentz subgroup (p = 1,2,3) is not a group action on the 7+1 manifold: its six generators preserve R^(1,3)_LL and act on it as so(3,1), but their Lie closure on R^(1,7)_LL has dimension 9 = so(3,1) + so(3) (exact), the extra so(3) rotating only the four directions i e4..i e7 (there the boosts act trivially and the rotations by the twisted spinor action)  -- closure on R^(1,7)_LL: 9, restricted to R^(1,3)_LL: 6
[PASS] N-6  the left multiplications generate Lie groups (their closures, exact): by imaginary bi-octonions, so(8,C) of real dimension 56 (every L_x is complex-linear and preserves the complex norm N_LL); by imaginary octonions on R^8, so(8) (28); by imaginary split units on Os = R^(4,4), so(4,4) (28; each generator preserves the split norm): any linear action of the loop is an action of these LIE groups  -- closures 56 / 28 / 28
[PASS] N-7  the super-loop manifold (N_LL(w) = 1, a complex 7-sphere of real dimension 14) is parallelized by the frame fields X_x(w) = w x (x imaginary), which are tangent (N_LL(w, w x) = 0); their Lie brackets are [X_x, X_y](w) = w [x,y] + 2 [w,x,y] (100 random w, x, y): at w = 1 this is the Malcev bracket, elsewhere the associator makes the structure functions non-constant: the manifold is parallelizable but carries no Lie group structure
[PASS] N-8  the 7-sphere S^7 of unit octonions (the compact real form of the super-loop) is the homogeneous space Spin(7)/G2: the products L_(e_p) L_(e_p') (p < p') span spin(7) (dimension 21, exact) acting on R^8, the stabilizer of 1 has dimension 14 (g2) and the orbit of 1 has dimension 7  -- spin(7): 21, stabilizer: 14, orbit: 7
[PASS] N-9  the split real form: inside the imaginary split octonions the elements e1, i e4, i e5 span a Lie subalgebra whose Killing form has signature (2 positive, 1 negative): sl(2,R) = so(2,1), the Lorentz algebra of 2+1 dimensions; the unit split octonions S^(3,4) contain SL(2,R) = Spin(2,1): on the real 4+4 manifold the super-loop's Lorentz part is the 2+1 Lorentz group, and SO(3,1) requires the complex form  -- signature (2, 1, 0)
\end{verbatim}

\hypertarget{command-superfluid}{%
\section{\texorpdfstring{7. Command
\texttt{superfluid}}{7. Command superfluid}}\label{command-superfluid}}

\textbf{Field content (S-0).} With
\(\mathsf E_{\rm LL}=i\sigma^{\rm LL}_2\), the matrix
\(\psi'_{\rm LL}=\mathsf E_{\rm LL}^{-1}M_{\rm LL}(\mathbf b)^T\) has
two columns (isospin up and down) on which the Lorentz derivations act
as \(D_{\rm LL}(J^j_{\rm LL})=-\tfrac{i}{2}\sigma^{\rm LL}_j\) and
\(D_{\rm LL}(K^j_{\rm LL})=+\tfrac12\sigma^{\rm LL}_j\) (using
\(\mathsf E_{\rm LL}^{-1}(\sigma^{\rm LL}_j)^T\mathsf E_{\rm LL}=-\sigma^{\rm LL}_j\)),
and the isospin generators act on the column index as
\(\tfrac12\tau^{\rm LL}_\iota\). The chiral Dirac matrices are
\[\gamma^0_{\rm LL}=\begin{pmatrix}0&\sigma^{\rm LL}_0\\\sigma^{\rm LL}_0&0\end{pmatrix},\qquad
\gamma^j_{\rm LL}=\begin{pmatrix}0&\sigma^{\rm LL}_j\\-\sigma^{\rm LL}_j&0\end{pmatrix},\]
\(\gamma^{\rm ch}_{\rm LL}=i\gamma^0_{\rm LL}\gamma^1_{\rm LL}\gamma^2_{\rm LL}\gamma^3_{\rm LL}=\mathrm{diag}(-1,-1,1,1)\),
\(C_{\rm LL}=i\gamma^2_{\rm LL}\gamma^0_{\rm LL}=\mathrm{diag}(\mathsf E_{\rm LL},-\mathsf E_{\rm LL})\),
\(C_{\rm LL}\gamma^{\rm ch}_{\rm LL}=-\mathrm{diag}(\mathsf E_{\rm LL},\mathsf E_{\rm LL})\).
The spinor charges are the block \(\gamma^{\rm ch}_{\rm LL}=+1\); for
each of the six Lorentz generators, the commutator of the four matrices
\(\gamma^p_{\rm LL}\) with the spinor generator reproduces the vector
action \(\mathcal V_{\rm LL}(\mathbf u)\) of the loop's conjugations on
\(\mathbb R^{1,3}_{\rm LL}\) (24 relations). The field \(\Psi_{\rm LL}\)
is a Dirac isodoublet (8 components).

\textbf{Model.}
\(\mathsf h_{\rm LL}(\vec q_{\rm LL})=\sum_j\alpha^j_{\rm LL}q^j_{\rm LL}+\beta_{\rm LL}m_{\rm LL}\)
with \(\alpha^j_{\rm LL}=\gamma^0_{\rm LL}\gamma^j_{\rm LL}\),
\(\beta_{\rm LL}=\gamma^0_{\rm LL}\) (times \(\sigma^{\rm LL}_0\) on the
isospin index); eigenvalues \(\pm\varepsilon_{\rm LL}\) with
\(\varepsilon_{\rm LL}=(q_{\rm LL}^2+m_{\rm LL}^2)^{1/2}\);
\[\mathcal H^{\rm sf}_{\rm LL}=\sum_{\vec q_{\rm LL}}\Psi^\dagger_{\rm LL}(\mathsf h_{\rm LL}-\mu_{\rm LL})\Psi_{\rm LL}-\tfrac14G_{\rm LL}\sum_\iota(\Pi^\iota_{\rm LL})^\dagger\Pi^\iota_{\rm LL},\qquad
\Pi^\iota_{\rm LL}=\sum_{\vec q_{\rm LL}}\Psi_{\rm LL}(\vec q_{\rm LL})^T\mathsf K^\iota_{\rm LL}\Psi_{\rm LL}(-\vec q_{\rm LL}),\]
momenta below \(\Lambda_{\rm LL}\). \textbf{Channels (S-1, S-2):} the
antisymmetric \(8\times8\) matrices \(\mathsf K\) for which the pair
\(\Psi^T_{\rm LL}\mathsf K\Psi_{\rm LL}\) is Lorentz invariant (for each
of the six Lorentz generators on Dirac \(\otimes\) isospin, the
transpose of the generator times \(\mathsf K\) plus \(\mathsf K\) times
the generator vanishes) are computed as a null space (real nullity 12, a
complex space of dimension 6), spanned by
\(\mathsf K^\iota_{\rm LL}=C_{\rm LL}\gamma^{\rm ch}_{\rm LL}\otimes\mathsf E_{\rm LL}\tau^{\rm LL}_\iota\)
and
\(\mathsf K'^\iota_{\rm LL}=C_{\rm LL}\otimes\mathsf E_{\rm LL}\tau^{\rm LL}_\iota\);
the isospin Casimir acting on them is 2 (a triplet), and the charges of
\(\mathsf K^3_{\rm LL}\),
\(\mathsf K^1_{\rm LL}\mp i\mathsf K^2_{\rm LL}\) are 0 and \(\pm1\).

\textbf{Mean field.}
\(d^\iota_{\rm LL}=-\tfrac12G_{\rm LL}\langle\Pi^\iota_{\rm LL}\rangle\),
\(\mathsf P_{\rm LL}=\sum_\iota d^\iota_{\rm LL}(\mathsf K^\iota_{\rm LL})^\dagger\),
and on the Nambu spinor
\(\mathcal N_{\rm LL}=(\Psi_{\rm LL}(\vec q_{\rm LL}),\Psi^*_{\rm LL}(-\vec q_{\rm LL}))\)
\[\mathsf H_{\rm LL}=\begin{pmatrix}\mathsf h_{\rm LL}(\vec q_{\rm LL})-\mu_{\rm LL}&\mathsf P_{\rm LL}\\\mathsf P_{\rm LL}^\dagger&-(\mathsf h_{\rm LL}(-\vec q_{\rm LL})-\mu_{\rm LL})^T\end{pmatrix},\qquad
\Omega_{\rm LL}=\frac{\lvert\vec d_{\rm LL}\rvert^2}{G_{\rm LL}}-\frac12\int\!\frac{d^3q_{\rm LL}}{(2\pi)^3}\sum_{\rm pos}(\ldots)-(\ldots)_{\vec d_{\rm LL}=0},\]
the sum running over the positive eigenvalues of \(\mathsf H_{\rm LL}\).
\textbf{Spectra (S-3, Theorem SF2):} in the scalar channel the positive
eigenvalues are
\(E^\pm_{\rm LL}=[(\varepsilon_{\rm LL}\pm\mu_{\rm LL})^2+\Delta_{\rm LL}^2]^{1/2}\)
with \(\Delta_{\rm LL}\) replaced by each of the gaps
\((\Delta^\pm_{\rm LL})^2=\lvert\vec d_{\rm LL}\rvert^2(1\pm\zeta_{\rm LL})\),
\(\zeta_{\rm LL}=\lvert\vec d_{\rm LL}\times\vec d^{\,*}_{\rm LL}\rvert/\lvert\vec d_{\rm LL}\rvert^2\),
each twice; in the pseudoscalar channel (unitary \(\vec d_{\rm LL}\))
\((\tilde E^\pm_{\rm LL})^2=\varepsilon_{\rm LL}^2+\mu_{\rm LL}^2+\Delta_{\rm LL}^2\pm2[\mu_{\rm LL}^2\varepsilon_{\rm LL}^2+m_{\rm LL}^2\Delta_{\rm LL}^2]^{1/2}\),
four times each. The program compares these with the eigenvalues of
\(\mathsf H_{\rm LL}\), diagonalized at 60 random points (largest
deviation \(7.11\times10^{-15}\)). The proof: with
\(\mathsf D_{\rm LL}=\sum_\iota d^\iota_{\rm LL}(\mathsf E_{\rm LL}\tau^{\rm LL}_\iota)^\dagger\),
\(\mathsf D_{\rm LL}\mathsf D_{\rm LL}^\dagger=\lvert\vec d_{\rm LL}\rvert^2\sigma^{\rm LL}_0+i\sum_\iota(\vec d_{\rm LL}\times\vec d^{\,*}_{\rm LL})^\iota\tau^{\rm LL}_\iota\),
and \(C_{\rm LL}\gamma^{\rm ch}_{\rm LL}\) maps
\(\mathsf h_{\rm LL}(-\vec q_{\rm LL})^T\) to
\(\mathsf h_{\rm LL}(\vec q_{\rm LL})\) and has the square
\(-\mathsf 1_4\).

\textbf{Gap equation and critical coupling (S-4, S-5).} For a unitary
scalar state
\[\frac1{G_{\rm LL}}=\int_{q_{\rm LL}<\Lambda_{\rm LL}}\frac{d^3q_{\rm LL}}{(2\pi)^3}\Big(\frac1{E^+_{\rm LL}}+\frac1{E^-_{\rm LL}}\Big),\qquad
G^{\rm crit}_{\rm LL}=\frac{2\pi^2}{\Lambda_{\rm LL}(\Lambda_{\rm LL}^2+m_{\rm LL}^2)^{1/2}-m_{\rm LL}^2\ln\frac{\Lambda_{\rm LL}+(\Lambda_{\rm LL}^2+m_{\rm LL}^2)^{1/2}}{m_{\rm LL}}}\]
at \(\mu_{\rm LL}=0\), \(\Delta_{\rm LL}\to0\):
\(G^{\rm crit}_{\rm LL}=19.739209\) for \(m_{\rm LL}=0\) and
\(21.286610\) for \(m_{\rm LL}=0.2\) (\(\Lambda_{\rm LL}=1\)). With a
Fermi surface the slope of \(\ln\Delta_{\rm LL}\) against
\(1/G_{\rm LL}\) at weak coupling is
\(-1/(2\nu^{\rm F}_{\rm LL})=-\pi^2/(k^{\rm F}_{\rm LL}\mu_{\rm LL})\),
\(k^{\rm F}_{\rm LL}=(\mu_{\rm LL}^2-m_{\rm LL}^2)^{1/2}\),
\(\nu^{\rm F}_{\rm LL}=k^{\rm F}_{\rm LL}\mu_{\rm LL}/(2\pi^2)\); the
program solves the gap equation for 40 couplings from \(6\) to \(29.4\)
(\(\mu_{\rm LL}=0.5\), \(m_{\rm LL}=0.2\); gaps from
\(2.568874\times10^{-3}\) to \(0.688144\), data file
\texttt{sl\_sf\_gap\_vs\_coupling.csv}) and takes the slope from the two
weakest couplings: \(-43.0637\) against \(-43.0745\) (ratio \(0.9997\)).

\textbf{Which state condenses (S-6, S-10, S-7).} With the two gaps of a
state of non-unitarity \(\zeta_{\rm LL}\), \(\Omega_{\rm LL}\) is the
sum over the gaps of \(\lvert\vec d_{\rm LL}\rvert^2/(2G_{\rm LL})\)
minus the integral of \(E^+_{\rm LL}+E^-_{\rm LL}\) at that gap,
measured from the normal state; the program minimizes it over
\(\lvert\vec d_{\rm LL}\rvert\) for \(\zeta_{\rm LL}=0,0.1,\dots,1\) at
\(G_{\rm LL}=12\), \(\mu_{\rm LL}=0.5\), \(m_{\rm LL}=0.2\) (data file
\texttt{sl\_sf\_omega\_zeta.csv}): the minimum rises from
\(-9.720605\times10^{-5}\) at \(\zeta_{\rm LL}=0\) to
\(-4.860303\times10^{-5}\) at \(\zeta_{\rm LL}=1\) (ratio \(0.500000\)),
as Theorem SF3 states. The scalar channel beats the pseudoscalar one at
\(m_{\rm LL}=0.2\) (\(-9.720605\times10^{-5}\) against
\(-4.685258\times10^{-5}\)), and both give \(-1.384064\times10^{-4}\) at
\(m_{\rm LL}=0\). The stabilizer of a condensate in
\(SU(2)_{\rm v}\times U(1)^{\rm F}_{\rm LL}\) is the null space of the
map that sends each of the four generators to its action on
\(\vec d_{\rm LL}\) (an isospin rotation acts by the cross product with
its axis, the fermion number \(F_{\rm LL}\) by multiplication with
\(2i\)): for the polar state \(\vec d_{\rm LL}\propto(0,0,1)\) and for
the charged state \(\vec d_{\rm LL}\propto(1,i,0)\) it is
one-dimensional, and three generators are broken (data file
\texttt{sl\_sf\_states.csv}).

\textbf{The super-loop singlet (S-8, S-9).}
\((\mathbf b\,\mathbf e_4)(\mathbf b'\mathbf e_4)=-\bar{\mathbf b}'\mathbf b\)
has the scalar part \(-N_{\rm LL}(\mathbf b',\mathbf b)\), invariant
under the Lorentz and isospin automorphisms and symmetric in the two
spinor charges, and an imaginary part in \(\mathfrak l_{\rm LL}\). For
two species,
\(\mathsf K^{\rm s}_{\rm LL}=C_{\rm LL}\gamma^{\rm ch}_{\rm LL}\otimes\mathsf E_{\rm LL}\otimes\mathsf E_{\rm LL}\)
is antisymmetric, and the \(32\times32\) Bogoliubov-de Gennes matrix has
the positive eigenvalues \(E^-_{\rm LL}\) and \(E^+_{\rm LL}\), eight
times each (\(0.347098\) and \(0.878250\) at the test point).

\textbf{Numerics.} The momentum integrals use Gauss-Legendre rules of 40
points on panels whose end points are 0, the cutoff, and 0.5, 0.9, 0.99,
1, 1.01, 1.1 and 1.5 times \(k^{\rm F}_{\rm LL}\) (where these lie
inside the cutoff); the gap equation is solved by Brent's method for
\(\Delta_{\rm LL}\) between \(10^{-9}\) and 20 (tolerance \(10^{-13}\));
the minima of \(\Omega_{\rm LL}\) are found by golden-section search
(120 or 150 steps) on the interval from 0 to three times the gap. The
structural checks S-0, S-1, S-2, S-7, S-8, S-9 use matrices with the
entries 0, \(\pm1\), \(\pm i\), \(\pm\tfrac12\) in double precision,
with tolerances between \(10^{-14}\) and \(10^{-10}\); a rank is the
number of eigenvalues of a Gram matrix above \(10^{-10}\) times the
largest. Data files: \texttt{sl\_sf\_dispersion.csv} (the quasi-particle
branches drawn in panel (a) of the figure \texttt{fig\_sl\_superfluid}),
\texttt{sl\_sf\_channels.csv} (the channels and their quantum numbers).

\begin{verbatim}
Revision_SuperLoop, command `superfluid`: the spin-zero spinor superfluid of the super-loop's spinor charges
11 of 11 checks passed

[PASS] S-0  field content from the bi-octonions: with the isomorphism C (x) H = M_2(C), e_j -> -i sigma_j (j = 1,2,3), the spinor charge bS, as the 2x2 matrix M_LL(bS), gives the Weyl isodoublet psi'_LL = eps2^-1 M_LL(bS)^T, whose column 1 (isospin up) and column 2 (isospin down) are Weyl spinors; on it the Lorentz derivations are D_LL(J_LL^j) = -i sigma_j/2, D_LL(K_LL^j) = +sigma_j/2 (on each column), and the isospin generators are T_LL^iota = tau_iota/2 (mixing the columns); the chiral Dirac matrices gamma_LL^p (p = 0..3) intertwine these spinor generators with the vector action V_LL of the loop's conjugations on R^(1,3)_LL (all 24 relations [gamma_LL^p, D] = sum_p' V_LL^p_p' gamma_LL^p'), the spinor charges being the gamma_ch = +1 block: Psi_LL = Dirac (x) isodoublet
[PASS] S-1  Theorem SF1 (spin-zero channels): the antisymmetric, Lorentz-invariant pair matrices on Dirac (x) isospin form a complex space of dimension 6, spanned by Kch_iota = C_LL gamma_ch (x) eps2 tau_iota (scalar 0+) and Kch'_iota = C_LL (x) eps2 tau_iota (pseudoscalar 0-), iota = 1,2,3; each is unitary and an isospin TRIPLET (Casimir 2): a spin-zero s-wave pair of spinor charges carries isospin 1 and electric charge Q_LL = +1, 0, -1  -- real nullity 12 (= 6 complex dimensions)
[PASS] S-2  electric charge of the spin-zero channels (charge of the pair operator Psi_LL^T Kch Psi_LL, which annihilates two spinor charges): Kch_3 pairs up with down, charge 0; Kch_1 - i Kch_2 and Kch_1 + i Kch_2 pair up-up and down-down, charges +1 and -1  -- Q_LL(Kch_3) = 0, Q_LL(Kch_1 - i Kch_2) = +1, Q_LL(Kch_1 + i Kch_2) = -1
[PASS] S-3  Theorem SF2 (spectra; 60 random momenta, masses, chemical potentials and order parameters dvec in C^3, double precision): in the scalar channel the 8 positive BdG energies are E^+(Delta_+), E^-(Delta_+), E^+(Delta_-), E^-(Delta_-) (each twice) with E^(+-)(Delta) = [(eps +- mu)^2 + Delta^2]^(1/2), eps = (q^2 + m^2)^(1/2), and the two gaps Delta_(+-)^2 = |dvec|^2 (1 +- zeta), zeta = |cross(dvec, dvec^*)| / |dvec|^2 (unitary states: zeta = 0; non-unitary: two gaps); in the pseudoscalar channel (unitary dvec) Etilde^(+-) = [q^2 + m^2 + mu^2 + Delta^2 +- 2 (mu^2 (q^2 + m^2) + m^2 Delta^2)^(1/2)]^(1/2), each four times  -- largest deviation 7.11e-15
[PASS] S-4  critical coupling of the scalar channel at mu = 0 (gap equation 1/G = int d^3q/(2 pi)^3 (1/E^+ + 1/E^-), momentum cutoff Lambda): G_crit = 2 pi^2/Lambda^2 = 19.7392 for m = 0, and G_crit = 2 pi^2/[Lambda (Lambda^2 + m^2)^(1/2) - m^2 ln((Lambda + (Lambda^2 + m^2)^(1/2))/m)] for m > 0 (m = 0.2: quadrature against the closed form)  -- G_crit at m = 0: 19.739209, at m = 0.2: 21.286610
[PASS] S-5  the gap of the unitary scalar state grows monotonically with the coupling (mu = 0.5, m = 0.2, 40 couplings), and at weak coupling ln Delta is linear in 1/G with the BCS slope -1/(2 nuF) = -pi^2/(kF mu), nuF = kF mu/(2 pi^2) the density of states of one species at the Fermi surface (ratio of the slopes within 5 %): a Fermi surface pairs for every G > 0  -- slope -43.0637, BCS value -43.0745, ratio 0.9997
[PASS] S-6  Theorem SF3 (the unitary state wins): at G = 12, mu = 0.5, m = 0.2 the minimum over |dvec| of Omega increases strictly with zeta from the unitary (polar) state zeta = 0 to the fully charged state zeta = 1 (concavity of E^(+-) in Delta^2): the ground state is a unitary isospin-triplet condensate  -- minimum of Omega at zeta = 0: -9.720605e-5, at zeta = 1: -4.860303e-5
[PASS] S-10  which spin-zero channel condenses: at the same coupling (G = 12, mu = 0.5) the minimum of Omega in the scalar (0+) channel lies below that of the pseudoscalar (0-) channel for m = 0.2, and the two minima coincide for m = 0 (chiral symmetry maps one channel onto the other): for massive spinor charges the scalar isospin-triplet condensate is the ground state  -- m = 0.2: -9.720605e-5 (0+) and -4.685258e-5 (0-); m = 0: -1.384064e-4 and -1.384064e-4
[PASS] S-7  symmetry breaking under SU(2)_v x U(1)_F_LL (isospin and fermion number, 4 generators): the polar state dvec = (0,0,1) keeps the rotations about the 3-axis (stabilizer of dimension 1; the electric charge Q_LL = T_LL^3 stays unbroken: a neutral superfluid), the charged state dvec = (1, i, 0) keeps one combination of T_LL^3 and F_LL (stabilizer 1; its condensed pairs Psi_LL^T (Kch_1 - i Kch_2) Psi_LL consist of two isospin-up spinor charges, charge Q_LL = +1); both break 3 generators  -- polar: broken 3, stabilizer 1, pair charge 0; charged: broken 3, stabilizer 1, pair charge +1
[PASS] S-8  the octonionic Cooper pair of two spinor-charge fields: (bS e4)(bS' e4) = -bS'bar bS splits into the scalar part -N_LL(bS', bS) (Lorentz scalar AND isospin singlet: spin 0, isospin 0, charge 0: the super-loop singlet) and the imaginary part, which lies in the Lorentz part l_LL (spin 1, isospin 0: the Malcev bracket); the scalar part is SYMMETRIC in the two fields, so it vanishes for one fermion field and needs two species (Theorem SF4)
[PASS] S-9  two-species singlet: for one field the isospin-singlet scalar C_LL gamma_ch (x) eps2 is symmetric (the pair vanishes); for two species the pair matrix Kch_singlet = C_LL gamma_ch (x) eps2 (x) eps2 (spin 0, isospin 0, species singlet: the third factor acts on the species index; antisymmetric as fermions require) gives the 16 positive BdG energies E^- (8 times) and E^+ (8 times) with E^(+-) = [(eps +- mu)^2 + Delta^2]^(1/2): a fully gapped super-loop-singlet superfluid that breaks only the fermion number and keeps isospin, electric charge and Lorentz invariance  -- E^- = 0.347098, E^+ = 0.878250
\end{verbatim}

\hypertarget{command-renorm}{%
\section{\texorpdfstring{8. Command
\texttt{renorm}}{8. Command renorm}}\label{command-renorm}}

\textbf{Power counting (R-1, R-2).} In \(D+1\) space-time dimensions
\([\Psi_{\rm LL}]=D/2\), \([\phi_{\rm LL}]=(D-1)/2\),
\([G_{\rm LL}]=1-D\), \([g^{\rm Y}_{\rm LL}]=(3-D)/2\),
\([\lambda^{(1)}_{\rm LL}]=[\lambda^{(2)}_{\rm LL}]=3-D\) (exact
half-integers). For a one-particle-irreducible graph,
\(n^{\rm loop}_{\rm LL}=n^{\rm int,F}_{\rm LL}+n^{\rm int,B}_{\rm LL}-n^{\rm v}_{\rm LL}+1\)
and
\(\deg_{\rm LL}=(D+1)\,n^{\rm loop}_{\rm LL}-n^{\rm int,F}_{\rm LL}-2\,n^{\rm int,B}_{\rm LL}\);
with the counts of line ends at the vertices this gives
\[\deg_{\rm LL}=(D+1)-\tfrac D2n^{\rm ext,F}_{\rm LL}+(D-1)\,n^{\rm v}_{\rm LL}\quad\text{(contact)},\]
\[\deg_{\rm LL}=(D+1)-\tfrac D2n^{\rm ext,F}_{\rm LL}-\tfrac{D-1}2n^{\rm ext,B}_{\rm LL}+(D-3)\big(\tfrac12n^{\rm v,Y}_{\rm LL}+n^{\rm v,4}_{\rm LL}\big)\quad\text{(Yukawa and quartic)}.\]
The program compares both formulas with the topological count on all
graphs with up to six vertices (data file
\texttt{sl\_power\_counting.csv}).

\textbf{The contact coupling (R-3).} At the fixed gap
\(\Delta_{\rm LL}=0.05\) (\(\mu_{\rm LL}=0.5\), \(m_{\rm LL}=0.2\)) the
gap equation gives \(1/G_{\rm LL}\) as an integral, evaluated with
Gauss-Legendre rules of 24 points on panels with end points at 0.5, 0.9,
0.99, 1, 1.01, 1.1 and 2 times \(k^{\rm F}_{\rm LL}\) and further panels
whose lengths double up to the cutoff (data file
\texttt{sl\_contact\_running.csv}):

\begin{longtable}[]{@{}
  >{\raggedright\arraybackslash}p{(\columnwidth - 4\tabcolsep) * \real{0.3333}}
  >{\raggedright\arraybackslash}p{(\columnwidth - 4\tabcolsep) * \real{0.3333}}
  >{\raggedright\arraybackslash}p{(\columnwidth - 4\tabcolsep) * \real{0.3333}}@{}}
\toprule\noalign{}
\begin{minipage}[b]{\linewidth}\raggedright
\(\Lambda_{\rm LL}\)
\end{minipage} & \begin{minipage}[b]{\linewidth}\raggedright
\(G_{\rm LL}\)
\end{minipage} & \begin{minipage}[b]{\linewidth}\raggedright
\(G_{\rm LL}\Lambda_{\rm LL}^2/(2\pi^2)\)
\end{minipage} \\
\midrule\noalign{}
\endhead
\bottomrule\noalign{}
\endlastfoot
1 & 10.258986 & 0.519726 \\
2 & 3.737871 & 0.757451 \\
4 & 1.121104 & 0.908732 \\
8 & 0.299414 & 0.970782 \\
16 & 0.076436 & 0.991311 \\
32 & 0.019228 & 0.997505 \\
64 & 0.004816 & 0.999298 \\
128 & 0.001205 & 0.999805 \\
\end{longtable}

\textbf{The bubble of the Yukawa completion (R-4).}
\[I^{\rm bub}_{\rm LL}=\int_{k^{\rm E}_{\rm LL}<\Lambda_{\rm LL}}\frac{d^4k^{\rm E}_{\rm LL}}{(2\pi)^4}\,\frac{k^{\rm E}_{\rm LL}\cdot(k^{\rm E}_{\rm LL}+p^{\rm E}_{\rm LL})+m_{\rm LL}^2}{[(k^{\rm E}_{\rm LL})^2+m_{\rm LL}^2][(k^{\rm E}_{\rm LL}+p^{\rm E}_{\rm LL})^2+m_{\rm LL}^2]}\]
is integrated in four-dimensional spherical coordinates: the length of
\(k^{\rm E}_{\rm LL}\) on panels that grow by the factor 1.5 from 0.25
up to the cutoff, and the angle between \(k^{\rm E}_{\rm LL}\) and
\(p^{\rm E}_{\rm LL}\) from 0 to \(\pi\) with the weight \(4\pi\) times
the square of its sine (the volume of the three-sphere of directions),
each with Gauss-Legendre rules of 48 points; \(p^{\rm E}_{\rm LL}=1\),
\(m_{\rm LL}=0.5\) (data file \texttt{sl\_yukawa\_bubble.csv}):

\begin{longtable}[]{@{}
  >{\raggedright\arraybackslash}p{(\columnwidth - 4\tabcolsep) * \real{0.3333}}
  >{\raggedright\arraybackslash}p{(\columnwidth - 4\tabcolsep) * \real{0.3333}}
  >{\raggedright\arraybackslash}p{(\columnwidth - 4\tabcolsep) * \real{0.3333}}@{}}
\toprule\noalign{}
\begin{minipage}[b]{\linewidth}\raggedright
\(\Lambda_{\rm LL}\)
\end{minipage} & \begin{minipage}[b]{\linewidth}\raggedright
\(I^{\rm bub}_{\rm LL}(0)\)
\end{minipage} & \begin{minipage}[b]{\linewidth}\raggedright
\(I^{\rm bub}_{\rm LL}(p^{\rm E}_{\rm LL})-I^{\rm bub}_{\rm LL}(0)\)
\end{minipage} \\
\midrule\noalign{}
\endhead
\bottomrule\noalign{}
\endlastfoot
10 & 0.62376810 & -0.01583480 \\
20 & 2.52134855 & -0.02021827 \\
40 & 10.11824336 & -0.02460619 \\
80 & 40.51240394 & -0.02899523 \\
160 & 162.09562965 & -0.03338454 \\
\end{longtable}

The slope of the last column against \(\ln\Lambda_{\rm LL}\) between the
two largest cutoffs is \(-6.332440\times10^{-3}\), against
\(-(p^{\rm E}_{\rm LL})^2/(16\pi^2)=-6.332574\times10^{-3}\) (ratio
\(0.99998\)).

\textbf{The couplings of the Yukawa completion (R-5).} The quartic
polynomials of degree 2 in \(\phi_{\rm LL}\) and 2 in its complex
conjugate are combinations of 36 monomials (6 symmetric products of two
components times 6); invariance under the three isospin rotations is a
linear system of rank 34, so the invariants form a space of dimension 2,
spanned by \((\phi^\dagger_{\rm LL}\phi_{\rm LL})^2\) and
\(\lvert\phi^T_{\rm LL}\phi_{\rm LL}\rvert^2\), both verified exactly.
\textbf{Loop and closure (R-6).} A quadratic form on the real
16-dimensional \(\mathbb B_{\rm LL}\) is invariant under a generator if
the transpose of the generator times the symmetric matrix of the form,
plus that matrix times the generator, vanishes. For the 14 left
multiplications \(L^{\rm LL}_{\mathbf x}\) and for the 56 elements of
their closure \(\mathfrak{so}(8,\mathbb C)\) the invariant forms are, in
both cases, the two-dimensional space of \(\mathrm{Re}\,N_{\rm LL}\) and
\(\mathrm{Im}\,N_{\rm LL}\).

\begin{verbatim}
Revision_SuperLoop, command `renorm`: power counting and renormalizability of the super-loop theories
6 of 6 checks passed

[PASS] R-1  mass dimensions in D + 1 space-time dimensions (D = number of space dimensions): [Psi_LL] = D/2, [phi] = (D-1)/2, contact coupling [G] = 1 - D, Yukawa [g_Y] = (3-D)/2, quartic [lam4] = 3 - D; D = 3: [G] = -2 (non-renormalizable), [g_Y] = [lam4] = 0 (renormalizable); D = 7 (7+1, and 4+4, where power counting sees the total dimension 8): [G] = -6, [g_Y] = -2, [lam4] = -4 (all non-renormalizable)  -- twice ([Psi_LL], [phi], [G], [g_Y], [lam4]): D = 3 (3, 2, -4, 0, 0), D = 7 (7, 6, -12, -4, -8)
[PASS] R-2  superficial degree of divergence deg (topological counting nloop = nint_F + nint_B - nvert + 1, deg = (D+1) nloop - nint_F - 2 nint_B, compared with the formulas on all diagrams with up to 6 vertices): contact theory deg = (D+1) - D next_F/2 + (D-1) nvert, Yukawa-quartic theory deg = (D+1) - D next_F/2 - (D-1) next_B/2 + (D-3)(nvert_Y/2 + nvert_4); for D = 3 the contact theory's deg grows by 2 per vertex (non-renormalizable) and the Yukawa theory's does not depend on the vertices (renormalizable); for D = 7 both grow (non-renormalizable)
[PASS] R-3  the contact coupling must depend on the cutoff like Lambda^-2: at a fixed physical gap (Delta = 0.05, mu = 0.5, m = 0.2) the bare coupling solving the gap equation obeys G Lambda^2/(2 pi^2) -> 1 (within 2 % at Lambda = 128): the cutoff cannot be removed with a finite number of logarithmic counterterms  -- G Lambda^2/(2 pi^2) at Lambda = 128: 0.99981
[PASS] R-4  Yukawa completion (pair field phi coupled to Psi_LL^T Kch Psi_LL): the one-loop bubble Ibub(pE) = int d^4kE/(2 pi)^4 [kE.(kE+pE) + m^2]/((kE^2 + m^2)((kE+pE)^2 + m^2)) (Euclidean, four-dimensional cutoff Lambda; pE = 1, m = 0.5) has a quadratically growing part Ibub(0) (absorbed by the mass counterterm) and a logarithmic momentum dependence: d[Ibub(pE) - Ibub(0)]/d ln Lambda -> -pE^2/(16 pi^2) (absorbed by the wave-function counterterm); no new structure is needed  -- slope -6.332440e-3, -pE^2/(16 pi^2) = -6.332574e-3
[PASS] R-5  the renormalizable completion has finitely many couplings: the quartic polynomials of a complex isospin-triplet pair field phi (degree 2 in phi and 2 in phi^*) invariant under the isospin rotations form a space of dimension exactly 2, spanned by (phi^dag phi)^2 and |phi^T phi|^2 (exact; with the Yukawa coupling and the masses this is the complete list)  -- dimension 2
[PASS] R-6  invariance under the loop's linear action equals invariance under its Lie closure: the quadratic forms on the real 16-dimensional B_LL invariant under the 14 left multiplications L_x and those invariant under the 56-dimensional closure so(8,C) form the same 2-dimensional space (Re N_LL and Im N_LL, verified exactly): every counterterm allowed by the closure is allowed by the loop and conversely, so the super-loop pattern of a Lagrangian is as stable under renormalization as the closure symmetry  -- closure 56, invariant forms 2 / 2
\end{verbatim}

\hypertarget{the-examples-of-the-explanation-folder-explanations}{%
\section{\texorpdfstring{9. The examples of the explanation (folder
\texttt{explanations})}{9. The examples of the explanation (folder explanations)}}\label{the-examples-of-the-explanation-folder-explanations}}

The document
\texttt{explanations/The\ spin-zero\ superfluid\ of\ the\ spinor\ charges\ -\ an\ explanation}
explains, word by word, the sentence of the conclusions of the
manuscript on the spin-zero superfluid. Every example in it is computed
by the Rust program \texttt{explanations/code/explain-rs} (binary
\texttt{explain\_superfluid}, 977 lines in \texttt{src/main.rs}; it uses
the complex numbers, the Hermitian diagonalization and the writer of
data files of the library crate \texttt{code/spacetimes-rs}). It must be
started in the folder \texttt{explanations}, into whose \texttt{data/}
it writes. It builds the Pauli matrices, \(\mathsf E_{\rm LL}\), the
chiral Dirac matrices, \(C_{\rm LL}\) and \(\gamma^{\rm ch}_{\rm LL}\),
the generators of the Lorentz group and of isospin on the field
\(\Psi_{\rm LL}\) (and on two species), the spin-zero pair matrices,
seven example order parameters and 200 random ones, and the
bi-quaternions with the Cayley-Dickson product, and verifies:

\begin{itemize}
\tightlist
\item
  X-1 the exchange signs of the Dirac, isospin and species factors, and
  that the sign of every Kronecker product is the product of the signs
  of its factors; X-2 that a pair of anticommuting components keeps only
  the antisymmetric part \(\tfrac12(\mathsf K-\mathsf K^T)\) of its
  matrix, and the number of surviving terms of every spin-zero pair of
  one field; X-3 the Lorentz invariance and the isospin (Casimir 2 or 0)
  of the eight spin-zero pair matrices; X-4 the charges \(+1\), 0,
  \(-1\) of the triplet pairs;
\item
  X-5 that the isospin wavefunction of the condensed pairs is
  \(\mathsf D_{\rm LL}=\sum_\iota d^\iota_{\rm LL}(\mathsf E_{\rm LL}\tau^{\rm LL}_\iota)^\dagger\),
  the fractions of pairs of each charge, the average charge, and the
  average isospin
  \(i\,\vec d_{\rm LL}\times\vec d^{\,*}_{\rm LL}/\lvert\vec d_{\rm LL}\rvert^2\)
  of length \(\zeta_{\rm LL}\); X-6 the two gaps
  \((\Delta^\pm_{\rm LL})^2=\lvert\vec d_{\rm LL}\rvert^2(1\pm\zeta_{\rm LL})\)
  and the unpaired direction of the charged states;
\item
  X-7 the action of the four generators of
  \(SU(2)_{\rm v}\times U(1)^{\rm F}_{\rm LL}\) on \(\vec d_{\rm LL}\),
  derived from their action on the pair matrices, and the unbroken
  generators of the examples;
\item
  X-8 the matrix \(M_{\rm LL}\) of a bi-quaternion, the
  non-associativity, the multiplicative norm, the automorphisms
  \(\varphi^{\rm LL}_{\mathbf g,\mathbf v}\), the product of two spinor
  charges, the singlet
  \(N_{\rm LL}(\mathbf b',\mathbf b)=\tfrac12\mathrm{tr}(M_{\rm LL}(\mathbf b)^T\mathsf E_{\rm LL}M_{\rm LL}(\mathbf b')\mathsf E_{\rm LL}^T)\),
  its symmetry and its invariance; X-9 the pairs of two species, among
  them
  \(\mathsf K^{\rm s}_{\rm LL}=C_{\rm LL}\gamma^{\rm ch}_{\rm LL}\otimes\mathsf E_{\rm LL}\otimes\mathsf E_{\rm LL}\);
  X-10 the counting of the two-particle states of spin \(\tfrac12\) and
  isospin \(\tfrac12\) (one kind and two species) by the traces of
  projectors.
\end{itemize}

Equalities are tested with the tolerance \(10^{-12}\); the random
numbers come from fixed seeds. The data files are
\texttt{explanations/data/ex\_pair\_matrices.csv},
\texttt{ex\_states.csv} and \texttt{ex\_two\_particle\_counting.csv};
the script \texttt{explanations/code/make\_explain\_figures.py} draws
the figures \texttt{explanations/figures/fig\_ex\_bookkeeping},
\texttt{fig\_ex\_pairs} and \texttt{fig\_ex\_energy} (.pdf and .png),
the last from \texttt{data/sl\_sf\_omega\_zeta.csv} and
\texttt{data/sl\_sf\_dispersion.csv} of the command \texttt{superfluid}.
The output of the last run:

\begin{verbatim}
Revision_SuperLoop, folder explanations, program `explain_superfluid`: the examples of the explanation of the sentence on the spin-zero superfluid of the spinor charges
10 of 10 checks passed

[PASS] X-1  exchange signs: the Dirac factors C_LL gamma_ch (0+) and C_LL (0-) of the spin-zero pairs are antisymmetric (in each chirality block they are eps2, the spin singlet up-down minus down-up), the isospin factors eps2 tau_1, eps2 tau_2, eps2 tau_3 are symmetric (isospin 1) and eps2 is antisymmetric (isospin 0), the species factors eps2 (antisymmetric) and sigma_0 (symmetric); gamma_ch = diag(-1,-1,1,1), C_LL = diag(eps2, -eps2), C_LL gamma_ch = -diag(eps2, eps2); and the exchange sign of each of the 24 Kronecker products of a Dirac, an isospin and possibly a species factor is the product of the signs of its factors  -- signs: Dirac [-1, -1], isospin [1, 1, 1, -1], species [-1, 1]
[PASS] X-2  the Pauli principle in numbers: for anticommuting components the pair Psi_LL^T Kch Psi_LL keeps only the antisymmetric part (Kch - Kch^T)/2 of its matrix; counting the products of two different components with a non-zero coefficient, each of the six isospin-triplet matrices (C_LL gamma_ch or C_LL times eps2 tau_iota) gives 4 terms, and each of the two isospin-singlet matrices (C_LL gamma_ch or C_LL times eps2) gives 0: the spin-zero isospin-singlet pair of one field vanishes identically  -- terms (C_LL gamma_ch with eps2 tau_1, eps2 tau_2, eps2 tau_3, eps2; then C_LL with the same): [4, 4, 4, 0, 4, 4, 4, 0]
[PASS] X-3  the eight spin-zero matrices are Lorentz invariant (generator^T Kch + Kch generator = 0 for the three rotations and the three boosts of Psi_LL), and the isospin Casimir sum_iota of the twofold change under T_LL^iota is 2 Kch = 1(1+1) Kch for the six matrices with eps2 tau_iota (isospin 1, triplet) and 0 for the two matrices with eps2 (isospin 0, singlet)
[PASS] X-4  electric charges of the pairs: the pair operator Psi_LL^T Kch Psi_LL annihilates two particles of total charge Q_LL(Kch), the number with generator^T Kch + Kch generator = Q_LL(Kch) Kch for the generator T_LL^3 = Q_LL: Kch_1 - i Kch_2 (isospin part [[2,0],[0,0]], up with up) has Q_LL(Kch) = +1, Kch_3 (isospin part [[0,-1],[-1,0]], up with down, symmetric) has 0, Kch_1 + i Kch_2 (isospin part [[0,0],[0,-2]], down with down) has -1, and the vanishing singlet C_LL gamma_ch (x) eps2 (up with down, antisymmetric) has 0  -- Q_LL(Kch) = +1, 0, -1, 0
[PASS] X-5  the pairs of seven example condensates (|dvec| = 1) and of 200 random ones: the condensed pairs have the isospin wavefunction -sum_iota d_iota tau_iota eps2 (= sum_iota d_iota eps2 tau_iota^*), which is the isospin part Dpair = sum_iota d_iota (eps2 tau_iota)^dag of the pairing block of the Bogoliubov-de Gennes matrix; the weights of its entries up-up, up-down, down-down are the fractions of pairs with the charges +1, 0, -1; the average charge of a pair, fraction(+1) - fraction(-1), equals the third component of i cross(dvec, dvec^*)/|dvec|^2, the whole vector is the average isospin of a pair, and its length is zeta: unitary states (zeta = 0) have pairs of average charge 0 and average isospin 0, the charged states (zeta = 1) consist of pairs of charge +1 only or -1 only  -- polar along the third axis: charges +1/0/-1 0/1/0, average charge 0, zeta 0; polar along the first axis: charges +1/0/-1 0.5/0/0.5, average charge 0, zeta 0; polar along (1,2,2)/3: charges +1/0/-1 0.277778/0.444444/0.277778, average charge 0, zeta 0; unitary with a phase: charges +1/0/-1 0.18/0.64/0.18, average charge 0, zeta 0; partly non-unitary: charges +1/0/-1 0.9/0/0.1, average charge 0.8, zeta 0.8; charged, up-up pairs: charges +1/0/-1 1/0/0, average charge 1, zeta 1; charged, down-down pairs: charges +1/0/-1 0/0/1, average charge -1, zeta 1
[PASS] X-6  gaps of the examples and of 200 random dvec: the eigenvalues of the 2 x 2 gap matrix times its Hermitian conjugate are |dvec|^2 (1 - zeta) and |dvec|^2 (1 + zeta), so unitary states have the single gap Delta_+ = Delta_- = |dvec| for both isospin directions, and the charged states have Delta_+ = sqrt2 |dvec| and Delta_- = 0, the gapless direction being isospin down for the up-up state and isospin up for the down-down state (those particles are not paired)  -- polar along the third axis: Delta_+ 1, Delta_- 1; polar along the first axis: Delta_+ 1, Delta_- 1; polar along (1,2,2)/3: Delta_+ 1, Delta_- 1; unitary with a phase: Delta_+ 1, Delta_- 1; partly non-unitary: Delta_+ 1.341641, Delta_- 0.447214; charged, up-up pairs: Delta_+ 1.414214, Delta_- 0 (unpaired: down); charged, down-down pairs: Delta_+ 1.414214, Delta_- 0 (unpaired: up)
[PASS] X-7  unbroken generators of SU(2)_v x U(1)_F_LL: the action of each generator on dvec is derived from its action on the pair matrices (generator^T Kch_iota + Kch_iota generator is again a combination of Kch_1, Kch_2, Kch_3; true for all 12 cases), and it is: T_LL^iota changes dvec by minus the cross product of the unit vector iota with dvec (a rotation of dvec about the axis iota), F_LL multiplies dvec by 2 i (20 random dvec); a generator is unbroken if it leaves dvec unchanged; for every unitary example exactly one generator is unbroken, the isospin rotation about the real direction of dvec (for the polar state along the third axis that is T_LL^3 = Q_LL: the electric charge is unbroken); for the charged state with up-up pairs the only unbroken generator is F_LL/2 - T_LL^3, the number of isospin-down particles (they are not paired), and for the down-down state F_LL/2 + T_LL^3; the partly non-unitary state breaks all four generators  -- polar along the third axis: isospin rotation about the axis (0, 0, 1): T_LL^3 = Q_LL; polar along the first axis: isospin rotation about the axis (1, 0, 0); polar along (1,2,2)/3: isospin rotation about the axis (0.333333, 0.666667, 0.666667); unitary with a phase: isospin rotation about the axis (0, 0.6, 0.8); partly non-unitary: none; charged, up-up pairs: F_LL/2 - T_LL^3, the number of isospin-down particles; charged, down-down pairs: F_LL/2 + T_LL^3, the number of isospin-up particles
[PASS] X-8  the singlet of the whole super-loop (200 random pairs of bi-quaternions bS, bS' with Gaussian-integer coefficients): M_LL is multiplicative and det M_LL(bS) = N_LL(bS); N_LL(bS', bS) = (1/2) tr(M_LL(bS)^T eps2 M_LL(bS') eps2^T), a contraction of the two isospin indices with eps2 (isospin 0) and of the two spinor indices with eps2 (spin 0); it is the scalar part of bS'bar bS and minus the scalar part of the bi-octonion product (bS e4)(bS' e4) (Cayley-Dickson product), which is SYMMETRIC under the exchange of the two spinor charges, while the imaginary part is antisymmetric; it is invariant under bS -> vI bS gLbar (vI = e3, gL = 1 + i e1 + e2, and random complex unit gL); examples: e4 e4 = -e0, e4 e5 = e1, e5 e4 = -e1; the product of B_LL is not associative ((e1 e2) e4 = e7 but e1 (e2 e4) = -e7), its norm is multiplicative (N_LL(x y) = N_LL(x) N_LL(y) for 200 random pairs of bi-octonions, so the unit bi-octonions are closed under the product), and varphi_(gL,vI) preserves the product (the same 200 pairs; vI = e3 and vI = (e0 + e1 + e2 + e3)/2)  -- determinant and contraction true, Cayley-Dickson product true, exchange symmetry true, examples true, non-associativity true, multiplicative norm true, automorphism true
[PASS] X-9  two species: the super-loop singlet Kch_singlet = C_LL gamma_ch (x) eps2 (x) eps2 (spin 0, isospin 0, species singlet: species Casimir 0) is antisymmetric (signs -1, -1, -1), Lorentz invariant, of charge 0, and keeps 8 terms; the same Dirac and isospin factors with the species factor sigma_0 (each species paired with itself) give a symmetric matrix and 0 terms, as for one field; the isospin triplet eps2 tau_3 with sigma_0 keeps 8 terms (each species pairs as a triplet) and with the species singlet eps2 vanishes
[PASS] X-10  the nucleon analogy (two particles of spin 1/2 and isospin 1/2 in a relative s-wave, which is symmetric under the exchange of the positions, so that the spin-isospin part must be antisymmetric): of the 16 spin-isospin states the antisymmetric ones are the 3 states of spin 0 and isospin 1 and the 3 states of spin 1 and isospin 0; spin 0 with isospin 0 (1 state) and spin 1 with isospin 1 (9 states) are symmetric and forbidden; with two species (64 states) the state of spin 0, isospin 0 and species singlet is antisymmetric and allowed (counted by the traces of the projectors on total spin, total isospin, total species spin and exchange symmetry)  -- one kind: [(0, 0, 0), (0, 1, 3), (1, 0, 3), (1, 1, 0)]; two species: [(0, 0, 0, 1), (0, 0, 1, 0), (0, 1, 0, 0), (0, 1, 1, 9), (1, 0, 0, 0), (1, 0, 1, 9), (1, 1, 0, 9), (1, 1, 1, 0)]
\end{verbatim}

\hypertarget{the-31-manifold-and-the-renormalizability-of-the-yukawa-form-folder-explanations}{%
\section{\texorpdfstring{10. The \(3+1\) manifold and the
renormalizability of the Yukawa form (folder
\texttt{explanations})}{10. The 3+1 manifold and the renormalizability of the Yukawa form (folder explanations)}}\label{the-31-manifold-and-the-renormalizability-of-the-yukawa-form-folder-explanations}}

The document
\texttt{explanations/The\ 3+1\ manifold\ of\ the\ super-loop\ and\ the\ renormalizability\ of\ the\ Yukawa\ form}
answers the test of the sentence of the conclusions ``In \(3+1\)
dimensions its Yukawa form is renormalizable''. Every number and curve
in it is computed by the Rust program \texttt{explanations/code/r13-rs}
(binary \texttt{r13\_yukawa}, 923 lines in \texttt{src/main.rs}). It
uses the complex numbers, the complex matrices, the Gauss-Legendre rules
and the writer of data files of the library crate
\texttt{code/spacetimes-rs}. It must be started in the folder
\texttt{explanations}, into whose \texttt{data/} it writes. Its
formulas:

\begin{itemize}
\tightlist
\item
  \textbf{The manifold (Z-1, Z-2).} A bi-quaternion
  \(\mathbf b=\sum_{p=0}^3\xi^{\rm LL}_p(\mathbf b)\,\mathbf e_p\) has
  the matrix
  \(M_{\rm LL}(\mathbf b)=\xi^{\rm LL}_0(\mathbf b)\,\sigma^{\rm LL}_0-i\sum_j\xi^{\rm LL}_j(\mathbf b)\,\sigma^{\rm LL}_j\).
  A point
  \(\mathbf x=x^0_{\rm LL}\mathbf e_0+i\sum_jx^j_{\rm LL}\mathbf e_j\)
  of \(\mathbb R^{1,3}_{\rm LL}\) has
  \(M_{\rm LL}(\mathbf x)=\sum_px^p_{\rm LL}\sigma^{\rm LL}_p\). Z-1
  compares the two matrices, tests Hermiticity,
  \(\det M_{\rm LL}(\mathbf x)=N_{\rm LL}(\mathbf x)=(x^0_{\rm LL})^2-\sum_j(x^j_{\rm LL})^2\)
  and the inverse formulas
  \(x^p_{\rm LL}=\tfrac12\mathrm{tr}(\sigma^{\rm LL}_pM_{\rm LL}(\mathbf x))\),
  for the four basis elements and 200 random points. Z-2 computes the
  matrix of the quadratic form by polarization: the entry \((p,p')\) is
  half the difference between \(N_{\rm LL}\) of the sum of the basis
  points \(p\) and \(p'\) and the two values \(N_{\rm LL}\) of the basis
  points alone.
\item
  \textbf{The Lorentz group (Z-3).} For a unit bi-quaternion
  \(\mathbf g\), \(\mathsf L_{\rm LL}(\mathbf g)\) has the entry
  \(\tfrac12\mathrm{Re}\,\mathrm{tr}(\sigma^{\rm LL}_pM_{\rm LL}(\mathbf g)\sigma^{\rm LL}_{p'}M_{\rm LL}(\mathbf g)^\dagger)\)
  in row \(p\) and column \(p'\). A random \(M_{\rm LL}(\mathbf g)\) has
  four entries whose real and imaginary parts are uniform in \([-1,1)\).
  It is accepted if the modulus of its determinant exceeds 0.1, and is
  then divided by a square root of the determinant. The check tests 200
  such matrices: preservation of the metric, determinant 1, the
  time-time entry at least 1,
  \(\mathsf L_{\rm LL}(\mathbf g)=\mathsf L_{\rm LL}(-\mathbf g)\), the
  action on the coordinates, and
  \(\mathsf L_{\rm LL}(\mathbf g\mathbf h)=\mathsf L_{\rm LL}(\mathbf g)\mathsf L_{\rm LL}(\mathbf h)\).
\item
  \textbf{The generators (Z-4).} For
  \(\mathbf u=J^j_{\rm LL}=\tfrac12\mathbf e_j\) and
  \(\mathbf u=K^j_{\rm LL}=\tfrac12i\mathbf e_j\) the program computes
  the matrix \(\mathcal V_{\rm LL}(\mathbf u)\) of
  \(M_{\rm LL}(\mathbf x)\mapsto M_{\rm LL}(\mathbf u)M_{\rm LL}(\mathbf x)+M_{\rm LL}(\mathbf x)M_{\rm LL}(\mathbf u)^\dagger\).
  It compares \(\mathcal V_{\rm LL}(\mathbf u)\) with the central
  difference
  \([\mathsf L_{\rm LL}(\exp(10^{-5}\mathbf u))-\mathsf L_{\rm LL}(\exp(-10^{-5}\mathbf u))]/(2\cdot10^{-5})\),
  and tests the preservation of the metric, the 27 commutation relations
  of \(\mathfrak{so}(3,1)\), and the rank 6 of the Gram matrix of the
  six matrices.
\item
  \textbf{The examples (Z-5).} The sample point \((2,1.2,1.5,1.0)\); the
  boost along \(x^1_{\rm LL}\) with the rapidity 0.8 of \((1,0,0,0)\);
  the rotation about \(x^3_{\rm LL}\) by \(\pi/3\) of \((0,1,0,0)\); the
  rotation by \(2\pi\); the null point
  \(\chi_{\rm LL}\chi_{\rm LL}^\dagger\) of \(\chi_{\rm LL}=(1,0)\) and
  its boost; 200 random spinors; and 200 random unit bi-quaternions
  \(\mathbf g\) for the identity
  \(\mathsf E_{\rm LL}^{-1}M_{\rm LL}(\bar{\mathbf g})^T\mathsf E_{\rm LL}=M_{\rm LL}(\mathbf g)\),
  which makes each column of \(\psi'_{\rm LL}\) a two-component spinor.
\item
  \textbf{Mass dimensions (Z-6).} Twice the dimensions, as integers,
  follow from the condition that every term has the dimension \(D+1\):
  \(2[\Psi_{\rm LL}]=D\), \(2[\phi_{\rm LL}]=D-1\),
  \(2[G_{\rm LL}]=2(D+1)-8[\Psi_{\rm LL}]\),
  \(2[g^{\rm Y}_{\rm LL}]=2(D+1)-2[\phi_{\rm LL}]-4[\Psi_{\rm LL}]\),
  \(2[\lambda^{(1)}_{\rm LL}]=2(D+1)-8[\phi_{\rm LL}]\), for \(D=3\) and
  \(D=7\). They are compared with the closed forms, and with
  \([G_{\rm LL}]=2[g^{\rm Y}_{\rm LL}]-2\), which follows from
  \(G_{\rm LL}=4(g^{\rm Y}_{\rm LL})^2/(m^\phi_{\rm LL})^2\).
\item
  \textbf{Power counting (Z-7).} For the Yukawa form the program
  enumerates every class of graphs: every number
  \(n^{\rm v,Y}_{\rm LL}\) of Yukawa vertices and
  \(n^{\rm v,4}_{\rm LL}\) of quartic vertices with
  \(1\le n^{\rm v}_{\rm LL}\le6\), every even
  \(n^{\rm ext,F}_{\rm LL}\le2n^{\rm v,Y}_{\rm LL}\), and every
  \(n^{\rm ext,B}_{\rm LL}\le n^{\rm v,Y}_{\rm LL}+4n^{\rm v,4}_{\rm LL}\)
  for which the difference is even. The internal lines are
  \(n^{\rm int,F}_{\rm LL}=(2n^{\rm v,Y}_{\rm LL}-n^{\rm ext,F}_{\rm LL})/2\)
  and
  \(n^{\rm int,B}_{\rm LL}=(n^{\rm v,Y}_{\rm LL}+4n^{\rm v,4}_{\rm LL}-n^{\rm ext,B}_{\rm LL})/2\);
  classes with no loop or no external line are skipped. For the contact
  form it enumerates \(1\le n^{\rm v}_{\rm LL}\le6\) and every even
  \(n^{\rm ext,F}_{\rm LL}\) from 2 to \(4n^{\rm v}_{\rm LL}\). It
  compares the direct count
  \(\deg_{\rm LL}=(D+1)n^{\rm loop}_{\rm LL}-n^{\rm int,F}_{\rm LL}-2n^{\rm int,B}_{\rm LL}\)
  with the closed formulas, in doubled integers, for \(D=3\) and
  \(D=7\). A function with \(n^{\rm ext,F}_{\rm LL}\) fermion lines
  (fermion number \(\pm1\) each) and \(n^{\rm ext,B}_{\rm LL}\)
  pair-field lines (\(\pm2\) each) is allowed by the fermion number if
  some choice of the signs adds up to zero.
\item
  \textbf{The bubble (Z-8).} The integrand of \(I^{\rm bub}_{\rm LL}\)
  is averaged over the directions of the loop momentum by a 48-point
  Gauss-Legendre rule in the angle between the two momenta, with the
  weight \(\tfrac2\pi\sin^2\) of that angle. The average is then
  integrated over the length of the loop momentum, with the measure
  \(\lvert k^{\rm E}_{\rm LL}\rvert^3/(8\pi^2)\), by the same rule on
  panels between the points 0, 0.25, 0.375, \ldots{} (each point 1.5
  times the previous one), up to the cutoff. The slope is
  \(\Lambda_{\rm LL}^4/(8\pi^2)\) times the difference of the averages
  at \(p^{\rm E}_{\rm LL}\) and at 0, on the sphere
  \(\lvert k^{\rm E}_{\rm LL}\rvert=\Lambda_{\rm LL}\). The renormalized
  bubble subtracts, under the integral, the average at 0,
  \(1/[(k^{\rm E}_{\rm LL})^2+m_{\rm LL}^2]\), and the term
  \(-(p^{\rm E}_{\rm LL})^2[(k^{\rm E}_{\rm LL})^2+2m_{\rm LL}^2]/(2[(k^{\rm E}_{\rm LL})^2+m_{\rm LL}^2]^3)\).
\end{itemize}

Equalities are tested with the tolerance \(10^{-12}\) unless the
statement of a check names another; the power counting uses exact
integers; the random numbers come from a fixed seed. The data files are
\texttt{explanations/data/ex\_r13\_points.csv},
\texttt{ex\_r13\_orbits.csv}, \texttt{ex\_r13\_dimensions.csv},
\texttt{ex\_r13\_power\_counting.csv},
\texttt{ex\_r13\_divergent\_counts.csv} and
\texttt{ex\_r13\_bubble.csv}. The script
\texttt{explanations/code/make\_r13\_figures.py} draws the figures
\texttt{explanations/figures/fig\_ex\_r13\_coordinates},
\texttt{fig\_ex\_r13\_lorentz}, \texttt{fig\_ex\_r13\_divergences} and
\texttt{fig\_ex\_r13\_bubble} (.pdf and .png) from them. The output of
the last run:

\begin{verbatim}
Revision_SuperLoop, folder explanations, program `r13_yukawa`: the 3+1 manifold R^(1,3)_LL of the super-loop and the renormalizability of the Yukawa form on it
8 of 8 checks passed

[PASS] Z-1  the coordinates of R^(1,3)_LL: the Hermitian bi-quaternion x = x_LL^0 e0 + i (x_LL^1 e1 + x_LL^2 e2 + x_LL^3 e3) with real coordinates has the matrix M_LL(x) = x_LL^0 sigma_0 + x_LL^1 sigma_1 + x_LL^2 sigma_2 + x_LL^3 sigma_3 (M_LL(e0) = sigma_0, M_LL(i e_j) = sigma_j), which is Hermitian; det M_LL(x) = N_LL(x) = (x_LL^0)^2 - (x_LL^1)^2 - (x_LL^2)^2 - (x_LL^3)^2; the coordinates are recovered as x_LL^0 = tr M_LL(x)/2 and x_LL^j = tr(sigma_j M_LL(x))/2  -- the four basis elements and 200 random points
[PASS] Z-2  R^(1,3)_LL is Minkowski space: the symmetric matrix of the quadratic form N_LL = det M_LL(x) in the coordinates x_LL^0..x_LL^3 (by polarization) is diag(1, -1, -1, -1): signature (1 positive, 3 negative)  -- matrix diag(1.000000, -1.000000, -1.000000, -1.000000)
[PASS] Z-3  the Lorentz group of R^(1,3)_LL: for 200 random unit bi-quaternions gL (random matrices M_LL(gL) in SL(2,C): four entries with real and imaginary parts uniform in [-1, 1), accepted if the modulus of the determinant exceeds 0.1, divided by a square root of the determinant; tolerances 1e-10 for the 2 x 2 matrices, 1e-9 for Lmat and the coordinates, 1e-12 for Lmat(gL) against Lmat(-gL)), M_LL(x) -> M_LL(gL) M_LL(x) M_LL(gL)^dag keeps the matrix Hermitian and keeps N_LL(x); its matrix Lmat in the coordinates satisfies Lmat^T diag(1,-1,-1,-1) Lmat = diag(1,-1,-1,-1), det Lmat = 1 and Lmat[0][0] >= 1 (a proper orthochronous Lorentz transformation); gL and -gL give the same Lmat (SL(2,C) is the double cover), and Lmat(gL hL) = Lmat(gL) Lmat(hL) (a group homomorphism)  -- largest change of N_LL: 6.3e-14
[PASS] Z-4  the six generators: the rotation generators J_LL^j = e_j/2 and the boost generators K_LL^j = i e_j/2 of the super-loop have M_LL(J_LL^j) = -i sigma_j/2 and M_LL(K_LL^j) = sigma_j/2; their action M_LL(x) -> M_LL(uL) M_LL(x) + M_LL(x) M_LL(uL)^dag on R^(1,3)_LL gives six real 4 x 4 matrices V_LL(J_LL^j), V_LL(K_LL^j) (equal, within 1e-8, to the derivatives at theta = 0 of Lmat of the finite rotations exp(-i theta sigma_j/2) and boosts exp(theta sigma_j/2), computed by central differences with the step 1e-5), each preserving the metric (V_LL(uL)^T diag(1,-1,-1,-1) + diag(1,-1,-1,-1) V_LL(uL) = 0), linearly independent (the Gram matrix of their 16 entries has rank 6 in a Gaussian elimination with the pivot threshold 1e-9), with the relations of so(3,1): [V_LL(J_LL^j), V_LL(J_LL^j')] = sum_j'' eps3_jj'j'' V_LL(J_LL^j''), [V_LL(J_LL^j), V_LL(K_LL^j')] = sum_j'' eps3_jj'j'' V_LL(K_LL^j''), [V_LL(K_LL^j), V_LL(K_LL^j')] = -sum_j'' eps3_jj'j'' V_LL(J_LL^j'') (two boosts close into a rotation with the minus sign that distinguishes the Lorentz algebra from so(4))  -- rank 6, 27 commutation relations
[PASS] Z-5  worked examples on R^(1,3)_LL: the sample point (2, 1.2, 1.5, 1.0) has M_LL(x) = [[3, 1.2 - 1.5 i], [1.2 + 1.5 i, 1]] and N_LL = -0.69 (spacelike); the boost along x_LL^1 with the rapidity 0.8 takes (1, 0, 0, 0) to (cosh 0.8, sinh 0.8, 0, 0); the rotation about x_LL^3 by pi/3 takes (0, 1, 0, 0) to (0, 1/2, sqrt3/2, 0); the rotation by 2 pi has M_LL(gL) = -sigma_0 and Lmat = identity; a two-component spinor chi gives the null point with M_LL(x) = chi chi^dag (chi = (1, 0): (1/2, 0, 0, 1/2)), and the boost of chi along x_LL^3 with the rapidity 0.8 moves it along the light cone to exp(0.8) (1/2, 0, 0, 1/2); chi chi^dag is a future null point for 200 random chi; the Lorentz transformation gL acts on each column of the Weyl isodoublet psi'_LL = eps2^-1 M_LL(bS)^T of a spinor charge (bS -> bS gLbar) as chi -> M_LL(gL) chi, because eps2^-1 M_LL(gLbar)^T eps2 = M_LL(gL) (200 random gL as in Z-3, tolerance 1e-10)  -- sample point: N_LL = -0.690000; boost: (1.337435, 0.888106, 0, 0); rotation: (0, 0.500000, 0.866025, 0); null points checked: 200; columns checked: 200
[PASS] Z-6  mass dimensions in D + 1 space-time dimensions (every term of the Lagrangian density has the dimension D + 1): [Psi_LL] = D/2, [phi] = (D-1)/2, the contact coupling [G] = 1 - D, the Yukawa coupling [g_Y] = (3-D)/2, the quartic couplings [lam4] = 3 - D; for D = 3 (the manifold R^(1,3)_LL): 3/2, 1, -2, 0, 0 -- the Yukawa form has only couplings of dimension 0 and is renormalizable by power counting, the contact form is not; for D = 7 (R^(1,7)_LL, and the 4+4 manifold, where power counting sees the total dimension 8): 7/2, 3, -6, -2, -4 -- every coupling has a negative dimension; below the mass m_phi of the pair field its exchange gives the contact coupling G = 4 g_Y^2/m_phi^2, consistent with [G] = 2[g_Y] - 2  -- D = 3: 3/2, 1, -2, 0, 0; D = 7: 7/2, 3, -6, -2, -4
[PASS] Z-7  power counting (all classes of graphs with 1 to 6 vertices and at least one loop; a class is a set of numbers next_F, next_B, nvert_Y, nvert_4 for which the line ends at the vertices pair into internal lines, and deg depends only on the class; topological count nloop = nint_F + nint_B - nvert + 1, deg = (D+1) nloop - nint_F - 2 nint_B, compared with the formulas deg = (D+1) - D next_F/2 - (D-1) next_B/2 + (D-3)(nvert_Y/2 + nvert_4) of the Yukawa form and deg = (D+1) - D next_F/2 + (D-1) nvert of the contact form): in 3+1 dimensions the Yukawa form has, for the graphs with at most nvert vertices and every nvert from 2 to 6, exactly four divergent functions that conserve the fermion number (the pair-field self-energy, deg 2; the fermion self-energy, deg 1; the Yukawa vertex, deg 0; the four-point function of the pair field, deg 0) -- each has a term of the Lagrangian density to absorb it; the functions with one or three pair-field lines and no fermion lines (deg 3 and 1) vanish by fermion-number conservation; the number of divergent functions of the contact form grows with the number of vertices, and so does that of both forms for D = 7  -- 408 graph classes compared; Yukawa form, D = 3, divergent and allowed: (next_F, next_B) = (0, 2) deg 2, (next_F, next_B) = (0, 4) deg 0, (next_F, next_B) = (2, 0) deg 1, (next_F, next_B) = (2, 1) deg 0
[PASS] Z-8  the one-loop bubble of the Yukawa form, Ibub(pE) = int_(|kE| < Lambda) d^4kE/(2 pi)^4 [kE.(kE+pE) + m^2]/((kE^2 + m^2)((kE+pE)^2 + m^2)) (pE = 1, m = 0.5, Lambda = 4 to 2048; the angular average of the integrand at pE = 0 equals 1/(kE^2 + m^2) within 1e-14): Ibub(0) equals the closed form (Lambda^2 - m^2 ln(1 + Lambda^2/m^2))/(16 pi^2) (relative tolerance 1e-10) and grows like Lambda^2 (ratio at 2048/1024 within 1e-3 of 4; absorbed by the mass counterterm of the pair field); the cutoff derivative of the momentum-dependent part, computed on the sphere |kE| = Lambda, tends to -pE^2/(16 pi^2) (relative tolerance 1e-5 at Lambda = 2048; a logarithmic divergence, absorbed by the wave-function counterterm); the renormalized bubble Iren = Ibub(pE) - Ibub(0) - pE^2 dIbub/d(pE^2) at pE = 0 converges as Lambda grows: power counting guarantees a remainder falling at least like Lambda^-2 (a factor 4 per doubling); here the term of order pE^4 of the angular average of the integrand is pE^4 m^4/(kE^2 + m^2)^5 (verified at kE = 1, 2, 4 with pE = 0.05, relative tolerance 1e-2), so each doubling divides the change by about 64 (each of the first three ratios between 56 and 72), and from Lambda = 32 on Iren is constant within the rounding errors (relative changes below 1e-8): after the two counterterms nothing depends on the cutoff  -- Ibub(0) ratio at 2048/1024: 4.000011; slope at Lambda = 2048: -6.332574e-3 (target -6.332574e-3); changes of Iren per doubling from Lambda = 4: 3.208453e-8, 4.968769e-10, 7.745994e-12, 1.200301e-13 (ratios 64.572, 64.146, 64.534); Iren at 2048: 1.560666e-3; largest relative deviation from it for Lambda >= 32: 1.1e-9; pE^4 term verified: true
\end{verbatim}

\hypertarget{the-checks-of-the-documents}{%
\section{11. The checks of the
documents}\label{the-checks-of-the-documents}}

The script \texttt{check\_symbols.py} verifies eight statements on the
symbols: LL-S1, no symbol and no plain-text name is registered twice;
LL-S2, the taken entries carry the meanings of the master registries, no
new symbol coincides with a symbol of a master registry, and the symbols
shared with the folder \texttt{Revision\_NoGo} are declared and denote
the same objects there; LL-S3, every column of every data file
(\texttt{data/sl\_*.csv} and \texttt{explanations/data/ex\_*.csv}) is a
registered name or a descriptive word; LL-S4, every document contains
the complete table of symbols; LL-S5, the symbols of the superfluid
\(\mu_{\rm LL}\), \(\Lambda_{\rm LL}\), \(\Omega_{\rm LL}\),
\(\nu^{\rm F}_{\rm LL}\), \(\Delta_{\rm LL}\), \(\zeta_{\rm LL}\),
\(\varepsilon_{\rm LL}\) always carry their suffix; LL-S6, no letter
that the project uses with another meaning occurs in a formula without
the suffix LL; LL-S7, every letter in bold, sans-serif, calligraphic,
blackboard or fraktur type in a formula is registered; LL-S8, every name
with the suffix \texttt{\_LL}, every name of an algebra element, every
coefficient and every pair matrix in a report or data file
(\texttt{data/sl\_*} and \texttt{explanations/data/ex\_*}) is
registered. The script \texttt{check\_self\_contained.py} searches every
document for the wordings by which a text sends its reader to another
document (LL-H1). Their output:

\begin{verbatim}
PROJECT spacetimes, Revision_SuperLoop: checks of the symbols (registry provenance/symbols_ll.csv)
8 of 8 checks passed

[PASS] LL-S1  every symbol text is registered once (87 entries) and every plain-text name of the programs, figures and data files belongs to one entry (310 names)
[PASS] LL-S2  the 2 taken entries carry exactly the meaning of the MASTER registry from which they are taken; none of the 203 items of the 85 new entries is an item of either master registry (503 items compared); the 22 items that the registry of Revision_NoGo also contains are exactly the 22 declared shared ones, each with the same object ($SU(2)$, $U(1)$, $\det$, $\langle\ \rangle$, $\mathbb C$, $\mathbb R$, $\mathfrak g_2$, $\mathfrak{so}(3)$, $\mathfrak{so}(3,1)$, $\mathfrak{so}(4)$, $\mathfrak{so}(4,4)$, $\mathfrak{so}(8)$, $\mathfrak{so}(8,\mathbb C)$, $\mathfrak{su}(2)$, $\mathrm G_2$, $\mathrm{Re}$, $\mathrm{diag}$, $\mathrm{tr}$, $\otimes$, $i$, $p$, ${}^T$)
[PASS] LL-S3  every column of the 23 data files data/sl_*.csv and explanations/data/ex_*.csv is a registered plain-text name or a descriptive word (120 different columns), and every file states its content in a line '# content:'
[PASS] LL-S4  each of the 6 documents of Revision_SuperLoop and the manuscript contains the complete table of the 87 registered symbols, so that every symbol it uses is defined in the document itself
[PASS] LL-S5  in the 3098 formulas of the documents, mu, Lambda, Omega, nu, Delta, zeta and varepsilon always carry the suffix LL
[PASS] LL-S6  in the 3620 formulas of the documents and of the new registry entries, no Latin letter except the index j, D, the imaginary unit i, the index p and the differential d stands without the suffix LL in its scripts, and the Greek letters that are symbols carry the suffix as well
[PASS] LL-S7  every one of the 3010 occurrences of a letter in bold, sans-serif, calligraphic, blackboard or fraktur type in the formulas of the documents is a registered symbol (47 registered styled letters)
[PASS] LL-S8  every one of the 3158 occurrences of a name with the suffix _LL, of an algebra element, of a family coefficient or of a pair matrix in the 32 reports and data files (data/sl_* and explanations/data/ex_*) is a registered plain-text name
\end{verbatim}

\begin{verbatim}
PROJECT spacetimes, Revision_SuperLoop: does any document send its reader to another document? (Hard Rule 7)
1 of 1 checks passed

[PASS] LL-H1  none of the 28 documents and document sources of Revision_SuperLoop refers its reader to another document for an equation, a symbol or a definition

Allowed exceptions:
  README: a provenance document named as a place to read (the README is the list of contents of the folder: naming its documents is its purpose)
\end{verbatim}

\hypertarget{data-files-and-figures}{%
\section{12. Data files and figures}\label{data-files-and-figures}}

\begin{longtable}[]{@{}
  >{\raggedright\arraybackslash}p{(\columnwidth - 4\tabcolsep) * \real{0.3333}}
  >{\raggedright\arraybackslash}p{(\columnwidth - 4\tabcolsep) * \real{0.3333}}
  >{\raggedright\arraybackslash}p{(\columnwidth - 4\tabcolsep) * \real{0.3333}}@{}}
\toprule\noalign{}
\begin{minipage}[b]{\linewidth}\raggedright
data file
\end{minipage} & \begin{minipage}[b]{\linewidth}\raggedright
content (header line of the file)
\end{minipage} & \begin{minipage}[b]{\linewidth}\raggedright
columns
\end{minipage} \\
\midrule\noalign{}
\endhead
\bottomrule\noalign{}
\endlastfoot
\texttt{sl\_associators.csv} & the non-zero associators {[}e\_p1, e\_p2,
e\_p3{]} = (e\_p1 e\_p2) e\_p3 - e\_p1 (e\_p2 e\_p3) of three units with
p1 \textless{} p2 \textless{} p3; each equals coefficient x the unit
result (all other orders follow by alternativity) & \texttt{unit1}
\texttt{unit2} \texttt{unit3} \texttt{result} \texttt{coefficient} \\
\texttt{sl\_contact\_running.csv} & bare contact coupling G versus the
cutoff Lambda at the fixed gap Delta = 0.05 (mu = 0.5, m = 0.2) &
\texttt{Lambda} \texttt{G} \texttt{G\_times\_Lambda2}
\texttt{G\_times\_Lambda2\_over\_2pi2} \\
\texttt{sl\_malcev\_families.csv} & members of the covariant families
verified exhaustively: Jacobi on all basis triples, linearized Malcev on
all basis 4-tuples & \texttt{family} \texttt{coefficients} \texttt{Lie}
\texttt{Malcev} \texttt{failing\_Malcev\_4\_tuples} \\
\texttt{sl\_malcev\_structure.csv} & non-zero commutators of the real
basis elements of the Malcev algebra m\_LL = C (x) Im O: {[}left,
right{]} = coefficient x result (left before right in the basis order
e1, e2, e3, ie1, ie2, ie3, e4..e7, ie4..ie7) & \texttt{left}
\texttt{right} \texttt{result} \texttt{coefficient} \\
\texttt{sl\_manifolds.csv} & manifolds that carry the Lorentz super-loop
& \texttt{manifold} \texttt{real\_dimension} \texttt{signature}
\texttt{role} \\
\texttt{sl\_multiplication\_table.csv} & products of the units e0..e7 of
B\_LL (row: left factor, column: right factor); the bi-octonions are the
complex combinations of the units & \texttt{left} \texttt{e0}
\texttt{e1} \texttt{e2} \texttt{e3} \texttt{e4} \texttt{e5} \texttt{e6}
\texttt{e7} \\
\texttt{sl\_power\_counting.csv} & superficial degree of divergence deg:
theory 0 = contact (next\_F = 4), theory 1 = Yukawa (next\_F = 2,
next\_B = 0), versus the number of vertices nvert, for D = 3 and D = 7
space dimensions & \texttt{D} \texttt{theory} \texttt{nvert}
\texttt{next\_F} \texttt{next\_B} \texttt{deg} \\
\texttt{sl\_quantum\_numbers.csv} & quantum numbers of the sectors of
B\_LL under the operators of the super-loop: Q\_LL (eigenvalues with
multiplicities), T\_LL\^{}2, Cas\_LL, Lorentz type, spinor number exp(2
pi D\_LL(J\_LL\^{}3)), twist charge Theta\_LL & \texttt{sector}
\texttt{real\_dimension} \texttt{Q\_LL} \texttt{T\_LL\^{}2}
\texttt{Cas\_LL} \texttt{Lorentz\_type} \texttt{spinor\_number}
\texttt{Theta\_LL} \\
\texttt{sl\_sf\_channels.csv} & pair channels of the spinor charges:
pair matrix, spin and parity, isospin, charge Q\_LL of the pair, remark
& \texttt{pair\_matrix} \texttt{spin\_parity} \texttt{isospin}
\texttt{pair\_charge} \texttt{remark} \\
\texttt{sl\_sf\_dispersion.csv} & quasiparticle energies versus q at G =
12, mu = 0.5, m = 0.2: polar state E\^{}-, E\^{}+; charged state (zeta =
1) gapped branch E\^{}-(sqrt(2) Delta) {[}its other branch is the normal
\textbar eps - mu\textbar{]}; normal state \textbar eps - mu\textbar;
pseudoscalar channel Etilde\^{}-, Etilde\^{}+ at the same Delta &
\texttt{q} \texttt{E\_minus\_polar} \texttt{E\_plus\_polar}
\texttt{E\_minus\_charged\_gapped} \texttt{E\_normal}
\texttt{Etilde\_minus} \texttt{Etilde\_plus} \\
\texttt{sl\_sf\_gap\_vs\_coupling.csv} & unitary scalar (0+)
isospin-triplet state: gap Delta and Omega(Delta) - Omega(0) versus the
coupling G (cutoff Lambda = 1, mu = 0.5, m = 0.2); G\_over\_Gcrit\_m0 is
G divided by G\_crit at mu = 0, m = 0 & \texttt{G}
\texttt{G\_over\_Gcrit\_m0} \texttt{Delta}
\texttt{Omega\_minus\_Omega0} \\
\texttt{sl\_sf\_omega\_zeta.csv} & minimum over \textbar dvec\textbar{}
of Omega(dvec) - Omega(0) at fixed zeta = \textbar cross(dvec,
dvec\^{}*)\textbar/\textbar dvec\textbar\^{}2 (G = 12, mu = 0.5, m =
0.2), the \textbar dvec\textbar{} of the minimum, and Omega at
\textbar dvec\textbar{} = the gap of the polar state & \texttt{zeta}
\texttt{abs\_dvec\_at\_min} \texttt{Omega\_min}
\texttt{Omega\_at\_polar\_gap} \\
\texttt{sl\_sf\_states.csv} & spin-zero isospin-triplet states of the
superfluid of the spinor charges: order parameter dvec, zeta, dimension
of the stabilizer in SU(2)\_v x U(1)\_F\_LL, broken generators, charge
of the condensed pairs, gaps & \texttt{state} \texttt{dvec}
\texttt{zeta} \texttt{stabilizer\_dim} \texttt{broken\_generators}
\texttt{pair\_charge} \texttt{gaps} \\
\texttt{sl\_yukawa\_bubble.csv} & Euclidean bubble Ibub of the Yukawa
completion versus the cutoff Lambda at pE = 1, m = 0.5; Ibub(0) grows
like Lambda\^{}2, Ibub(pE) - Ibub(0) like -(pE\^{}2/16 pi\^{}2) ln
Lambda & \texttt{Lambda} \texttt{Ibub\_0} \texttt{Ibub\_pE}
\texttt{Ibub\_pE\_minus\_Ibub\_0} \\
\end{longtable}

\begin{longtable}[]{@{}
  >{\raggedright\arraybackslash}p{(\columnwidth - 4\tabcolsep) * \real{0.3333}}
  >{\raggedright\arraybackslash}p{(\columnwidth - 4\tabcolsep) * \real{0.3333}}
  >{\raggedright\arraybackslash}p{(\columnwidth - 4\tabcolsep) * \real{0.3333}}@{}}
\toprule\noalign{}
\begin{minipage}[b]{\linewidth}\raggedright
figure (\texttt{figures/}, .pdf and .png)
\end{minipage} & \begin{minipage}[b]{\linewidth}\raggedright
what it shows
\end{minipage} & \begin{minipage}[b]{\linewidth}\raggedright
data
\end{minipage} \\
\midrule\noalign{}
\endhead
\bottomrule\noalign{}
\endlastfoot
\texttt{fig\_sl\_tables} & (a) the products of the units of
\(\mathbb B_{\rm LL}\); (b) the triples of imaginary units by sector and
the number with a non-zero associator &
\texttt{sl\_multiplication\_table.csv}, \texttt{sl\_associators.csv} \\
\texttt{fig\_sl\_superloop} & the Malcev algebra of the super-loop, its
brackets and its outer operators & the results of checks L and Q \\
\texttt{fig\_sl\_charges} & the quantum numbers of the three sectors of
\(\mathbb B_{\rm LL}\) & \texttt{sl\_quantum\_numbers.csv} \\
\texttt{fig\_sl\_classification} & Theorem M: what the bracket of two
spinor charges can contain, by family & the results of checks G and M \\
\texttt{fig\_sl\_manifolds} & the manifolds that carry the super-loop &
the results of checks N \\
\texttt{fig\_sl\_superfluid} & (a) quasi-particle energies; (b) the gap
versus the coupling; (c) the weak-coupling law; (d) the minimum of
\(\Omega_{\rm LL}\) versus \(\zeta_{\rm LL}\) &
\texttt{sl\_sf\_dispersion.csv}, \texttt{sl\_sf\_gap\_vs\_coupling.csv},
\texttt{sl\_sf\_omega\_zeta.csv} \\
\texttt{fig\_sl\_renorm} & (a) the superficial degree of divergence; (b)
the contact coupling versus the cutoff; (c) the bubble &
\texttt{sl\_power\_counting.csv}, \texttt{sl\_contact\_running.csv},
\texttt{sl\_yukawa\_bubble.csv} \\
\end{longtable}

\hypertarget{checksums}{%
\section{13. Checksums}\label{checksums}}

SHA-256 of the data files, the figures and the programs from which the
documents were built (the last block: the folder \texttt{explanations}):

\begin{verbatim}
280787dfd4c8336cb63a314a332559cdb91ab4de4aacfa9f3547ab748d1d9a22  sl_algebra_checks.txt
b0c7312ba54c997c485492d313dcba905f2e4fb09ed9cb1222382579acd5f3e3  sl_associators.csv
290e26ae38ad95f39199a10c75264cb0f326a75cb8973254aad5a0ba32da9073  sl_contact_running.csv
3239d96fd7eeafb436fea956d6ad4e73760b7f205a701b05ddfb5ebce9131008  sl_groebner_bases.txt
b4d99fc5fd7eefce05bbac017563ab864e787b80663915c5cbf50574cff318a6  sl_groebner_checks.txt
e26d9914e9d7a30eb9e75ea3e8776b9f87516ccaae225631666ad23e511e238d  sl_malcev_checks.txt
dff39dffa2f958ca1d0c2d5944393b5756db4b2816c5043334ab1af0a60d86f4  sl_malcev_families.csv
34e8d33fc5e375bcb35f4023ba732d622db74b530e651caea7e9603b6ada941a  sl_malcev_structure.csv
ab062b4a66b8a38abe410ffafdb0a95521cf30ba4e095eca6053d625627c323d  sl_manifolds.csv
13336d96de70487a043d6c5f06d5536f16269b1e19157ee445189847321ec16c  sl_manifolds_checks.txt
b1cfd3dfc2434a613909e50cc3a23ca0edfb7b16007129f6137e53c3755f53a0  sl_multiplication_table.csv
1892ba5e74d6afc965fce4dfd36b5b8d8e28f715e4e57b4b8f533c5dd42912d4  sl_power_counting.csv
4014ea2dd2f73847097be883f5ee73a6329f08e689f6e6784be3161c93cdf18f  sl_quantum_numbers.csv
4e18b556b2e26a20b3946694e192780fcafabdca0c92da4ede8835cb23960014  sl_renorm_checks.txt
0f3635dcb9d1f5156c2cb2183192a9600bcb34430f1285f28c8c7dac7a40b4f5  sl_sf_channels.csv
dea8b2c8284a70334768cc492bdd32d9d0ac7c723abcb31a1f2e95c8e3459495  sl_sf_dispersion.csv
8bd27baccf82bd3fe666951b10ee99e81b6de1d1299a1362bf774bc1fb83527b  sl_sf_gap_vs_coupling.csv
6dbce6906928d983723928c990f61e9fe9aab72c52de96dfe7bca9f2d7f6edd3  sl_sf_omega_zeta.csv
27d74793da4c0e42cbc1f0cf7a4bb0ef39f5b2e8dfc52bb7c836303bca86cb8a  sl_sf_states.csv
f6692aa0efd3b52ef97406c4f0853f0ea2163434526398653004f979f4657c62  sl_superfluid_checks.txt
171af3b33706ea76db8b74c266d31f86674d8e34ac0f3d61e244287cb7ae1508  sl_yukawa_bubble.csv
2bc7a676537fa73fd08c8d30cc057d2055ebf53df2f813ed4e72d6cc71afe205  fig_sl_charges.pdf
ac70bdd826a0d65c59fb79b94d300393fc68015bbb4d0189fc053fe8e0802980  fig_sl_classification.pdf
aa7132242af09217d639b821c1c4acd1e79eb49085c8db55628ad08051db21bf  fig_sl_manifolds.pdf
e0f3c69f58033cc4248d9dee033dd0ebc19e2e520d0a2e243ce93f33f1725497  fig_sl_renorm.pdf
0a76d5ed9e6d22ddd2e0ab0c615562514220623621094a885f19a600e68481ed  fig_sl_superfluid.pdf
e02fc2cb28f14725fc0605e0ee87b6d85b3e99fc5fc8403722cf7ab6633a6e24  fig_sl_superloop.pdf
4f73b86dd47fe47973a6cc16e0ba68e9f50049c7510c484cacd75fbe0b52240d  fig_sl_tables.pdf
bd886920b5c69449d0bfed2064654e91f6b6b8bcf98c75eeb2dec38a9538ed67  code/superloop-rs/src/algebra.rs
01e0867bd3f3bad05b555249fe79732fb71cdcb909fe16b9885f3dca734fcdea  code/superloop-rs/src/bioct.rs
ecc3d05e48b6a49eefe1fae96dcf340d7c54ca275d995db30ad018e7d651e2cd  code/superloop-rs/src/checks.rs
cc2abe0ad40727da94e17b2cc78891b41b564bde4ba56f460a5f7319f2ad6e78  code/superloop-rs/src/exact.rs
124f653bbc959e6d591e693dcbbb3812cd564dbd00047652029255c65ced70d9  code/superloop-rs/src/main.rs
def7dbfc144559e324d3337cd16d851fa0b0d8a17114fe81eadf830581af1a0e  code/superloop-rs/src/malcev.rs
d9f46fbb4a2bcbccfd2107fd6fbf41091e78352c1346d943a9ce6385423dba04  code/superloop-rs/src/manifolds.rs
b0097ad508f55fe34ce5d8ad72c4ac7a3d7cf7d919661ca1422d1dbc6b747d88  code/superloop-rs/src/renorm.rs
61a6fb93c9591631e820c538d003a551447945a2bafe2ae84652dcd056383765  code/superloop-rs/src/superfluid.rs
84ce114dda4a896f3a7286572de3301961e6f1d82cd208379b093187e64a4c4c  code/superloop-rs/Cargo.toml
f3abd67e3bc8b707583adcb4dd5aa5bdb5ec233237bb115dcf1d1f4163df436f  code/scripts/assemble_docs.py
fd9de7ac8794a7d0316d7c07d791807a180ce1eccfa53c19eed8f69ba4022ea1  code/scripts/build_manuscript.py
465063731c88f51a1d8af013d40f1cae7613007e63d3256e5ab26b364bd3b455  code/scripts/check_layout.py
864b86dfce7e615b250577659cda8226f4e1bea4e843e01fe05c75565943f353  code/scripts/check_self_contained.py
c2ab05157dbe7984d93e48c6de4a7d06a0ed80285c92d1df1d2278cf8cb8d235  code/scripts/check_symbols.py
d915832d77bfc194020e1bbed3740c8420e309cbf54832b2e5b1439a1b4cdc27  code/scripts/make_figures.py
cf5b78e608bd8ed2c880ca46425d5b6a5daca31331ca90af1697e5885d95e621  code/scripts/malcev_groebner.py
cfdc99c1e280e6956e641694c6105e6758d2cb38a6a8054aa4c4f7da11c21eb3  code/scripts/tex2md.py
ac90dfe506b51836429d6a76c798e846cdc8646293675015af0e49baab506303  code/scripts/write_doc_tables.py
d97abaf12df2b47dd67550b98d865d5e08b8b8a95d9199a5666c78c708dca02d  code/scripts/write_symbol_registry.py
783183f2a23f1fb50b3b430d6243f2f214e6e40d7b1e4ec971a59ad0a9f161f9  code/scripts/build_docs.sh
3252f0a1d90606c4d1255d917ee7d5856541f44f8513b89cdd7ea010dbd67d6e  code/scripts/check_all.sh
cad52ccfbd5b20a51f59cd044cc1bfb05cae4268d45581539f66d47145b7e24d  code/scripts/doc_header.tex
\end{verbatim}

\begin{verbatim}
3c4366a5d11a44c8b8dd21d05fceb379646f55db407cc41321c5501e595fe4d8  explanations/data/ex_pair_matrices.csv
f35b57b83de79a1de0b45f6f1d73c8be77ddf5ba0f13799460fb62039572687e  explanations/data/ex_r13_bubble.csv
71bc8ac7fb68d7be2ea856bb01d377bd41daf1d79d158b421542bd6bd6ed12f6  explanations/data/ex_r13_checks.txt
50cb49b9704a94ca75ae079a7ceba2500701094c050f9e568a196f53e0ef35c0  explanations/data/ex_r13_dimensions.csv
60ee44d86bc281eec7760086a5e8daeeed8bb09addd5cb5e2751e8003fa36f81  explanations/data/ex_r13_divergent_counts.csv
3798cda1fd127c65e1f28a1a2642469dc8c8e6ceb70326c10f446244a5da76ce  explanations/data/ex_r13_orbits.csv
6c90420cdc8456ee62b29f52deaf8b5effff9b090e19a101699212a916090ef9  explanations/data/ex_r13_points.csv
d159f9c35283f947a298ad9bc73174c65b6726638ca51b6238e3b62b35648415  explanations/data/ex_r13_power_counting.csv
0cd9ab8b4b09768f58ef444cb11ad65eb16a065b33946369472e26f6bc9a859d  explanations/data/ex_states.csv
68f5badc7b9182d2226e6f8d36650eb324462eed9ddfa65e039a16f3b0b2978a  explanations/data/ex_superfluid_checks.txt
77cdbabb5afa005bc90d98735c371a595858e22e79c1a8d5c9689c9a4c9361af  explanations/data/ex_two_particle_counting.csv
5f5026bbbaea06c1bf85883730514138297454ee5fd8681d1f509d66e2347c0a  explanations/figures/fig_ex_bookkeeping.pdf
f25a98e31da53ef2c337460ff305c58f17756ec4e571390ba10d7bea58f340f7  explanations/figures/fig_ex_energy.pdf
e37d43b994421650e4528f6905c3250b0e3ee073213f843f923e98957b36c06a  explanations/figures/fig_ex_pairs.pdf
a7f08ad3abfac08855badb197d1c818e5b163670e840dacc76c9f1bd497cadb1  explanations/figures/fig_ex_r13_bubble.pdf
ca9deee3ab200422dbc8d0b8e88b333d6a84c5779a8a5b66382da7d71adf8222  explanations/figures/fig_ex_r13_coordinates.pdf
7cc9977ce308777e49e1725690cb8e2c5bd9bf49f9019233e8c5d3116029b871  explanations/figures/fig_ex_r13_divergences.pdf
84f38351c97118e63c15738abd0f8b6c4ac825ed7b3a8eaefd429ee200498a56  explanations/figures/fig_ex_r13_lorentz.pdf
9ebe88b30c37e04d72636b771d9045fe00bf8cadb355acc808ae6247601c7bf3  explanations/code/explain-rs/src/main.rs
ea0102d48fa91faba90a2f3df9090eb71be4d01d548fec47a38a296628a06e41  explanations/code/explain-rs/Cargo.toml
d54d506b7d88a51c7152057d89a12800edff59d1a83fc3c5ab1c5f67862a808f  explanations/code/make_explain_figures.py
d8a9fe9245e9b5b4d9eff79eca49fbd9d81ceaafd966face6a9b7ea47aa35da8  explanations/code/r13-rs/src/main.rs
563167b8f4074b5576af1fda712b23ac00695b3e1866377b0213eae834902a4e  explanations/code/r13-rs/Cargo.toml
d8a96a63b09fba62c1ef5900a690cd5993fc8138068aeee7f870fd69148c2952  explanations/code/make_r13_figures.py
\end{verbatim}

\hypertarget{the-symbols}{%
\section{14. The symbols}\label{the-symbols}}

\textbf{1 Indices, conventions and number sets}

\begin{longtable}[]{@{}
  >{\raggedright\arraybackslash}p{(\columnwidth - 6\tabcolsep) * \real{0.2200}}
  >{\raggedright\arraybackslash}p{(\columnwidth - 6\tabcolsep) * \real{0.5000}}
  >{\raggedright\arraybackslash}p{(\columnwidth - 6\tabcolsep) * \real{0.0800}}
  >{\raggedright\arraybackslash}p{(\columnwidth - 6\tabcolsep) * \real{0.2000}}@{}}
\toprule\noalign{}
\begin{minipage}[b]{\linewidth}\raggedright
symbol
\end{minipage} & \begin{minipage}[b]{\linewidth}\raggedright
meaning
\end{minipage} & \begin{minipage}[b]{\linewidth}\raggedright
origin
\end{minipage} & \begin{minipage}[b]{\linewidth}\raggedright
plain-text names
\end{minipage} \\
\midrule\noalign{}
\endhead
\bottomrule\noalign{}
\endlastfoot
\(i,j\) & index of space-like coordinates; values 1, 2, 3 (ordinary
space) unless a sum is said to run over all space-like coordinates of
the spacetime. In Revision\_SuperLoop only \(j\) is used as an index,
with the values 1, 2, 3, and with primes (\(j'\), \(j''\)) where two or
three such indices are needed: it labels the three imaginary quaternion
units \(\mathbf e_j\), which correspond to the three space directions of
\(\mathbb R^{1,3}_{\rm LL}\), the coordinates \(x^j_{\rm LL}\), the
momentum components \(q^j_{\rm LL}\), the rotation and boost generators
\(J^j_{\rm LL}\), \(K^j_{\rm LL}\), and the Pauli matrices
\(\sigma^{\rm LL}_j\) and Dirac matrices \(\alpha^j_{\rm LL}\),
\(\gamma^j_{\rm LL}\). The letter \(i\) is the imaginary unit (entry
\(i\)); the isospin index is \(\iota\) (a different letter). & taken &
\texttt{j} \texttt{j\textquotesingle{}}
\texttt{j\textquotesingle{}\textquotesingle{}} \\
\(D\) & the number of space dimensions in which the particles of matter
have momentum: 7 (all space coordinates of 7+1), 4 (fields that do not
depend on \(x5,x6,x7\)) or 3 (fields that depend on ordinary space
only). \(V_D=\pi^{D/2}/\Gamma_{\rm E}(D/2+1)\) is the volume and
\(A_D=D\,V_D\) the area of the unit ball in \(D\) dimensions
(\(\Gamma_{\rm E}\): Euler's gamma function). In Revision\_SuperLoop
\(D=3\) for the four-dimensional world of the superfluid and \(D=7\) for
the eight-dimensional manifolds (on \(4+4\) power counting sees the
total dimension 8 as well); the value 4, the coordinates \(x5,x6,x7\) of
the eight-dimensional spacetime of the project (on which such fields do
not depend), and the volume and the area of the unit ball do not occur.
& taken & \texttt{D} \\
\(i\) & the imaginary unit, \(i^2=-1\); in \(\mathbb B_{\rm LL}\) it is
the imaginary unit of the complex coefficients: it commutes and
associates with every bi-octonion and is not one of the units
\(\mathbf e_1,\dots,\mathbf e_7\) & new & \\
\(\mathbb R\), \(\mathbb C\) & the real and the complex numbers (plain
text: \texttt{R\^{}8}, \texttt{C\^{}3}, \texttt{C\^{}8} are the real or
complex vector spaces of that dimension) & new & \texttt{R} \texttt{C}
\texttt{R\^{}8} \texttt{C\^{}3} \texttt{C\^{}8} \\
\(p\), \(p'\) & a summation index over the basis of an algebra (values
\(0,\dots,7\), or \(0,\dots,3\) for the quaternions): in
Revision\_SuperLoop it labels the units \(\mathbf e_p\) of
\(\mathbb B_{\rm LL}\), the coefficients \(\xi^{\rm LL}_p(\mathbf x)\),
the coordinates \(x^p_{\rm LL}\) of \(\mathbb R^{1,7}_{\rm LL}\) and
\(\mathbb R^{1,3}_{\rm LL}\) and the Dirac matrices
\(\gamma^p_{\rm LL}\) (\(p=0,1,2,3\)); \(p'\) is a second index of the
same kind (plain text \texttt{p1}, \texttt{p2}, \texttt{p3}: three such
indices) & new & \texttt{p} \texttt{p\textquotesingle{}} \texttt{p1}
\texttt{p2} \texttt{p3} \\
\(\iota\), \(\iota'\), \(\iota''\) & the isospin index, values 1, 2, 3:
it labels the isospin generators \(T^\iota_{\rm LL}\), the Pauli
matrices \(\tau^{\rm LL}_\iota\) on the isospin index of a field, the
pair channels \(\mathsf K^\iota_{\rm LL}\),
\(\mathsf K'^\iota_{\rm LL}\), the pair operators \(\Pi^\iota_{\rm LL}\)
and the components \(d^\iota_{\rm LL}\) of the order parameter; it is an
internal direction, not a direction in space (that is \(j\)) & new &
\texttt{iota} \\
\(\epsilon^{\rm LL}_{jj'j''}\) & the Levi-Civita symbol: \(+1\) if
\((j,j',j'')\) is an even permutation of \((1,2,3)\), \(-1\) if it is an
odd one, \(0\) otherwise; also with the isospin indices,
\(\epsilon^{\rm LL}_{\iota\iota'\iota''}\) & new & \texttt{eps3} \\
\(\sigma^{\rm LL}_0\), \(\sigma^{\rm LL}_j\) & the \(2\times2\) unit
matrix \(\sigma^{\rm LL}_0\) and the Pauli matrices
\(\sigma^{\rm LL}_1=\begin{pmatrix}0&1\\1&0\end{pmatrix}\),
\(\sigma^{\rm LL}_2=\begin{pmatrix}0&-i\\i&0\end{pmatrix}\),
\(\sigma^{\rm LL}_3=\begin{pmatrix}1&0\\0&-1\end{pmatrix}\), acting on
the two components of a Weyl spinor or on the species index & new &
\texttt{sigma\_0} \texttt{sigma\_j} \texttt{sigma\_1} \texttt{sigma\_2}
\texttt{sigma\_3} \\
\(\tau^{\rm LL}_\iota\) & the same three Pauli matrices acting on the
isospin index of a field (the index of the two components up and down of
an isodoublet) & new & \texttt{tau\_iota} \texttt{tau\_1}
\texttt{tau\_2} \texttt{tau\_3} \\
\(\mathsf E_{\rm LL}\) & the real antisymmetric \(2\times2\) matrix
\(i\sigma^{\rm LL}_2=\begin{pmatrix}0&1\\-1&0\end{pmatrix}\), with
\(\mathsf E_{\rm LL}^{-1}=\mathsf E_{\rm LL}^T=-\mathsf E_{\rm LL}\); it
acts on a Weyl, isospin or species index, as stated at each place & new
& \texttt{eps2} \\
\(\mathsf 1_4\) & the \(4\times4\) unit matrix & new & \\
\end{longtable}

\textbf{2 The bi-octonions and the super-loop}

\begin{longtable}[]{@{}
  >{\raggedright\arraybackslash}p{(\columnwidth - 6\tabcolsep) * \real{0.2200}}
  >{\raggedright\arraybackslash}p{(\columnwidth - 6\tabcolsep) * \real{0.5000}}
  >{\raggedright\arraybackslash}p{(\columnwidth - 6\tabcolsep) * \real{0.0800}}
  >{\raggedright\arraybackslash}p{(\columnwidth - 6\tabcolsep) * \real{0.2000}}@{}}
\toprule\noalign{}
\begin{minipage}[b]{\linewidth}\raggedright
symbol
\end{minipage} & \begin{minipage}[b]{\linewidth}\raggedright
meaning
\end{minipage} & \begin{minipage}[b]{\linewidth}\raggedright
origin
\end{minipage} & \begin{minipage}[b]{\linewidth}\raggedright
plain-text names
\end{minipage} \\
\midrule\noalign{}
\endhead
\bottomrule\noalign{}
\endlastfoot
\(\mathbb H\), \(\mathbb O\) & the quaternions and the octonions: the
real combinations of \(\mathbf e_0,\dots,\mathbf e_3\) and of
\(\mathbf e_0,\dots,\mathbf e_7\); \({\rm Im}\,\mathbb H\) and
\({\rm Im}\,\mathbb O\) are their imaginary parts, the real combinations
of \(\mathbf e_1,\mathbf e_2,\mathbf e_3\) and of
\(\mathbf e_1,\dots,\mathbf e_7\) (plain text \texttt{Im\ H},
\texttt{Im\ O}; \texttt{C\ (x)\ H} is \(\mathbb C\otimes\mathbb H\), the
bi-quaternions) & new & \texttt{H} \texttt{O} \\
\(\mathbb B_{\rm LL}\) & the bi-octonions \(\mathbb C\otimes\mathbb O\):
the complex combinations
\(\mathbf x=\sum_{p=0}^7\xi^{\rm LL}_p(\mathbf x)\,\mathbf e_p\) of the
eight units, with the product of the units extended bilinearly; real
dimension 16 & new & \texttt{B\_LL} \\
\(\mathbf e_0,\dots,\mathbf e_7\); \(\mathbf e_p\), \(\mathbf e_j\),
\(\mathbf e_\iota\) & the units of \(\mathbb B_{\rm LL}\), of
\(\mathbb O\) and of \(\mathbb H\): \(\mathbf e_0=1\);
\(\mathbf e_1,\mathbf e_2,\mathbf e_3\) the quaternion units,
\(\mathbf e_j^2=-1\),
\(\mathbf e_1\mathbf e_2=-\mathbf e_2\mathbf e_1=\mathbf e_3\) and
cyclically; \(\mathbf e_4\) the unit of the Cayley-Dickson doubling,
\(\mathbf e_4^2=-1\); \(\mathbf e_5=\mathbf e_1\mathbf e_4\),
\(\mathbf e_6=\mathbf e_2\mathbf e_4\),
\(\mathbf e_7=\mathbf e_3\mathbf e_4\) (\(\mathbf e_\iota\): a
quaternion unit labelled by the isospin index; plain text \texttt{ie1},
\ldots, \texttt{ie7}: the bi-octonions
\(i\mathbf e_1,\dots,i\mathbf e_7\)) & new & \texttt{e0} \texttt{e1}
\texttt{e2} \texttt{e3} \texttt{e4} \texttt{e5} \texttt{e6} \texttt{e7}
\texttt{e\_p} \texttt{e\_p\textquotesingle{}} \texttt{e\_j}
\texttt{e\_iota} \texttt{e\_p1} \texttt{e\_p2} \texttt{e\_p3}
\texttt{ie1} \texttt{ie2} \texttt{ie3} \texttt{ie4} \texttt{ie5}
\texttt{ie6} \texttt{ie7} \\
\(\xi^{\rm LL}_p(\mathbf x)\) & the complex coefficient of the unit
\(\mathbf e_p\) in the bi-octonion \(\mathbf x\) & new & \texttt{xi\_p}
\texttt{xi\_0} \texttt{xi\_1} \texttt{xi\_2} \texttt{xi\_3} \\
\(\mathbf x\), \(\mathbf y\), \(\mathbf z\), \(\mathbf w\) & elements of
\(\mathbb B_{\rm LL}\) (bi-octonions); \(\mathbf w\) is also a point of
the super-loop manifold \(\mathbb L_{\rm LL}\) & new & \texttt{x}
\texttt{y} \texttt{z} \texttt{w} \\
\(\mathbf a\), \(\mathbf b\), \(\mathbf c\), \(\mathbf d\);
\(\mathbf b'\) & bi-quaternions: elements of
\(\mathbb C\otimes\mathbb H\), the complex combinations of
\(\mathbf e_0,\mathbf e_1,\mathbf e_2,\mathbf e_3\); every bi-octonion
is written \(\mathbf x=\mathbf a+\mathbf b\,\mathbf e_4\)
(Cayley-Dickson form), and \(\mathbf b\,\mathbf e_4\),
\(\mathbf b'\,\mathbf e_4\) are spinor charges & new & \texttt{aH}
\texttt{bS} \texttt{cH} \texttt{dS} \texttt{bS\textquotesingle{}} \\
\(\bar{\mathbf x}\), \(\mathbf x^*\), \(\mathbf x^\dagger\) & the three
conjugations (also of bi-quaternions): \(\bar{\mathbf x}\) changes the
signs of the coefficients of \(\mathbf e_1,\dots,\mathbf e_7\) (it is
complex-linear,
\(\overline{\mathbf x\mathbf y}=\bar{\mathbf y}\,\bar{\mathbf x}\), and
\(\overline{\mathbf a+\mathbf b\mathbf e_4}=\bar{\mathbf a}-\mathbf b\mathbf e_4\));
\(\mathbf x^*\) replaces every complex coefficient by its complex
conjugate; \(\mathbf x^\dagger=(\bar{\mathbf x})^*\) (plain text:
\texttt{xbar}, \texttt{x\^{}*}, \texttt{x\^{}dag}, and likewise
\texttt{aHbar}, \texttt{gL\^{}dag}, \ldots) & new & \texttt{xbar}
\texttt{x\^{}*} \texttt{x\^{}dag} \texttt{aHbar} \texttt{bSbar}
\texttt{bS\textquotesingle{}bar} \texttt{cHbar} \texttt{dSbar}
\texttt{gLbar} \texttt{hLbar} \texttt{rPbar}
\texttt{rP\textquotesingle{}bar} \texttt{bS\^{}*}
\texttt{bS\textquotesingle{}\^{}*} \texttt{gL\^{}dag}
\texttt{uL\^{}dag} \\
\(N_{\rm LL}(\mathbf x)\), \(N_{\rm LL}(\mathbf x,\mathbf y)\) & the
complex norm
\(N_{\rm LL}(\mathbf x)=\sum_p\xi^{\rm LL}_p(\mathbf x)^2\), with
\(\mathbf x\bar{\mathbf x}=\bar{\mathbf x}\mathbf x=N_{\rm LL}(\mathbf x)\,\mathbf e_0\),
and its polarization
\(N_{\rm LL}(\mathbf x,\mathbf y)=\sum_p\xi^{\rm LL}_p(\mathbf x)\,\xi^{\rm LL}_p(\mathbf y)\),
\(N_{\rm LL}(\mathbf x,\mathbf x)=N_{\rm LL}(\mathbf x)\)
(complex-valued; for bi-quaternions the determinant of the \(2\times2\)
matrix \(M_{\rm LL}\)) & new & \texttt{N\_LL} \\
\([\mathbf x,\mathbf y]\); \([\mathbf x,\mathbf y,\mathbf z]\) & the
commutator \(\mathbf x\mathbf y-\mathbf y\mathbf x\) and the associator
\((\mathbf x\mathbf y)\mathbf z-\mathbf x(\mathbf y\mathbf z)\) of
bi-octonions (the same brackets of other elements are defined at their
place: the bracket of an algebra \(\mathcal F^{\rm A}_{\rm LL}\),
\(\mathcal F^{\rm B}_{\rm LL}\), \(\mathcal F^{\rm C}_{\rm LL}\), and
the commutator of two matrices or operators, the product in one order
minus the product in the other) & new & \texttt{{[}x,y{]}}
\texttt{{[}x,y,z{]}} \\
\(\mathrm{Jac}_{\rm LL}\) & the Jacobiator
\(\mathrm{Jac}_{\rm LL}(\mathbf x,\mathbf y,\mathbf z)=[[\mathbf x,\mathbf y],\mathbf z]+[[\mathbf y,\mathbf z],\mathbf x]+[[\mathbf z,\mathbf x],\mathbf y]\)
of a bracket; it vanishes identically exactly in a Lie algebra & new &
\texttt{Jac} \\
\(L^{\rm LL}_{\mathbf x}\), \(R^{\rm LL}_{\mathbf x}\) & the left and
right multiplications \(\mathbf y\mapsto\mathbf x\mathbf y\) and
\(\mathbf y\mapsto\mathbf y\mathbf x\) by a bi-octonion, real
\(16\times16\) matrices (on \(\mathbb O\) or
\(\mathbb O^{\rm s}_{\rm LL}\): \(8\times8\)) & new & \texttt{L\_x}
\texttt{L\_y} \texttt{L\_e1} \texttt{L\_e4} \texttt{L\_{[}e1,e4{]}}
\texttt{L\_(e\_p)} \texttt{L\_(e\_p\textquotesingle{})} \texttt{R\_x}
\texttt{R\_y} \\
\(\mathbb L_{\rm LL}\) & the Lorentz super-loop: the unit bi-octonions,
\(\mathbb L_{\rm LL}=\{\mathbf w\in\mathbb B_{\rm LL}:N_{\rm LL}(\mathbf w)=1\}\),
with the product of \(\mathbb B_{\rm LL}\); a Moufang loop (not a
group), and a manifold of real dimension 14 & new & \texttt{L\_LL} \\
\(\mathbf g\), \(\mathbf h\) & elements of \(SL(2,\mathbb C)\),
i.e.~unit bi-quaternions (\(N_{\rm LL}(\mathbf g)=1\)), a subgroup of
\(\mathbb L_{\rm LL}\) & new & \texttt{gL} \texttt{hL} \\
\(\mathbf v\) & a unit (bi-)quaternion that rotates the isospin of the
spinor charges (from the left) & new & \texttt{vI} \\
\(\mathfrak m_{\rm LL}\), \(\mathfrak l_{\rm LL}\),
\(\mathfrak s_{\rm LL}\) & the Malcev algebra of the super-loop,
\(\mathfrak m_{\rm LL}=\mathbb C\otimes{\rm Im}\,\mathbb O\) with the
commutator (real dimension 14), its Lorentz part
\(\mathfrak l_{\rm LL}=\mathbb C\otimes{\rm Im}\,\mathbb H\) (6, a Lie
algebra isomorphic to \(\mathfrak{so}(3,1)\)), and the space of spinor
charges
\(\mathfrak s_{\rm LL}=(\mathbb C\otimes\mathbb H)\,\mathbf e_4\) (8);
\(\mathfrak m_{\rm LL}=\mathfrak l_{\rm LL}\oplus\mathfrak s_{\rm LL}\)
& new & \texttt{m\_LL} \texttt{l\_LL} \texttt{s\_LL} \\
\(J^j_{\rm LL}\), \(K^j_{\rm LL}\) & the rotation and boost generators
of \(\mathfrak l_{\rm LL}\), \(J^j_{\rm LL}=\mathbf e_j/2\) and
\(K^j_{\rm LL}=i\,\mathbf e_j/2\) & new & \texttt{J\_LL\^{}j}
\texttt{K\_LL\^{}j} \texttt{J\_LL\^{}j\textquotesingle{}}
\texttt{K\_LL\^{}j\textquotesingle{}}
\texttt{J\_LL\^{}j\textquotesingle{}\textquotesingle{}}
\texttt{K\_LL\^{}j\textquotesingle{}\textquotesingle{}}
\texttt{J\_LL\^{}3} \\
\(\mathbf u\), \(\mathbf u'\), \(\mathbf u_1\), \(\mathbf u_2\) &
elements of the Lorentz part \(\mathfrak l_{\rm LL}\) (Lorentz
generators): imaginary bi-quaternions & new & \texttt{uL}
\texttt{uL\textquotesingle{}} \texttt{uL1} \texttt{uL2} \\
\(\mathbf t\), \(\mathbf t'\) & elements of \({\rm Im}\,\mathbb H\)
(real imaginary quaternions): isospin generators, which act on the
spinor charges from the left & new & \texttt{tI}
\texttt{tI\textquotesingle{}} \\
\(W_{\rm LL}(\mathbf u)\) & the Weyl representation of
\(\mathfrak l_{\rm LL}\) on bi-quaternions,
\(W_{\rm LL}(\mathbf u)\mathbf b=-\mathbf b\,\mathbf u\): a Lie
representation,
\([W_{\rm LL}(\mathbf u_1),W_{\rm LL}(\mathbf u_2)]=W_{\rm LL}([\mathbf u_1,\mathbf u_2])\)
& new & \texttt{W\_LL} \\
\(\rho_{\rm LL}(\mathbf u)\) & the inner action of a Lorentz element on
the spinor charges,
\([\mathbf u,\mathbf b\,\mathbf e_4]=(\rho_{\rm LL}(\mathbf u)\mathbf b)\,\mathbf e_4\)
with \(\rho_{\rm LL}=-2W_{\rm LL}\); not a representation & new &
\texttt{rho\_LL} \\
\(\Theta_{\rm LL}\) & the twist charge: the linear map of
\(\mathfrak m_{\rm LL}\) that is \(1\) on \(\mathfrak l_{\rm LL}\) and
\(-2\) on \(\mathfrak s_{\rm LL}\); the inner action is
\([\mathbf u,\cdot\,]=D_{\rm LL}(\mathbf u)\circ\Theta_{\rm LL}\) (as a
number: the eigenvalue \(1\) or \(-2\), or the value of
\(c^{\rm tw}_{\rm LL}\) in the families) & new & \texttt{Theta\_LL}
\texttt{Theta\_LL\^{}2} \\
\(D_{\rm LL}(\mathbf u)\), \(D^{\rm v}_{\rm LL}(\mathbf t)\) & the
Lorentz and isospin derivations of \(\mathbb B_{\rm LL}\):
\(D_{\rm LL}(\mathbf u)(\mathbf a+\mathbf b\mathbf e_4)=[\mathbf u,\mathbf a]+(-\mathbf b\mathbf u)\mathbf e_4\)
and
\(D^{\rm v}_{\rm LL}(\mathbf t)(\mathbf a+\mathbf b\mathbf e_4)=(\mathbf t\mathbf b)\mathbf e_4\)
(both linear in their argument) & new & \texttt{D\_LL}
\texttt{Dv\_LL} \\
\(\mathcal D_{\rm LL}\) & a derivation of \(\mathbb B_{\rm LL}\): a
linear map with
\(\mathcal D_{\rm LL}(\mathbf x\mathbf y)=(\mathcal D_{\rm LL}\mathbf x)\mathbf y+\mathbf x(\mathcal D_{\rm LL}\mathbf y)\),
the infinitesimal form of an automorphism & new & \\
\(\varphi^{\rm LL}_{\mathbf g,\mathbf v}\) & the automorphism
\(\mathbf a+\mathbf b\mathbf e_4\mapsto\mathbf g\mathbf a\bar{\mathbf g}+(\mathbf v\mathbf b\bar{\mathbf g})\mathbf e_4\)
of \(\mathbb B_{\rm LL}\) (\(\mathbf g\) a Lorentz transformation,
\(\mathbf v\) an isospin rotation) & new & \texttt{varphi} \\
\(\varrho_{\rm LL}\); \(\mathsf O_1\), \(\mathsf O_2\), \(\mathsf O_3\)
& a linear map from \(\mathfrak m_{\rm LL}\) to the linear operators
(matrices) on a vector space that turns brackets into commutators,
\(\varrho_{\rm LL}([\mathbf x,\mathbf y])=\varrho_{\rm LL}(\mathbf x)\varrho_{\rm LL}(\mathbf y)-\varrho_{\rm LL}(\mathbf y)\varrho_{\rm LL}(\mathbf x)\)
(Theorem R); \(\mathsf O_1,\mathsf O_2,\mathsf O_3\) three such
operators & new & \\
\(\mathrm{Kil}_{\rm LL}\); \(\mathrm{ad}\) & the Killing form
\(\mathrm{Kil}_{\rm LL}(\mathbf x,\mathbf y)=\mathrm{tr}(\mathrm{ad}\,\mathbf x\,\mathrm{ad}\,\mathbf y)\)
of a Lie algebra, where \(\mathrm{ad}\,\mathbf x\) is the matrix of the
map \(\mathbf y\mapsto[\mathbf x,\mathbf y]\) & new & \texttt{Killing}
\texttt{ad} \\
\end{longtable}

\textbf{3 Operators on the super-loop and their eigenvalues}

\begin{longtable}[]{@{}
  >{\raggedright\arraybackslash}p{(\columnwidth - 6\tabcolsep) * \real{0.2200}}
  >{\raggedright\arraybackslash}p{(\columnwidth - 6\tabcolsep) * \real{0.5000}}
  >{\raggedright\arraybackslash}p{(\columnwidth - 6\tabcolsep) * \real{0.0800}}
  >{\raggedright\arraybackslash}p{(\columnwidth - 6\tabcolsep) * \real{0.2000}}@{}}
\toprule\noalign{}
\begin{minipage}[b]{\linewidth}\raggedright
symbol
\end{minipage} & \begin{minipage}[b]{\linewidth}\raggedright
meaning
\end{minipage} & \begin{minipage}[b]{\linewidth}\raggedright
origin
\end{minipage} & \begin{minipage}[b]{\linewidth}\raggedright
plain-text names
\end{minipage} \\
\midrule\noalign{}
\endhead
\bottomrule\noalign{}
\endlastfoot
\(Q_{\rm LL}\) & the electric charge, \(Q_{\rm LL}=T^3_{\rm LL}\):
eigenvalues \(0\) on \(\mathbb C\otimes\mathbb H\) and \(\pm\tfrac12\)
on the spinor charges & new & \texttt{Q\_LL} \\
\(T^\iota_{\rm LL}\), \((T_{\rm LL})^2\) & the isospin generators
\(T^\iota_{\rm LL}=i\,D^{\rm v}_{\rm LL}(\mathbf e_\iota/2)\),
Hermitian, \([T^1_{\rm LL},T^2_{\rm LL}]=iT^3_{\rm LL}\) and cyclically,
and the isospin Casimir
\((T_{\rm LL})^2=\sum_\iota(T^\iota_{\rm LL})^2\) & new &
\texttt{T\_LL\^{}iota} \texttt{T\_LL\^{}3} \texttt{T\_LL\^{}2} \\
\(\mathrm{Cas}_{\rm LL}\) & the Casimir operator of the rotations,
\(\mathrm{Cas}_{\rm LL}=-\sum_jD_{\rm LL}(J^j_{\rm LL})^2\); its
eigenvalue is 0 for spin 0, 3/4 for spin 1/2 and 2 for spin 1 & new &
\texttt{Cas\_LL} \\
\end{longtable}

\textbf{4 The covariant families and their classification}

\begin{longtable}[]{@{}
  >{\raggedright\arraybackslash}p{(\columnwidth - 6\tabcolsep) * \real{0.2200}}
  >{\raggedright\arraybackslash}p{(\columnwidth - 6\tabcolsep) * \real{0.5000}}
  >{\raggedright\arraybackslash}p{(\columnwidth - 6\tabcolsep) * \real{0.0800}}
  >{\raggedright\arraybackslash}p{(\columnwidth - 6\tabcolsep) * \real{0.2000}}@{}}
\toprule\noalign{}
\begin{minipage}[b]{\linewidth}\raggedright
symbol
\end{minipage} & \begin{minipage}[b]{\linewidth}\raggedright
meaning
\end{minipage} & \begin{minipage}[b]{\linewidth}\raggedright
origin
\end{minipage} & \begin{minipage}[b]{\linewidth}\raggedright
plain-text names
\end{minipage} \\
\midrule\noalign{}
\endhead
\bottomrule\noalign{}
\endlastfoot
\(\mathcal F^{\rm A}_{\rm LL}\), \(\mathcal F^{\rm B}_{\rm LL}\),
\(\mathcal F^{\rm C}_{\rm LL}\) & the three families of covariant
brackets: \(\mathcal F^{\rm A}_{\rm LL}\) on
\(\mathfrak l_{\rm LL}\oplus\mathfrak s_{\rm LL}\) (dimension 14),
\(\mathcal F^{\rm B}_{\rm LL}\) on
\(\mathfrak l_{\rm LL}\oplus\mathfrak p_{\rm LL}\oplus\mathfrak s_{\rm LL}\)
(18), \(\mathcal F^{\rm C}_{\rm LL}\) on
\(\mathfrak l_{\rm LL}\oplus\mathfrak{su}(2)\oplus\mathfrak s_{\rm LL}\)
(17) & new & \texttt{F\_A} \texttt{F\_B} \texttt{F\_C} \\
\(\mathfrak p_{\rm LL}\); \(\mathbf r\), \(\mathbf r'\) & the
translations of the family \(\mathcal F^{\rm B}_{\rm LL}\): the real
space of the Hermitian bi-quaternions \(\mathbf r=\mathbf r^\dagger\)
(dimension 4), on which \(\mathfrak l_{\rm LL}\) acts by
\([\mathbf u,\mathbf r]=\mathbf u\mathbf r+\mathbf r\mathbf u^\dagger\)
& new & \texttt{p\_LL} \texttt{rP} \texttt{rP\textquotesingle{}} \\
\(c^{\rm tw}_{\rm LL}\), \(c^{\rm sl}_{\rm LL}\),
\(c^{\rm PP}_{\rm LL}\), \(c^{\rm Ps}_{\rm LL}\),
\(c^{\rm sP}_{\rm LL}\), \(c^{\rm TT}_{\rm LL}\),
\(c^{\rm Ts}_{\rm LL}\), \(c^{\rm sT}_{\rm LL}\) & the coefficients of
the families:
\([\mathbf u,\mathbf b\mathbf e_4]=c^{\rm tw}_{\rm LL}(W_{\rm LL}(\mathbf u)\mathbf b)\mathbf e_4\)
(the twist),
\([\mathbf b\mathbf e_4,\mathbf b'\mathbf e_4]\supset c^{\rm sl}_{\rm LL}(\bar{\mathbf b}\mathbf b'-\bar{\mathbf b}'\mathbf b)\),
\([\mathbf r,\mathbf r']=c^{\rm PP}_{\rm LL}(\mathbf r\bar{\mathbf r}'-\mathbf r'\bar{\mathbf r})\),
\([\mathbf r,\mathbf b\mathbf e_4]=c^{\rm Ps}_{\rm LL}(\mathbf b^*\bar{\mathbf r})\mathbf e_4\),
\([\mathbf b\mathbf e_4,\mathbf b'\mathbf e_4]\supset c^{\rm sP}_{\rm LL}\,i(\bar{\mathbf b}\mathbf b'^*-\bar{\mathbf b}'\mathbf b^*)\),
\([\mathbf t,\mathbf t']=c^{\rm TT}_{\rm LL}(\mathbf t\mathbf t'-\mathbf t'\mathbf t)\),
\([\mathbf t,\mathbf b\mathbf e_4]=c^{\rm Ts}_{\rm LL}(\mathbf t\mathbf b)\mathbf e_4\),
\([\mathbf b\mathbf e_4,\mathbf b'\mathbf e_4]\supset\mathrm{Re}[c^{\rm sT}_{\rm LL}(\mathbf b\bar{\mathbf b}'-\mathbf b'\bar{\mathbf b})]\);
complex except \(c^{\rm sP}_{\rm LL}\), \(c^{\rm TT}_{\rm LL}\) (real);
plain text \texttt{ctw\_re}, \texttt{ctw\_im}: real and imaginary part
of \(c^{\rm tw}_{\rm LL}\), and likewise & new & \texttt{c\_tw}
\texttt{c\_sl} \texttt{c\_PP} \texttt{c\_Ps} \texttt{c\_sP}
\texttt{c\_TT} \texttt{c\_Ts} \texttt{c\_sT} \texttt{csP} \texttt{cTT}
\texttt{ctw\_re} \texttt{ctw\_im} \texttt{csl\_re} \texttt{csl\_im}
\texttt{cPP\_re} \texttt{cPP\_im} \texttt{cPs\_re} \texttt{cPs\_im}
\texttt{cTs\_re} \texttt{cTs\_im} \texttt{csT\_re} \texttt{csT\_im} \\
\(\lambda^{\rm sc}_{\rm LL}\) & a complex number \(\neq0\) by which the
spinor charges are rescaled,
\(\mathbf b\mapsto\lambda^{\rm sc}_{\rm LL}\mathbf b\) & new &
\texttt{lambda\_sc} \\
\(\mathcal I^{\rm Jac}_{\rm LL}\), \(\mathcal I^{\rm Mal}_{\rm LL}\);
\(w^{\rm sat}_{\rm LL}\) & the ideals of polynomials in the real
coefficients of a family generated by the Jacobi and by the Malcev
identity; \(w^{\rm sat}_{\rm LL}\) the auxiliary variable of the
saturation (the points of the variety at which a given polynomial does
not vanish are the points of the variety of the ideal enlarged by that
polynomial times \(w^{\rm sat}_{\rm LL}\), minus 1) & new &
\texttt{w\_sat} \\
\end{longtable}

\textbf{5 The manifolds}

\begin{longtable}[]{@{}
  >{\raggedright\arraybackslash}p{(\columnwidth - 6\tabcolsep) * \real{0.2200}}
  >{\raggedright\arraybackslash}p{(\columnwidth - 6\tabcolsep) * \real{0.5000}}
  >{\raggedright\arraybackslash}p{(\columnwidth - 6\tabcolsep) * \real{0.0800}}
  >{\raggedright\arraybackslash}p{(\columnwidth - 6\tabcolsep) * \real{0.2000}}@{}}
\toprule\noalign{}
\begin{minipage}[b]{\linewidth}\raggedright
symbol
\end{minipage} & \begin{minipage}[b]{\linewidth}\raggedright
meaning
\end{minipage} & \begin{minipage}[b]{\linewidth}\raggedright
origin
\end{minipage} & \begin{minipage}[b]{\linewidth}\raggedright
plain-text names
\end{minipage} \\
\midrule\noalign{}
\endhead
\bottomrule\noalign{}
\endlastfoot
\(\mathbb R^{1,7}_{\rm LL}\), \(\mathbb R^{1,3}_{\rm LL}\);
\(x^p_{\rm LL}\) & the Hermitian bi-octonions
\(\mathbf x=\mathbf x^\dagger=x^0_{\rm LL}\mathbf e_0+i\sum_{p=1}^7x^p_{\rm LL}\mathbf e_p\)
(real \(x^p_{\rm LL}\)): the \(7+1\) Minkowski space with the real norm
\(N_{\rm LL}(\mathbf x)=(x^0_{\rm LL})^2-\sum_p(x^p_{\rm LL})^2\); and
its subspace of Hermitian bi-quaternions (\(p\le3\)), Minkowski space of
\(3+1\) dimensions & new & \texttt{R\^{}(1,7)\_LL}
\texttt{R\^{}(1,3)\_LL} \texttt{x\_LL\^{}0} \texttt{x\_LL\^{}p}
\texttt{x\_LL\^{}1} \texttt{x\_LL\^{}2} \texttt{x\_LL\^{}3}
\texttt{x\_LL\^{}j} \\
\(\mathbb O^{\rm s}_{\rm LL}\) & the split octonions inside
\(\mathbb B_{\rm LL}\): the real space
\(\mathbb H+(i\mathbb H)\mathbf e_4\) (real combinations of
\(\mathbf e_0,\dots,\mathbf e_3,i\mathbf e_4,\dots,i\mathbf e_7\)),
closed under the product, on which \(N_{\rm LL}\) has signature
\((4,4)\): the \(4+4\) manifold (plain text \texttt{R\^{}(4,4)}) & new &
\texttt{Os} \texttt{R\^{}(4,4)} \\
\(S^7_{\rm LL}\), \(S^{3,4}_{\rm LL}\) & the unit elements
(\(N_{\rm LL}=1\)) of \(\mathbb O\) (the seven-sphere) and of
\(\mathbb O^{\rm s}_{\rm LL}\) (the pseudo-sphere): the two real forms
of \(\mathbb L_{\rm LL}\), both Moufang loops & new & \texttt{S\^{}7}
\texttt{S\^{}(3,4)} \\
\(X^{\rm LL}_{\mathbf x}\) & the frame field
\(X^{\rm LL}_{\mathbf x}(\mathbf w)=\mathbf w\mathbf x\) on
\(\mathbb L_{\rm LL}\) for imaginary \(\mathbf x\) (a vector field: the
velocity of the curve \(\mathbf w(\theta_{\rm LL})\) with
\(d\mathbf w/d\theta_{\rm LL}=\mathbf w\mathbf x\)) & new &
\texttt{X\_x} \texttt{X\_y} \\
\(\mathsf L_{\rm LL}(\mathbf g)\) & the real \(4\times4\) matrix of the
Lorentz transformation
\(\mathbf x\mapsto\mathbf g\mathbf x\mathbf g^\dagger\) of
\(\mathbb R^{1,3}_{\rm LL}\) in the coordinates
\(x^0_{\rm LL},\dots,x^3_{\rm LL}\), for a unit bi-quaternion
\(\mathbf g\): its entry in row \(p\) and column \(p'\)
(\(p,p'=0,\dots,3\)) is
\(\tfrac12\,\mathrm{Re}\,\mathrm{tr}(\sigma^{\rm LL}_pM_{\rm LL}(\mathbf g)\,\sigma^{\rm LL}_{p'}M_{\rm LL}(\mathbf g)^\dagger)\),
with \(\sigma^{\rm LL}_0\) the unit matrix & new & \texttt{Lmat} \\
\(\chi_{\rm LL}\) & a two-component spinor (a column of two complex
numbers) on which a unit bi-quaternion \(\mathbf g\) acts by
\(\chi_{\rm LL}\mapsto M_{\rm LL}(\mathbf g)\chi_{\rm LL}\); the
Hermitian matrix \(\chi_{\rm LL}\chi_{\rm LL}^\dagger\) is the matrix
\(M_{\rm LL}(\mathbf x)\) of a point \(\mathbf x\) of the future light
cone of \(\mathbb R^{1,3}_{\rm LL}\) & new & \texttt{chi} \\
\(\theta_{\rm LL}\) & the real parameter of a one-parameter family of
transformations or of a curve (an angle, a rapidity) & new &
\texttt{theta} \\
\(\mathcal V_{\rm LL}(\mathbf u)\) & the vector action of a Lorentz
element on \(\mathbb R^{1,3}_{\rm LL}\),
\(\mathcal V_{\rm LL}(\mathbf u)\mathbf x=\mathbf u\mathbf x+\mathbf x\mathbf u^\dagger\),
a real \(4\times4\) matrix with entries
\(\mathcal V_{\rm LL}(\mathbf u)^p{}_{p'}\) in the coordinates
\(x^p_{\rm LL}\) (plain text \texttt{V\_LL\^{}p\_p\textquotesingle{}}) &
new & \texttt{V\_LL} \\
\end{longtable}

\textbf{6 The spin-zero superfluid}

\begin{longtable}[]{@{}
  >{\raggedright\arraybackslash}p{(\columnwidth - 6\tabcolsep) * \real{0.2200}}
  >{\raggedright\arraybackslash}p{(\columnwidth - 6\tabcolsep) * \real{0.5000}}
  >{\raggedright\arraybackslash}p{(\columnwidth - 6\tabcolsep) * \real{0.0800}}
  >{\raggedright\arraybackslash}p{(\columnwidth - 6\tabcolsep) * \real{0.2000}}@{}}
\toprule\noalign{}
\begin{minipage}[b]{\linewidth}\raggedright
symbol
\end{minipage} & \begin{minipage}[b]{\linewidth}\raggedright
meaning
\end{minipage} & \begin{minipage}[b]{\linewidth}\raggedright
origin
\end{minipage} & \begin{minipage}[b]{\linewidth}\raggedright
plain-text names
\end{minipage} \\
\midrule\noalign{}
\endhead
\bottomrule\noalign{}
\endlastfoot
\(M_{\rm LL}(\mathbf b)\) & the complex \(2\times2\) matrix of a
bi-quaternion, \(M_{\rm LL}(\mathbf e_0)=\sigma^{\rm LL}_0\),
\(M_{\rm LL}(\mathbf e_j)=-i\sigma^{\rm LL}_j\), extended
complex-linearly;
\(M_{\rm LL}(\mathbf b\mathbf b')=M_{\rm LL}(\mathbf b)M_{\rm LL}(\mathbf b')\),
\(\det M_{\rm LL}(\mathbf b)=N_{\rm LL}(\mathbf b)\) (plain text
\texttt{M\_2(C)}: all complex \(2\times2\) matrices) & new &
\texttt{M\_LL} \texttt{M\_2(C)} \\
\(\psi'_{\rm LL}\) & the Weyl isodoublet of a spinor charge
\(\mathbf b\mathbf e_4\): the \(2\times2\) matrix
\(\psi'_{\rm LL}=\mathsf E_{\rm LL}^{-1}M_{\rm LL}(\mathbf b)^T\), whose
first column (isospin up) and second column (isospin down) are Weyl
spinors & new & \texttt{psi\textquotesingle{}\_LL} \\
\(\Psi_{\rm LL}\) & the Dirac isodoublet of the spinor charges: a field
of 8 complex Grassmann components (Dirac \(\otimes\) isospin), whose
\(\gamma^{\rm ch}_{\rm LL}=+1\) half is built from
\(\mathfrak s_{\rm LL}\) and whose other half from the complex-conjugate
super-loop & new & \texttt{Psi\_LL} \texttt{Psi\_LL\^{}T}
\texttt{Psi\_LL\^{}dag} \\
\(\gamma^p_{\rm LL}\), \(\gamma^{\rm ch}_{\rm LL}\), \(C_{\rm LL}\),
\(\alpha^j_{\rm LL}\), \(\beta_{\rm LL}\) & the \(4\times4\) Dirac
matrices of \(\mathbb R^{1,3}_{\rm LL}\) in the chiral basis,
\(\gamma^0_{\rm LL}=\begin{pmatrix}0&\sigma^{\rm LL}_0\\\sigma^{\rm LL}_0&0\end{pmatrix}\),
\(\gamma^j_{\rm LL}=\begin{pmatrix}0&\sigma^{\rm LL}_j\\-\sigma^{\rm LL}_j&0\end{pmatrix}\);
the chirality matrix
\(\gamma^{\rm ch}_{\rm LL}=i\gamma^0_{\rm LL}\gamma^1_{\rm LL}\gamma^2_{\rm LL}\gamma^3_{\rm LL}=\mathrm{diag}(-1,-1,1,1)\);
the charge-conjugation matrix
\(C_{\rm LL}=i\gamma^2_{\rm LL}\gamma^0_{\rm LL}\);
\(\alpha^j_{\rm LL}=\gamma^0_{\rm LL}\gamma^j_{\rm LL}\),
\(\beta_{\rm LL}=\gamma^0_{\rm LL}\) (on Dirac \(\otimes\) isospin each
acts as \(\otimes\sigma^{\rm LL}_0\)) & new & \texttt{gamma\_LL\^{}p}
\texttt{gamma\_LL\^{}p\textquotesingle{}} \texttt{gamma\_ch}
\texttt{C\_LL} \texttt{alpha\_LL\^{}j} \texttt{beta\_LL} \\
\(\mathsf K^\iota_{\rm LL}\), \(\mathsf K'^\iota_{\rm LL}\),
\(\mathsf K^{\rm s}_{\rm LL}\) & the \(8\times8\) pair matrices of the
spin-zero channels: scalar
\(\mathsf K^\iota_{\rm LL}=C_{\rm LL}\gamma^{\rm ch}_{\rm LL}\otimes\mathsf E_{\rm LL}\tau^{\rm LL}_\iota\)
(\(0^+\)), pseudoscalar
\(\mathsf K'^\iota_{\rm LL}=C_{\rm LL}\otimes\mathsf E_{\rm LL}\tau^{\rm LL}_\iota\)
(\(0^-\)), and the \(16\times16\) singlet of two species
\(\mathsf K^{\rm s}_{\rm LL}=C_{\rm LL}\gamma^{\rm ch}_{\rm LL}\otimes\mathsf E_{\rm LL}\otimes\mathsf E_{\rm LL}\)
(the last factor on the species index); the plain-text name Kch alone,
and the letter \(\mathsf K\) without scripts (in Theorem SF1 and in the
explanation of the folder explanations), stand for any pair matrix
(\(8\times8\) for one field, \(16\times16\) for two species), the pair
of two spinor charges being \(\Psi^T_{\rm LL}\mathsf K\Psi_{\rm LL}\);
\(\mathsf K_{12}\) is the entry of \(\mathsf K\) in row 1 and column 2 &
new & \texttt{Kch} \texttt{Kch\_iota} \texttt{Kch\_1} \texttt{Kch\_2}
\texttt{Kch\_3} \texttt{Kch\textquotesingle{}\_iota}
\texttt{Kch\_singlet} \\
\(\Pi^\iota_{\rm LL}\) & the pair operators
\(\Pi^\iota_{\rm LL}=\sum_{\vec q_{\rm LL}}\Psi_{\rm LL}(\vec q_{\rm LL})^T\mathsf K^\iota_{\rm LL}\Psi_{\rm LL}(-\vec q_{\rm LL})\),
which annihilate two spinor charges of opposite momenta & new &
\texttt{Pi\_iota} \\
\(\vec q_{\rm LL}\), \(q_{\rm LL}\), \(q^j_{\rm LL}\) & the momentum of
a particle in the space of \(\mathbb R^{1,3}_{\rm LL}\), its length and
its components & new & \texttt{q} \\
\(m_{\rm LL}\), \(\mu_{\rm LL}\), \(\Lambda_{\rm LL}\) & the mass of the
spinor charges, the chemical potential, and the momentum cutoff
(\(q_{\rm LL}<\Lambda_{\rm LL}\); \(\Lambda_{\rm LL}=1\) sets the units,
except where it varies) & new & \texttt{m} \texttt{mu}
\texttt{Lambda} \\
\(\varepsilon_{\rm LL}\), \(E^\pm_{\rm LL}\), \(\tilde E^\pm_{\rm LL}\)
& the energy \(\varepsilon_{\rm LL}=(q_{\rm LL}^2+m_{\rm LL}^2)^{1/2}\),
the quasi-particle energies of the scalar channel
\(E^\pm_{\rm LL}=[(\varepsilon_{\rm LL}\pm\mu_{\rm LL})^2+\Delta_{\rm LL}^2]^{1/2}\),
a function of the gap (written \(E^\pm_{\rm LL}(\Delta^+_{\rm LL})\) and
\(E^\pm_{\rm LL}(\Delta^-_{\rm LL})\) when \(\Delta_{\rm LL}\) is
replaced by one of the two gaps of a general state), and those of the
pseudoscalar channel,
\((\tilde E^\pm_{\rm LL})^2=\varepsilon_{\rm LL}^2+\mu_{\rm LL}^2+\Delta_{\rm LL}^2\pm2[\mu_{\rm LL}^2\varepsilon_{\rm LL}^2+m_{\rm LL}^2\Delta_{\rm LL}^2]^{1/2}\)
& new & \texttt{eps} \texttt{E\^{}+} \texttt{E\^{}-} \texttt{E\^{}(+-)}
\texttt{E\_minus\_polar} \texttt{E\_plus\_polar} \texttt{Etilde\^{}+}
\texttt{Etilde\^{}-} \texttt{Etilde\^{}(+-)} \texttt{Etilde\_minus}
\texttt{Etilde\_plus} \\
\(G_{\rm LL}\), \(G^{\rm crit}_{\rm LL}\) & the coupling of the
attractive spin-zero contact interaction (dimension mass\(^{-2}\)) and
its critical value at \(\mu_{\rm LL}=0\), above which the vacuum
condenses & new & \texttt{G} \texttt{G\_crit}
\texttt{G\_over\_Gcrit\_m0} \texttt{G\_crit\_mu0\_m0} \\
\(\mathcal H^{\rm sf}_{\rm LL}\) & the Hamiltonian of the superfluid,
\(\sum_{\vec q_{\rm LL}}\Psi^\dagger_{\rm LL}(\mathsf h_{\rm LL}-\mu_{\rm LL})\Psi_{\rm LL}-\tfrac14G_{\rm LL}\sum_\iota(\Pi^\iota_{\rm LL})^\dagger\Pi^\iota_{\rm LL}\)
& new & \texttt{Hsf} \\
\(\mathsf h_{\rm LL}(\vec q_{\rm LL})\) & the \(8\times8\) Hamiltonian
matrix of a plane wave,
\(\sum_j\alpha^j_{\rm LL}q^j_{\rm LL}+\beta_{\rm LL}m_{\rm LL}\) (times
\(\sigma^{\rm LL}_0\) on the isospin index); eigenvalues
\(\pm\varepsilon_{\rm LL}\) & new & \texttt{h(q)} \\
\(\vec d_{\rm LL}\), \(d^\iota_{\rm LL}\); \(\zeta_{\rm LL}\) & the
order parameter, the isospin vector
\(\vec d_{\rm LL}=(d^1_{\rm LL},d^2_{\rm LL},d^3_{\rm LL})\in\mathbb C^3\)
with
\(d^\iota_{\rm LL}=-\tfrac12G_{\rm LL}\langle\Pi^\iota_{\rm LL}\rangle\),
and its non-unitarity
\(\zeta_{\rm LL}=\lvert\vec d_{\rm LL}\times\vec d^{\,*}_{\rm LL}\rvert/\lvert\vec d_{\rm LL}\rvert^2\in[0,1]\)
(plain text \texttt{cross(dvec,\ dvec\^{}*)}: the cross product) & new &
\texttt{dvec} \texttt{d\_iota} \texttt{abs\_dvec\_at\_min}
\texttt{zeta} \\
\(\Delta_{\rm LL}\), \(\Delta^\pm_{\rm LL}\) & the gap:
\(\Delta_{\rm LL}=\lvert\vec d_{\rm LL}\rvert\) for a unitary state, and
the two gaps
\((\Delta^\pm_{\rm LL})^2=\lvert\vec d_{\rm LL}\rvert^2(1\pm\zeta_{\rm LL})\)
of a general state & new & \texttt{Delta} \texttt{Delta\_+}
\texttt{Delta\_-} \texttt{Delta\_(+-)} \\
\(\mathsf P_{\rm LL}\); \(\mathsf H_{\rm LL}(\vec q_{\rm LL})\);
\(\mathcal N_{\rm LL}(\vec q_{\rm LL})\) & the pairing block
\(\mathsf P_{\rm LL}=\sum_\iota d^\iota_{\rm LL}(\mathsf K^\iota_{\rm LL})^\dagger\),
the \(16\times16\) Bogoliubov-de Gennes matrix
\(\mathsf H_{\rm LL}=\begin{pmatrix}\mathsf h_{\rm LL}(\vec q_{\rm LL})-\mu_{\rm LL}&\mathsf P_{\rm LL}\\\mathsf P_{\rm LL}^\dagger&-(\mathsf h_{\rm LL}(-\vec q_{\rm LL})-\mu_{\rm LL})^T\end{pmatrix}\)
and the Nambu spinor
\(\mathcal N_{\rm LL}=(\Psi_{\rm LL}(\vec q_{\rm LL}),\Psi_{\rm LL}^{*}(-\vec q_{\rm LL}))\)
on which it acts & new & \texttt{Ppair} \texttt{H\_BdG} \\
\(\mathsf D_{\rm LL}\) & the isospin part of the pairing block of the
scalar channel, the \(2\times2\) matrix
\(\mathsf D_{\rm LL}=\sum_\iota d^\iota_{\rm LL}(\mathsf E_{\rm LL}\tau^{\rm LL}_\iota)^\dagger\),
so that
\(\mathsf P_{\rm LL}=(C_{\rm LL}\gamma^{\rm ch}_{\rm LL})^\dagger\otimes\mathsf D_{\rm LL}\);
it is also the isospin wavefunction of the condensed pairs (rows and
columns: isospin up, down); plain text \texttt{Dpair} (not
\texttt{D\_LL}, which is the Lorentz derivation) & new &
\texttt{Dpair} \\
\(\Omega_{\rm LL}\) & the grand potential per volume at zero
temperature, measured from the normal state:
\(\Omega_{\rm LL}=\lvert\vec d_{\rm LL}\rvert^2/G_{\rm LL}-\tfrac12\int\!d^3q_{\rm LL}(2\pi)^{-3}\sum(\text{positive eigenvalues of }\mathsf H_{\rm LL})\)
minus its value at \(\vec d_{\rm LL}=0\) & new & \texttt{Omega}
\texttt{Omega\_min} \texttt{Omega\_minus\_Omega0}
\texttt{Omega\_at\_polar\_gap} \\
\(k^{\rm F}_{\rm LL}\), \(\nu^{\rm F}_{\rm LL}\) & the Fermi momentum
\((\mu_{\rm LL}^2-m_{\rm LL}^2)^{1/2}\) and the density of states of one
species and one spin state at the Fermi surface,
\(\nu^{\rm F}_{\rm LL}=k^{\rm F}_{\rm LL}\mu_{\rm LL}/(2\pi^2)\) & new &
\texttt{kF} \texttt{nuF} \\
\(F_{\rm LL}\); \(U(1)^{\rm F}_{\rm LL}\); \(SU(2)_{\rm v}\) & the
generator of the fermion number
\(\Psi_{\rm LL}\to e^{i\theta_{\rm LL}}\Psi_{\rm LL}\) (it multiplies
\(\vec d_{\rm LL}\) by \(2i\)) and its group; the isospin group, which
rotates the isospin index of \(\Psi_{\rm LL}\) and \(\vec d_{\rm LL}\)
as a vector & new & \texttt{F\_LL} \texttt{U(1)\_F\_LL}
\texttt{SU(2)\_v} \\
\(\Psi^{(1)}_{\rm LL}\), \(\Psi^{(2)}_{\rm LL}\) & two species of the
field \(\Psi_{\rm LL}\), combined into one field of 16 components with a
species index & new & \\
\end{longtable}

\textbf{7 Power counting and renormalizability}

\begin{longtable}[]{@{}
  >{\raggedright\arraybackslash}p{(\columnwidth - 6\tabcolsep) * \real{0.2200}}
  >{\raggedright\arraybackslash}p{(\columnwidth - 6\tabcolsep) * \real{0.5000}}
  >{\raggedright\arraybackslash}p{(\columnwidth - 6\tabcolsep) * \real{0.0800}}
  >{\raggedright\arraybackslash}p{(\columnwidth - 6\tabcolsep) * \real{0.2000}}@{}}
\toprule\noalign{}
\begin{minipage}[b]{\linewidth}\raggedright
symbol
\end{minipage} & \begin{minipage}[b]{\linewidth}\raggedright
meaning
\end{minipage} & \begin{minipage}[b]{\linewidth}\raggedright
origin
\end{minipage} & \begin{minipage}[b]{\linewidth}\raggedright
plain-text names
\end{minipage} \\
\midrule\noalign{}
\endhead
\bottomrule\noalign{}
\endlastfoot
\([\ \cdot\ ]\) (one argument) & the mass dimension of a field or a
coupling, in units in which \(\hbar\) and the speed of light are 1 (a
bracket with two or three arguments is a commutator or an associator) &
new & \texttt{{[}Psi\_LL{]}} \texttt{{[}phi{]}} \texttt{{[}G{]}}
\texttt{{[}g\_Y{]}} \texttt{{[}lam4{]}} \\
\(\phi_{\rm LL}\), \(g^{\rm Y}_{\rm LL}\), \(\lambda^{(1)}_{\rm LL}\),
\(\lambda^{(2)}_{\rm LL}\) & the Yukawa completion: a complex
isospin-triplet pair field
\(\phi_{\rm LL}=(\phi^1_{\rm LL},\phi^2_{\rm LL},\phi^3_{\rm LL})\)
coupled to \(\Psi^T_{\rm LL}\mathsf K^\iota_{\rm LL}\Psi_{\rm LL}\) with
the Yukawa coupling \(g^{\rm Y}_{\rm LL}\), and the couplings
\(\lambda^{(1)}_{\rm LL}\), \(\lambda^{(2)}_{\rm LL}\) of the quartic
terms \((\phi^\dagger_{\rm LL}\phi_{\rm LL})^2\) and
\(\lvert\phi^T_{\rm LL}\phi_{\rm LL}\rvert^2\) (plain text
\texttt{lam4}: either of them) & new & \texttt{phi} \texttt{g\_Y}
\texttt{lam4} \texttt{lam4\_1} \texttt{lam4\_2} \\
\(\mathcal L^{\rm Y}_{\rm LL}\), \(\mathcal L^{\rm c}_{\rm LL}\) & the
Lagrangian densities of the Yukawa form and of the contact form of the
field theory of the spinor charges on \(\mathbb R^{1,3}_{\rm LL}\) & new
& \texttt{Lagr\_Y} \texttt{Lagr\_c} \\
\(m^\phi_{\rm LL}\) & the mass of the pair field \(\phi_{\rm LL}\) in
the Yukawa form & new & \texttt{m\_phi} \\
\(I^{\rm ren}_{\rm LL}\) & the renormalized one-loop bubble,
\(I^{\rm ren}_{\rm LL}(p^{\rm E}_{\rm LL})=I^{\rm bub}_{\rm LL}(p^{\rm E}_{\rm LL})-I^{\rm bub}_{\rm LL}(0)-(p^{\rm E}_{\rm LL})^2\,\partial I^{\rm bub}_{\rm LL}/\partial(p^{\rm E}_{\rm LL})^2\)
(the derivative taken at \(p^{\rm E}_{\rm LL}=0\)): the bubble minus its
mass and wave-function counterterms & new & \texttt{Iren} \\
\(\deg_{\rm LL}\) & the superficial degree of divergence of a
one-particle-irreducible graph: the power of the cutoff with which it
grows when all its loop momenta grow together & new & \texttt{deg} \\
\(n^{\rm loop}_{\rm LL}\), \(n^{\rm int,F}_{\rm LL}\),
\(n^{\rm int,B}_{\rm LL}\), \(n^{\rm ext,F}_{\rm LL}\),
\(n^{\rm ext,B}_{\rm LL}\), \(n^{\rm v}_{\rm LL}\),
\(n^{\rm v,Y}_{\rm LL}\), \(n^{\rm v,4}_{\rm LL}\) & the numbers of
loops, of internal fermion and boson lines, of external fermion and
boson lines, of vertices, of Yukawa vertices and of quartic vertices of
a graph & new & \texttt{nloop} \texttt{nint\_F} \texttt{nint\_B}
\texttt{next\_F} \texttt{next\_B} \texttt{nvert} \texttt{nvert\_Y}
\texttt{nvert\_4} \\
\(I^{\rm bub}_{\rm LL}\); \(k^{\rm E}_{\rm LL}\), \(p^{\rm E}_{\rm LL}\)
& the one-loop bubble of the Yukawa completion,
\(I^{\rm bub}_{\rm LL}(p^{\rm E}_{\rm LL})=\int_{k^{\rm E}_{\rm LL}<\Lambda_{\rm LL}}\frac{d^4k^{\rm E}_{\rm LL}}{(2\pi)^4}\frac{k^{\rm E}_{\rm LL}\cdot(k^{\rm E}_{\rm LL}+p^{\rm E}_{\rm LL})+m_{\rm LL}^2}{[(k^{\rm E}_{\rm LL})^2+m_{\rm LL}^2][(k^{\rm E}_{\rm LL}+p^{\rm E}_{\rm LL})^2+m_{\rm LL}^2]}\),
with the Euclidean loop and external four-momenta
\(k^{\rm E}_{\rm LL}\), \(p^{\rm E}_{\rm LL}\) (and their lengths) & new
& \texttt{Ibub} \texttt{Ibub\_0} \texttt{Ibub\_pE}
\texttt{Ibub\_pE\_minus\_Ibub\_0} \texttt{kE} \texttt{pE} \\
\end{longtable}

\textbf{8 Groups, Lie algebras, operations and labels}

\begin{longtable}[]{@{}
  >{\raggedright\arraybackslash}p{(\columnwidth - 6\tabcolsep) * \real{0.2200}}
  >{\raggedright\arraybackslash}p{(\columnwidth - 6\tabcolsep) * \real{0.5000}}
  >{\raggedright\arraybackslash}p{(\columnwidth - 6\tabcolsep) * \real{0.0800}}
  >{\raggedright\arraybackslash}p{(\columnwidth - 6\tabcolsep) * \real{0.2000}}@{}}
\toprule\noalign{}
\begin{minipage}[b]{\linewidth}\raggedright
symbol
\end{minipage} & \begin{minipage}[b]{\linewidth}\raggedright
meaning
\end{minipage} & \begin{minipage}[b]{\linewidth}\raggedright
origin
\end{minipage} & \begin{minipage}[b]{\linewidth}\raggedright
plain-text names
\end{minipage} \\
\midrule\noalign{}
\endhead
\bottomrule\noalign{}
\endlastfoot
\(SL(2,\mathbb C)\), \(Spin(3,1)^+\), \(SO(4,\mathbb C)\),
\(\mathrm G_2\), \(\mathrm G_2(\mathbb C)\), \(SO(1,7)\), \(Spin(7)\),
\(Spin(7,\mathbb C)\), \(SL(2,\mathbb R)\), \(Spin(2,1)\), \(SU(2)\),
\(U(1)\), \(\mathbb Z_2\) & groups: the complex \(2\times2\) matrices of
determinant 1 (the unit bi-quaternions), equal to the connected double
cover of the Lorentz group of \(3+1\) dimensions; the complex orthogonal
group of four dimensions; the compact exceptional group \(\mathrm G_2\)
(automorphisms of \(\mathbb O\)) and its complex form (automorphisms of
\(\mathbb B_{\rm LL}\)); the Lorentz group of \(7+1\) dimensions; the
spin groups of seven dimensions (compact and complex); the real
\(2\times2\) matrices of determinant 1, the double cover of the Lorentz
group of \(2+1\) dimensions; the special unitary group of two
dimensions; the phases; the group of two elements & new &
\texttt{SL(2,C)} \texttt{Spin(3,1)\^{}+} \texttt{SO(4,C)} \texttt{G2}
\texttt{G2(C)} \texttt{SO(1,7)} \texttt{SO(3,1)} \texttt{Spin(7)}
\texttt{Spin(7,C)} \texttt{SL(2,R)} \texttt{Spin(2,1)} \texttt{SU(2)}
\texttt{U(1)} \texttt{Z2} \\
\(\mathfrak{so}(3,1)\), \(\mathfrak{so}(1,7)\), \(\mathfrak{so}(4,1)\),
\(\mathfrak{so}(8)\), \(\mathfrak{so}(4,4)\),
\(\mathfrak{so}(8,\mathbb C)\), \(\mathfrak{so}(3)\),
\(\mathfrak{so}(4)\), \(\mathfrak{so}(2,1)\),
\(\mathfrak{sl}(2,\mathbb R)\), \(\mathfrak{sl}(2)\),
\(\mathfrak{su}(2)\), \(\mathfrak{spin}(7)\), \(\mathfrak g_2\),
\(\mathfrak g_2(\mathbb C)\) & Lie algebras: of the Lorentz groups of
\(3+1\) and \(7+1\) dimensions; of the de Sitter group (signature of the
Killing form: 4 positive, 6 negative); of the linear maps that preserve
a definite real, a \((4,4)\) and a complex quadratic form in eight
dimensions; of the rotations of three-dimensional space and of
four-dimensional Euclidean space; of the Lorentz group of \(2+1\)
dimensions, equal to \(\mathfrak{sl}(2,\mathbb R)\);
\(\mathfrak{sl}(2)\) the simple Lie algebra of dimension 3 over any
field; \(\mathfrak{su}(2)\) isospin; \(\mathfrak{spin}(7)\) of dimension
21; the exceptional Lie algebra \(\mathfrak g_2\) (compact, dimension
14) and its complex form (real dimension 28) & new & \texttt{so(3,1)}
\texttt{so(1,7)} \texttt{so(4,1)} \texttt{so(8)} \texttt{so(4,4)}
\texttt{so(8,C)} \texttt{so(3)} \texttt{so(4)} \texttt{so(2,1)}
\texttt{sl(2,R)} \texttt{sl(2)} \texttt{su(2)} \texttt{spin(7)}
\texttt{g2} \texttt{g2(C)} \\
\(\mathrm{tr}\), \(\det\), \(\mathrm{diag}\), \(\mathrm{Re}\),
\(\mathrm{Im}\) & trace, determinant, the diagonal matrix with the
listed entries, real part, imaginary part (also: the imaginary part of
an algebra, the span of its imaginary units); plain text \texttt{exp},
\texttt{ln}, \texttt{sqrt}, \texttt{sum}, \texttt{int}, \texttt{span},
\texttt{dim}: the exponential, natural logarithm, square root, a sum, an
integral, the set of real linear combinations, a dimension & new &
\texttt{tr} \texttt{det} \texttt{diag} \texttt{Re} \texttt{Im}
\texttt{exp} \texttt{ln} \texttt{sqrt} \texttt{sum} \texttt{int}
\texttt{span} \texttt{dim} \\
\({}^T\); \(\otimes\), \(\oplus\); \(\circ\); \(\langle\ \rangle\); h.c.
& transpose; tensor product and direct sum (plain text \texttt{(x)} for
\(\otimes\)); composition of maps; expectation value in the ground
state; the Hermitian conjugate of the preceding terms & new &
\texttt{(x)} \texttt{\^{}T} \texttt{o} \\
A-1, \ldots, R-6; G-1, \ldots, G-12; LL-S1, \ldots, LL-H1; X-1, \ldots,
X-10; Z-1, \ldots, Z-8; Theorems R, L, M, SF1, \ldots, SF4 & the
identifiers of the machine checks (letters: A associativity, L Lorentz
super-loop, Q operators and eigenvalues, C the inversion element, M
Malcev classification, G Groebner bases, N manifolds, S superfluid, R
renormalizability, LL-S symbols, LL-H self-containedness, X the examples
of the explanation of the spin-zero superfluid in the folder
explanations, Z the 3+1 manifold and the renormalizability of the Yukawa
form in the folder explanations) and the labels of the theorems proved
in the manuscript & new & \\
H-grp, H-op, H-Lie, H-gr, H-dyn & the labels of the hypotheses of the
no-go theorems: the symmetries form a group (H-grp), act as linear
operators on the space of states (H-op), have generators that form a Lie
algebra (H-Lie) or a Lie superalgebra (H-gr), and the dynamical
hypotheses on the scattering matrix (H-dyn) & new & \\
\end{longtable}

\end{document}
